% !TeX program = pdflatex
% ─────────────────────────────────────────────────────────────────────────────
%  fisica.tex — A FÍSICA ARANIANA. Compila SOZINHO (pdflatex fisica.tex) e
%  também dentro do livro. O preâmbulo é o estilo.tex.
%
%  O QUE ESTE FICHEIRO É: a matéria da aranha estigmérgica ORGANIZADA nos seis
%  lugares onde ela opera — álgebra, topologia, análise, geometria, mecânica,
%  óptica. Nada foi inventado e nada foi acrescentado: os enunciados, as provas
%  e os medidos são os que a obra já tem, postos cada um no seu sítio.
%
%  AUTOCONTIDO: toda a matemática de que precisa está DENTRO. Onde um irmão da
%  obra dizia a mesma coisa, o enunciado foi DUPLICADO e não referenciado — este
%  documento lê-se sem abrir outro. As únicas coisas de fora são o estilo da
%  capa e os MEDIDORES citados, que são programas e não texto.
%
%  E NÃO ENTRA CONJUNTO NUMÉRICO NENHUM. O documento CONSTRÓI os degraus: B são
%  os dois símbolos que o Teor. das tabelas força, I é a lista, X_2 é o segundo
%  degrau. Escrever Z, Q ou R seria usar o que se vai construir.
%
%  OS RÓTULOS levam todos o prefixo `fis:`: os corpos entram no mesmo volume e
%  há `sec:geo` em mais do que um. Sem prefixo, o hyperref liga a secção errada.
% ─────────────────────────────────────────────────────────────────────────────
\documentclass[livro.tex]{subfiles}

\begin{document}

% ─── Licença ──────────────────────────────────────────────────────────────
% Proprietário / comercial. Copyright (c) 2026 Aarão Melo Lopes.
% Uso de código ou textos mediante contrato e pagamento — ver LICENSE na raiz.
% Contacto: aarao.melo.lopes@gmail.com

% ── o que o estilo.tex não dá, e este documento usa ──
% O estilo dá teorema/proposicao/corolario/definicao/obs, \code, \medido e
% \gkcapa. Não dá `lema`, nem `\colunas`, nem a cor `cinza`. Tudo guardado: um
% \newtheorem repetido é erro FATAL quando os corpos entram no mesmo documento.
% O teste é `\lema` e NÃO `\c@lema`: o lema partilha o contador `teorema`, pelo
% que `\c@lema` nunca chega a existir e o guarda seria sempre falso.
\ifcsname lema\endcsname\else
  \theoremstyle{plain}\newtheorem{lema}[teorema]{Lema}
\fi
\providecolor{cinza}{gray}{0.35}
\providecommand{\abs}[1]{\lvert #1\rvert}
%% O SEGUNDO DEGRAU, e não um conjunto importado: este documento CONSTRÓI-O
%% (§fis:escada, X_2 = X_1^2/~). Escrever Z seria usar o que se vai construir.
\providecommand{\Zc}{\ensuremath{X_2}}
\providecommand{\sepcol}{\quad\penalty0$\big|$\quad\penalty0}
\providecommand{\colunas}[3]{%
\par\smallskip\noindent{\footnotesize\raggedright
\textbf{prova}: #1\sepcol\textbf{realiza}: #2\sepcol\textbf{verifica}: #3\par}\smallskip}

\ifSubfilesClassLoaded{%
  \gkcapa{A Física Araniana}
         {multiplicidade e memória}
         {álgebra, topologia, análise, geometria, mecânica, óptica, partículas e núcleo,
          termodinâmica, cosmologia, redes --- dez lugares, um sistema}
  \maketitle
  \tableofcontents
}{%
  \part*{A Física}\addcontentsline{toc}{part}{A Física}
}

\begin{abstract}\noindent
\textbf{De onde parte.} De \textbf{dois extremos e de uma dobra}: a acção que os
troca, $\partial x=a+b-x$. O seu ponto fixo é o \emph{meio}, e dobrando outra vez
de cada lado --- a mesma estrutura nos dois --- a indução dá todos os pontos. Nos
extremos sozinhos não há meio, e é essa falta que obriga a \emph{duas}
composições, com as duas equações $\partial x+x=0$ e $\partial x\cdot x=1$: é
nelas que $0$ e $1$ nascem, como o que a dobra devolve em cada face. As tabelas
saem daí forçadas, sem que se escreva um sinal, e delas a soma com transporte; a
\textbf{lista} $I$ é o objeto finito que daí se obtém. Só então entra uma
\textbf{realização} $\pi\colon I\to X$ num conjunto $X$ com igualdade decidível,
e não se supõe métrica, nem dimensão, nem estrutura algébrica em $X$. A
realização identifica índices distintos na mesma célula; essa identificação
\emph{é} a dobra, e a \textbf{multiplicidade} $G(x)=\abs{\pi^{-1}(x)}$ conta-a.

\medskip\noindent\textbf{Para onde vai.} Para o par de operações que a aranha
executa, e que são uma inversa da outra: escrever o campo é a
\textbf{convolução} com $\zeta$ na ordem $(I,\le)$; recuperar o passo é a
\textbf{deconvolução} com $\mu=\zeta^{-1}$, que nessa ordem é a diferença finita.
Enquanto a família $\{G_t\}$ existe, a trajetória inteira volta sem nada dela ter
sido guardado; retendo só $G$, o que se perde repõe-se com \emph{uma} coordenada
de folha, seja qual for a dimensão, e com volta exata. Pelo caminho: a dimensão
de $X$ entra num sítio único --- a vizinhança, $\abs{\mathcal V(x)}=2n$ na grade
e $1$ na árvore; sobre ultramétricos a célula \emph{é} a bola; numa órbita
determinista $G$ conta o período; e a dobra lê-se na norma do espectro, sem
visitar a trajetória.

\medskip\noindent\textbf{E o eixo é uma dualidade, que não foi escolhida:
deduziu-se.} Sobre dois extremos $\partial$ não fixa nada, e por isso uma
composição não chega --- são precisas \textbf{duas}, com $\partial$ a ser o
\emph{oposto} numa e o \emph{inverso} na outra. Cada face tem a sua dobra e cada
dobra tem o seu meio: a \textbf{aritmética} $\frac{a+b}{2}$ de um lado, a
\textbf{geométrica} $\sqrt{ab}$ do outro; a reta e o círculo; o cartesiano e o
polar, onde as fases somam e os módulos multiplicam. E as duas \emph{fecham}:
batendo alternadas, os intervalos encaixam e o processo \textbf{termina numa
igualdade} --- o ponto onde para é o ponto fixo que as duas faces passam a
partilhar, e é ele o \textbf{corte}. Não converge para o ponto: termina nele.

\medskip\noindent\textbf{A organização deste documento.} A matéria é uma só e
opera em \textbf{dez} lugares, e é por eles que ela se reparte aqui. A
\textbf{álgebra} constrói: a dobra, as duas faces, os dois símbolos, a lista, a
multiplicidade, o par $\zeta/\mu$, a transformada, a escada e a caixa. A
\textbf{topologia} conta o que a construção identifica: a ultramétrica que
particiona, a característica que lê a dobra em ciclos, o grau que proíbe, e a
folga vista pelos dois mapas. A \textbf{análise} faz o encaixe terminar, e dá o
corte e o número $\pi$. A \textbf{geometria} realiza: o semicírculo, o
hipercubo, a base ortonormal e o dragão. A \textbf{mecânica} põe a construção a
correr: o ciclo, a lei local, o levantamento, o tempo e a régua que a memória
curva. A \textbf{óptica} lê os espelhos: a caixa dos sistemas sem perda, a lei
do discriminante, as ordens que uma rede admite e a borda onde a régua colapsa.

\medskip\noindent E os quatro últimos lugares são o mesmo sistema em quatro
escalas. As \textbf{partículas e o núcleo}: o átomo como cadeia espectral, a
massa como distância ao círculo, a ligação como folga, e o mapa espectral que diz
que átomos a álgebra permite. A \textbf{termodinâmica} e a \textbf{cosmologia},
que são um par --- \emph{entropia é o que se desfaz, expansão é o que se abre} ---
com o rendimento a sair do $\oint$ fechado, o custo a ser a dobra, a segunda lei a
valer de \emph{um} lado, e a taxa de expansão a ser a régua invertida. E as
\textbf{redes}: o neurónio como cisão e soma, a sobreposição como prefixo comum,
a interferência como dobra, e as duas torres --- a que mede e a que ordena.

\medskip\noindent\textbf{O alcance.} De um operador com duas equações e de uma
aplicação com igualdade no contradomínio saem, com demonstração, o percurso, a
memória, a multiplicidade, o tempo, a métrica, a transformada, o corte, a
classificação dos sistemas sem perda, o espectro de um átomo, a expansão de um
universo, a memória de uma rede --- e as dualidades que os ligam.

\medskip\noindent Todo enunciado usado é enunciado e demonstrado aqui.
\end{abstract}

% ═══════════════════════════════════════════════════════════════════════════
\section*{As perguntas, e onde estão respondidas}\label{fis:roteiro}

\begin{center}\small
\begin{tabular}{@{}p{0.44\textwidth}p{0.32\textwidth}l@{}}
\toprule
a pergunta & a resposta & onde \\
\midrule
O que é o ponto de partida? & dois extremos e a dobra que os troca & Def.~\ref{fis:def:op} \\
Donde vêm as operações e os dois dígitos? & da dobra: sem meio, são precisas duas faces & Teor.~\ref{fis:thm:duas},~\ref{fis:thm:B} \\
E as coordenadas? & a tupla ordenada: a \emph{representação} & Teor.~\ref{fis:thm:coord} \\
Como constrói? & a aresta: ela \emph{é} a memória da divisão & §\ref{fis:hipercubo} \\
Como sabe por onde passou? & não sabe --- a marca está no espaço & Teor.~\ref{fis:thm:mult}(3) \\
O que acontece quando volta? & a dobra: $G>1$ & Def.~\ref{fis:def:objeto} \\
Como se representam os passos? & por uma palavra no alfabeto das arestas & Teor.~\ref{fis:thm:palavra} \\
E se forem muitos construtores? & o campo não os distingue & Teor.~\ref{fis:thm:agentes} \\
Como se desfaz a sobreposição? & por uma coordenada de folha & Teor.~\ref{fis:thm:levant} \\
Onde entra a dimensão? & na vizinhança, e só aí & Teor.~\ref{fis:thm:dim} \\
De onde vem o tempo? & da contagem das aplicações do operador & Def.~\ref{fis:def:tempo} \\
O que curva a régua? & o campo: $g_{uv}=G\,h_{uv}$ & §\ref{fis:tempo} \\
O que classifica um sistema sem perda? & $\Delta$, e não o traço & Teor.~\ref{fis:thm:leidisc} \\
Que ordens admite uma rede? & $\{1,2,3,4,6\}$, e nunca $5$ & Lema~\ref{fis:lem:cristal} \\
\bottomrule
\end{tabular}
\end{center}

% ═══════════════════════════════════════════════════════════════════════════
\part{Álgebra}
% ═══════════════════════════════════════════════════════════════════════════

\section{O operador, e as duas faces}\label{fis:dobra}

\noindent Em quatro frases, e cada uma é um resultado abaixo:
\[
\boxed{
\begin{array}{l}
\text{O finito é \textbf{codificado}.}\\
\text{O infinito é \textbf{caminhado}.}\\
\text{A representação \textbf{comprime}.}\\
\text{A dobra \textbf{mede} o que foi perdido.}
\end{array}}
\]

\medskip\noindent O documento parte de \emph{dois} extremos e de \emph{uma}
operação sobre eles: dobrar ao meio. Não se supõem números, nem símbolos, nem
quantos são --- tudo o que segue é o que a dobra produz, e produz por indução.

\begin{definicao}[o operador]\label{fis:def:op}
Dado um intervalo de extremos $a$ e $b$, a \textbf{dobra} é a acção que os
\emph{troca}:
\[
\boxed{\;\partial x \;=\; a+b-x\;}
\]
Ela \textbf{conserva}: para todo $x$, a composição $x\circ\partial x$ vale sempre
o mesmo --- a soma dos extremos --- e a esse valor comum chama-se o
\textbf{invariante}. Pede-se ainda que $\partial$ seja \emph{efetiva}
($\partial\neq\mathrm{id}$) e que tenha \textbf{volta}: de $\partial x$
recupera-se $x$. Sem a volta não há o que conservar, e é ela que distingue uma
dobra de um esmagamento.

\medskip\noindent Normalizando os extremos --- que é só escolher a régua --- fica
$\partial x=1-x$, e o invariante é uma \textbf{unidade}:
\[
\boxed{\;\abs{\,x\circ\partial x\,}=1\;}
\]
O $\abs{\cdot}$ diz apenas isso. Não diz qual, e não traz sinal nenhum: nada do
que segue precisa de escrever um número negativo.
\end{definicao}

\begin{teorema}[uma composição não chega]\label{fis:thm:duas}
Com $\partial$ da Def.~\ref{fis:def:op}:
\begin{enumerate}\itemsep3pt
\item[(1)] \textbf{O invariante é o neutro.} O valor comum $e$ cumpre
$e\circ x=x\circ e=x$.
\item[(2)] \textbf{$\partial$ é involutiva, e tem um ponto fixo: o MEIO.} São
factos da dobra e não dependem de composição nenhuma:
\[
\partial^{2}=\mathrm{id},\qquad \partial m=m,\qquad m=\tfrac{a+b}{2} .
\]
O meio é \emph{um só} e não é acrescentado ao conjunto --- é \textbf{produzido}
pela dobra.
\item[(3)] \textbf{E onde houver composição, é ele o neutro dela.} Por (1) o
invariante é o neutro, e por (2) o invariante de uma composição compatível com
$\partial$ tem de ser fixo por $\partial$; havendo um só fixo, é o meio.
\emph{Enquanto não houver a face, o meio é só o centro} --- a identificação com
o neutro nasce com a composição, não antes dela.
\item[(4)] \textbf{Nos extremos sozinhos não há meio, e por isso uma composição
não chega.} Sobre dois objectos apenas, $\partial$ troca-os e não fixa nenhum;
então (3) não pode valer numa composição só --- o invariante teria de ser fixo,
e não há fixo nenhum. É preciso uma \textbf{segunda}, e o seu invariante é o
\emph{outro} extremo:
\[
\boxed{\;\partial x + x = 0\;}\qquad\qquad\boxed{\;\partial x\cdot x = 1\;}
\]
--- $\partial$ é o \emph{oposto} numa face e o \emph{inverso} na outra, e o
invariante de cada uma é o neutro da outra.
\item[(5)] \textbf{Dobrando, a estrutura repete-se dos DOIS lados.} O meio parte
o intervalo em $[a,m]$ e $[m,b]$, e cada um, reescalado, é uma cópia exata do
inteiro --- mesma dobra, mesmo meio, mesmos extremos. \emph{O ponto fixo de uma
dobra é extremo da seguinte.}
\item[(6)] \textbf{E a indução dá todos os pontos.} Ao fim de $n$ dobras há
$2^{n}$ intervalos e $2^{n}+1$ pontos, e eles são exatamente os $k/2^{n}$ com
$0\le k\le 2^{n}$. Nada se escolhe pelo caminho: cada nível é o anterior dobrado,
e a escolha de \emph{qual lado} em cada passo é a coordenada que o
Teor.~\ref{fis:thm:coord} vai ler.
\end{enumerate}
\end{teorema}

\begin{proof}
(1) Aplicando a definição a $x$ e a $\partial x$ vem
$x\circ\partial x=e=\partial x\circ\partial\partial x$. Compondo
$e\circ x=(x\circ\partial x)\circ x$ e usando a associatividade, o valor não
depende de $x$; avaliando em $x=e$ obtém-se $e\circ e=e$, e daí $e\circ x=x$.

(2) $\partial\partial x=a+b-(a+b-x)=x$, directamente; e $a+b-x=x$ tem a única
solução $x=\tfrac{a+b}{2}$, que existe sempre que o $2$ se inverta. Note-se que
\emph{escrever} a dobra já usa o dois --- é o número de extremos --- pelo que
tê-lo é condição de a construção operar; e tê-lo é ter o seu dual $\tfrac12$,
que é o próprio meio. Nada disto menciona composição: são factos da reflexão.

(3) Dada uma composição $\circ$ que cumpra a Def.~\ref{fis:def:op}, o seu
invariante $e$ é neutro por (1), e de $e\circ\partial e=e$ com
$e\circ\partial e=\partial e$ vem $\partial e=e$: o invariante é fixo por
$\partial$. Como o único fixo é $m$, é $m$. A implicação corre neste sentido e
não no outro --- é a face que obriga o seu neutro a ser o centro, e não o centro
que se declara neutro por si.

(4) Se $\partial$ não fixa nada, (3) não pode valer numa composição só. Com dois
objetos trocados, o invariante teria de ser fixo e nenhum o é; restam
\emph{duas} composições, cada uma com o seu invariante, e os dois invariantes são
os dois objetos. Chamando-lhes $0$ e $1$, e $+$ e $\cdot$ às composições, a
Def.~\ref{fis:def:op} lê-se nas duas equações do enunciado.

(5) A aplicação afim que leva $[a,m]$ em $[a,b]$ transporta a dobra na dobra,
porque $\partial$ é afim e os afins compõem-se; o mesmo do outro lado. E $m$ é
fixo em $[a,b]$ por (2) e é extremo de cada pedaço por construção --- é a mesma
frase dita nos dois níveis.

(6) Por indução: o nível $0$ tem $1$ intervalo e $2$ pontos; dobrar cada
intervalo duplica-os e acrescenta um ponto por intervalo, o que dá $2^{n}$ e
$2^{n}+1$. Que os pontos sejam os $k/2^{n}$ sai de (3) aplicado a cada pedaço,
cujo meio é a média de dois diádicos consecutivos.
\end{proof}

\medido{\code{tests/meio.c} §M5--§M6 e \code{tests/aranha\_n.c} §AN36. Sobre três
objetos, as composições com $\partial$ efetivo cujo invariante é neutro são
$\mathbf{81}$, e em \emph{todas} o neutro é fixo por $\partial$ --- que é (2)
medido; ao passo que sobre dois, com $\partial$ a trocá-los, não há
\textbf{nenhuma}, que é (4). Com duas composições de invariantes distintos há
exactamente $\mathbf{2}$ soluções, com as entradas forçadas. E o nível $1$, que
já contém o seu centro, precisa de \emph{uma} só: as composições que conservam
com o invariante por neutro são nove, em todas ele é o meio e nunca um extremo,
e exigindo comutatividade e associatividade sobra \textbf{uma} --- o ciclo de
ordem três, com o meio de ordem $1$ e os extremos de ordem $3$. Não são dois
teoremas: é o mesmo, medido dos dois lados. Os primeiros níveis de (6) em §M3.}

\section{Os dois símbolos, a lista, e o par centro/desvio}\label{fis:binario}

\begin{teorema}[as soluções são únicas]\label{fis:thm:B}
No nível $0$ --- os dois extremos, sem meio --- as duas composições do
Teor.~\ref{fis:thm:duas}(4) ficam \textbf{forçadas}, e sem que se escreva um
sinal:
\[
\begin{array}{c|cc}
\cdot & 0 & 1\\\hline
0 & 0 & 0\\
1 & 0 & 1
\end{array}
\qquad\qquad
\begin{array}{c|cc}
+ & 0 & 1\\\hline
0 & 0 & 1\\
1 & 1 & 0
\end{array}
\]
Em particular $1+1=0$: \textbf{o oposto de $1$ é $1$}, e o $-1$ nunca precisa de
ser escrito --- o que noutras aritméticas seria um símbolo novo é aqui o próprio
$1$. Cada objeto é o seu próprio oposto, que é a primeira equação lida sobre os
dois.
\end{teorema}

\begin{proof}
O invariante da segunda face é o neutro dela, logo $1\cdot x=x$; e
$\partial 1=0$, pelo que $0$ fica absorvente, $0\cdot x=0$ --- as quatro entradas
ficam determinadas. Para a primeira, $0$ é neutro e falta $1+1$: da
Def.~\ref{fis:def:op} em $x=1$ vem $\partial 1+1=0$, isto é $0+1=1$, e em $x=0$
vem $1+0=1$; compondo $\partial$ consigo lê-se $\partial\partial=1$, donde
$1+1=0$.

Sobre dois símbolos com invariantes distintos sobrevivem exatamente \emph{duas}
soluções, e a segunda é a primeira com $0$ e $1$ trocados: a mesma estrutura
noutra nomeação. Partir do $0$ e construir o $1$, ou o contrário, é o mesmo
percurso lido nos dois sentidos.
\end{proof}

\begin{definicao}[o binário]\label{fis:def:B}
$B$ é o conjunto dos dois símbolos do Teor.~\ref{fis:thm:B}.
\end{definicao}

\medido{\code{tests/meio.c} §M7: varridas as tabelas comutativas com esses
neutros há exatamente \textbf{uma} de cada. E \code{tests/aranha\_n.c} §AN33: as
quatro aplicações, as duas invertíveis e involutivas, e a única tabela de
produto em dezasseis.}

\begin{teorema}[a lista, do binário]\label{fis:thm:BI}
Sobre palavras em $B$, com $\oplus$ e $\wedge$ as duas operações do
Teor.~\ref{fis:thm:B} aplicadas posição a posição, vale
\[
\boxed{\;a+b=(a\oplus b)+2\,(a\wedge b)\;}
\]
e a iteração $(a,b)\mapsto\bigl(a\oplus b,\,2(a\wedge b)\bigr)$ termina, porque a
segunda componente sobe de posição a cada passo e a largura é finita. A soma
assim obtida enumera as palavras, e o \textbf{encaixe por prefixos} ordena os
andares uns dentro dos outros,
\[
\{0,1\}\subset\{0,\dots,3\}\subset\{0,\dots,15\}\subset\{0,\dots,255\},
\]
com as palavras de largura $w$ em número de $2^{w}$ e a dobra a duplicar a
largura. A \textbf{lista} $I$ é essa ordem.

\medskip\noindent\textbf{E a mesma identidade lê-se a partir do meio.} O meio da
Def.~\ref{fis:def:op} é $m=\tfrac{a+b}{2}$; escrevendo o \textbf{desvio}
$\delta=\tfrac{a-b}{2}$ \emph{com sinal},
\[
\boxed{\;a \;=\; m+\delta\,,\qquad\qquad b \;=\; m-\delta\;}
\]
--- igualdade e não aproximação: $(m,\delta)$ e $(a,b)$ são a mesma coisa escrita
duas vezes, e a volta devolve o par \emph{ordenado} com resíduo $0$. É a dobra
outra vez: $m$ é o que $\partial$ fixa, $\delta$ o que ele reflete.

\medskip\noindent\textbf{Só depois entram $\vee$ e $\wedge$}, e é aí que se
perde: $a\vee b=m+\abs{\delta}$ e $a\wedge b=m-\abs{\delta}$. O módulo não é o
objeto --- é a \emph{intensidade}, e a intensidade é \textbf{metade}, porque
deixa cair o sinal de $\delta$, isto é, \emph{qual dos dois é qual}.

\medskip\noindent\textbf{E daqui sai a conservação, que não foi pedida.}
Expandindo, $a\,b=(m+\delta)(m-\delta)=m^{2}-\delta^{2}$; e $\partial$, que troca
$a$ com $b$, é $\delta\mapsto-\delta$ com $m$ fixo. Como $\delta$ entra ao
\textbf{quadrado}, o que $\partial$ move não sobrevive, e o produto é
\emph{invariante}. A conservação da norma não é hipótese posta à entrada --- é o
que a decomposição dá depois de as leis estarem no sítio.
\end{teorema}

\begin{proof}
Posição a posição, para $a_j,b_j\in B$ verifica-se nos quatro casos que
$a_j+b_j=(a_j\oplus b_j)+2(a_j\wedge b_j)$; multiplicando por $2^{j}$ e somando
vem a identidade. Na iteração, $a\wedge b$ multiplicado por $2$ desloca-se uma
posição para cima, de modo que ao fim de tantos passos quantos os bits da largura
a segunda componente é nula. O encaixe é a restrição das operações de um andar ao
anterior, e a ordem resultante é total.

A forma do meio é imediata e não usa a identidade: de $m=\tfrac{a+b}{2}$ e
$\delta=\tfrac{a-b}{2}$ vem $m+\delta=a$ e $m-\delta=b$, somando e subtraindo. O
sistema é triangular e inverte-se sem condição nenhuma --- é por isso que a volta
é exata, e é por isso que \emph{não} aparece aqui módulo nenhum. Os extremos
$\vee$ e $\wedge$ pedem-no: escolher o maior é perguntar o \emph{sinal} de
$\delta$ e deitá-lo fora, ficando com $\abs{\delta}$. Por fim
$ab=(m+\delta)(m-\delta)=m^{2}-\delta^{2}$ é o produto expandido pela
distributividade, e a troca $a\leftrightarrow b$ é $\delta\mapsto-\delta$ com $m$
fixo, que a potência par absorve.
\end{proof}

\medido{\code{tests/meio.c} §M3--§M4: os nove pares ordenados de
$\{0,\tfrac12,1\}$ caem em \textbf{seis} imagens por $(a\vee b,a\wedge b)$ --- e
as três colisões são os pares trocados --- contra \textbf{nove} em nove por
$(m,\delta)$. O controlo exige que a troca mova de facto o desvio, em seis dos
nove pares. E \code{tests/aranha\_n.c} §AN34: a codificação em
$\lceil\log_2 c\rceil$ bits para $c$ de $1$ a $64$, e nenhum a menos.}

\begin{teorema}[o processo dá as coordenadas, e elas são a representação]\label{fis:thm:coord}
Aplicando $\partial$ sucessivamente --- ao argumento, depois ao que dele saiu,
depois ao que daí saiu --- fica registada, por ordem, qual metade se tomou em
cada passo. Essa sucessão é a \textbf{tupla}
\[
\boxed{\;x=(x_1,x_2,\dots,x_w)\in B^{w}\;}
\]
e sobre ela valem três coisas:
\begin{enumerate}\itemsep3pt
\item[(1)] \textbf{Ordenada, e a ordem é a do operador.} O índice $j$ não é uma
etiqueta escolhida: é a \emph{contagem das aplicações} de $\partial$. Permutar as
componentes muda o objeto representado.
\item[(2)] \textbf{É a representação, não o objeto.} A tupla é a lista das
decisões que levaram ao objeto; o objeto é aquilo a que elas levaram.
\item[(3)] \textbf{E a passagem de uma à outra perde.} Duas tuplas distintas
podem representar o mesmo objeto, e uma largura finita obriga a que representem
--- é a \textbf{dobra}, medida por $G$ (Def.~\ref{fis:def:objeto}).
\end{enumerate}
\end{teorema}

\begin{proof}
(1) Por construção: o passo $j$ só existe depois do passo $j-1$, porque divide o
que ele produziu; a sucessão é portanto ordenada pelo próprio processo, e é essa
a única ordem que há. (2) A tupla é um objeto de $B^{w}$ e o representado vive
onde o processo o pôs; a aplicação de um para o outro é o que a
Def.~\ref{fis:def:objeto} chama realização, e nada obriga a que seja injetiva.
(3) É a contagem: em largura $w$ há $2^{w}$ tuplas, e realizar mais elementos do
que esse número identifica-os --- que é a cláusula (2) do
Teor.~\ref{fis:thm:mult}.
\end{proof}

\medskip\noindent\textbf{Donde as três palavras deste documento, já todas
presentes.} O \emph{operador} decide; a \emph{tupla} regista as decisões, por
ordem; e o \emph{objeto} é onde elas caem. A tupla ordenada de $w$ decisões é o
que adiante se chama a \textbf{palavra} (Teor.~\ref{fis:thm:palavra}) --- lá o
alfabeto é o das direções de aresta, e a divisão que $\partial$ faz no seu
argumento é a que o vértice faz nas suas saídas; a \textbf{aresta é a memória
dessa divisão}. E $G$ conta o que a passagem da tupla ao objeto deitou fora.

\medido{\code{tests/aranha\_n.c} §AN35: as duas nomeações a darem tabelas
isomorfas pela troca, com a identidade a NÃO servir como controlo ($6$
discordâncias); e permutar a tupla a mudar o objeto em $\mathbf{32}$ das $64$,
com as outras $32$ a serem exatamente as de componentes iguais.}

\section{A realização, e o que ela perde}\label{fis:mult}

\begin{definicao}[trajetória, realização, multiplicidade]\label{fis:def:objeto}
Seja $I$ a lista do Teor.~\ref{fis:thm:BI} --- ou qualquer restrição finita dela
--- e $X$ um conjunto com igualdade decidível. Uma \textbf{realização} é uma
aplicação
\[
\pi\colon I\longrightarrow X ,
\]
lida como «o índice $i$ ocupa a célula $\pi(i)$». A \textbf{multiplicidade
geométrica} de $\pi$ é o cardinal da fibra,
\[
G(x)=\bigl\lvert\pi^{-1}(x)\bigr\rvert
=\bigl\lvert\{i\in I:\pi(i)=x\}\bigr\rvert ,
\]
finito por $I$ o ser. O \textbf{suporte} de $G$ é
$\operatorname{supp}G=\{x: G(x)>0\}=\operatorname{im}\pi$.

\medskip\noindent A relação $i\sim j \iff \pi(i)=\pi(j)$ é uma equivalência em
$I$. Chama-se \textbf{dobra} de $\pi$, e as suas classes são as fibras. É a
\emph{mesma} palavra da Def.~\ref{fis:def:op}, e de propósito: lá a dobra é a
\textbf{acção} que traz um extremo sobre o outro; aqui é a
\textbf{identificação} que essa acção produz --- o que se sobrepõe passa a ser
indistinguível, e é isso que a equivalência regista. Uma é o gesto, a outra é a
marca que ele deixa.

\medskip\noindent Note-se o que \emph{não} está na definição: nem estrutura
algébrica em $X$, nem métrica, nem dimensão. Só igualdade. Esta economia não é
elegância --- é o conteúdo do Teor.~\ref{fis:thm:mult}, e é o que se paga no
Teor.~\ref{fis:thm:dim}, onde a dimensão finalmente aparece e se vê que aparece
uma vez só.
\end{definicao}

\begin{lema}[conservação]\label{fis:lem:conserva}
$\displaystyle\sum_{x\in X}G(x)=\abs{I}$, e a soma é finita porque
$\operatorname{supp}G$ tem no máximo $\abs{I}$ elementos.
\end{lema}

\begin{proof}
As fibras $\pi^{-1}(x)$, $x\in\operatorname{im}\pi$, particionam $I$, e a
cardinalidade é aditiva sobre partições finitas:
$\sum_x\lvert\pi^{-1}(x)\rvert=\lvert I\rvert$. Fora do suporte os termos são
nulos.
\end{proof}

\begin{teorema}[a medida conserva-se, a fibra perde-se]\label{fis:thm:medida}
Sejam $\pi,\pi'\colon I\to X$ duas realizações com o mesmo $\abs{I}$. Então
\[
\sum_x G_\pi(x)=\sum_x G_{\pi'}(x)=\abs{I}
\]
sempre, quaisquer que sejam as duas. Em particular a medida de contagem não
distingue realizações: existem $\pi$ e $\pi'$ com a mesma medida e dobras
diferentes, isto é, com $\max G_\pi\neq\max G_{\pi'}$.
\end{teorema}

\begin{proof}
A igualdade é o Lema~\ref{fis:lem:conserva} aplicado duas vezes. Para a segunda
afirmação basta uma testemunha: a curva de Heighway de $2^{12}$ passos e a
espiral quadrada do mesmo comprimento têm ambas $\abs{I}=4097$ e
$\sum G=\abs{I}$, e $\max G=2$ contra $\max G=1025$ (§\ref{fis:realizacoes}).
\end{proof}

\medskip\noindent\textbf{É esta a razão de ser de $G$.} A projeção identifica
índices, e o que dela se conserva --- a medida --- é cego a essa identificação:
duas patas produzem uma marca, e a marca não diz quantas patas a produziram. $G$
reconta exatamente o que a medida esquece, e por isso a informação tem de ficar
\textbf{no espaço}: não é do agente, que não a guarda, nem do procedimento, que
não a tem. É a marca, e a marca está no chão.

\medskip\noindent Daqui em diante, sempre que uma quantidade se propuser a
distinguir realizações, a pergunta é se ela é ou não invariante da medida ---
porque as que são, não distinguem.

\begin{teorema}[multiplicidade geométrica]\label{fis:thm:mult}
Com a Def.~\ref{fis:def:objeto}, sejam $G_t(x)=\lvert\{i\le t:\pi(i)=x\}\rvert$
os campos parciais. Então:
\begin{enumerate}\itemsep3pt
\item[(1)] \textbf{Duplicidade é dobra.} Para $i\neq j$ tem-se $\pi(i)=\pi(j)$
se e só se $i$ e $j$ pertencem à mesma classe de $\sim$. A dobra da
Def.~\ref{fis:def:objeto} é a fibra de $\pi$ --- não uma metáfora dela.
\item[(2)] \textbf{$G$ regista a dobra.} $G(x)>1$ se e só se existem $i\neq j$
com $\pi(i)=\pi(j)=x$. Em particular $\pi$ é injetiva se e só se $G\le1$.
\item[(3)] \textbf{A memória não pertence ao agente.} Vale a recorrência local
\[
G_{t+1}(x)=G_t(x)+\mathbf{1}_{\{x=\pi(t+1)\}},\qquad G_{-1}\equiv0,
\]
e $G=G_N$. Cada passo altera $G$ num único ponto, e esse ponto é a célula
corrente. Nenhum passo consulta $\pi(0),\dots,\pi(t)$.
\item[(4)] \textbf{Sentir é ler $G$.} A leitura local
$\mathcal{P}_t(x)=\bigl(G_t(y)\bigr)_{y\in \mathcal{V}(x)}$ sobre uma vizinhança
$\mathcal{V}(x)$ distingue $G=0$ (célula ainda não realizada) de $G>0$ (já
realizada). A distinção é propriedade do \emph{espaço}, e não propriedade
cognitiva do agente.
\end{enumerate}
\end{teorema}

\begin{proof}
(1) e (2) são a definição de fibra e a sua cardinalidade.

(3) Fixe-se $x$. Por definição $G_{t+1}(x)=\lvert\{i\le t+1:\pi(i)=x\}\rvert$, e
o conjunto contado difere do de $G_t(x)$ exatamente pelo elemento $t+1$, que lá
está se e só se $\pi(t+1)=x$. Donde a recorrência. A cláusula «num único ponto» é
a leitura da mesma igualdade: para $x\neq\pi(t+1)$ a indicatriz anula-se e
$G_{t+1}(x)=G_t(x)$. O termo direito só depende de $t+1$ através de $\pi(t+1)$,
que é o dado corrente: a história não é consultada.

(4) Por (2), $G_t(y)=0$ significa que nenhum índice $\le t$ realizou $y$, e
$G_t(y)>0$ que algum realizou. A leitura de $\mathcal{P}_t(x)$ decide isso para
cada $y\in\mathcal{V}(x)$ consultando o valor escrito \emph{em} $y$ --- valor que
está no espaço, e não no agente, por (3). Nenhum passo desta leitura consulta
$\pi(0),\dots,\pi(t)$.
\end{proof}

\medskip\noindent\textbf{O que a prova usou.} As cláusulas (1)--(3) usam só a
igualdade em $X$: para decidir $x=\pi(t+1)$ e para formar as fibras. Nenhuma
delas menciona mais estrutura. A cláusula~(4) pede uma coisa a mais, e uma só ---
a \emph{vizinhança} $\mathcal{V}(x)$ ---, e é lá que a estrutura de $X$ entra, na
§\ref{fis:grau}. É este o resultado: a memória estigmérgica não é um fenómeno do
plano, e o sítio por onde o plano entrava está isolado numa cláusula.

\medskip\noindent A memória estigmérgica é a composição
\[
\boxed{\ \text{trajetória}\ \xrightarrow{\ \pi\ }\ \text{dobra}\
\xrightarrow{\ \lvert\pi^{-1}\rvert\ }\ \text{multiplicidade}\
\xrightarrow{\ \text{leitura local}\ }\ \text{percepção}\ }
\]
e a última seta termina em \textbf{percepção}, não em decisão: o que a leitura
produz é a distinção realizada/não realizada, e essa distinção é do espaço.

\medido{\code{tests/aranha\_n.c} §AN1--§AN2: campo incremental igual à
recontagem e $\sum G=\abs{I}$ nos postos $1$ a $4$. §AN16: a mesma medida com
dobras diferentes --- Heighway e a espiral, com a reta como controlo.}

\section{O directo é $\zeta$, o inverso é $\mu$}\label{fis:zetamu}

\noindent As duas operações da aranha --- escrever o campo e recuperar o passo
--- não são duas invenções. São a \textbf{convolução} e a \textbf{deconvolução}
da álgebra de incidência da ordem $(I,\le)$.

\begin{definicao}[os dois operadores]\label{fis:def:ops}
No domínio, o \textbf{deslocamento} $S$ da tupla, que leva cada símbolo no
seguinte da lista. No contradomínio, a \textbf{adjacência} $A$ do grafo,
$(Af)(x)=\sum_{y\in\mathcal{V}(x)}f(y)$.
\end{definicao}

\begin{teorema}[o que $S$ gera]\label{fis:thm:shift}
Sobre sucessões que se anulam abaixo de $0$, a série $\zeta=\sum_{j\ge0}S^{j}$
está bem definida --- em cada índice tem um número finito de termos --- e
\[
\zeta\,a\,(t)=\sum_{u\le t}a(u),\qquad (1-S)\,\zeta=\mathrm{id}.
\]
Isto é: $S$ gera a \textbf{acumulação} $\zeta=(1-S)^{-1}$, cujo inverso
$\mu:=1-S$ é a \textbf{diferença finita}.
\end{teorema}

\begin{proof}
$(\sum_{j\ge0}S^j a)(t)=\sum_{j\ge0}a(t-j)$, e os termos com $j>t$ anulam-se,
donde a soma é $\sum_{u\le t}a(u)$. Aplicando $1-S$ a esse resultado telescopa:
$\sum_{u\le t}a(u)-\sum_{u\le t-1}a(u)=a(t)$.
\end{proof}

\begin{definicao}[$W$, e $S$ como matriz]\label{fis:def:incid}
$W=\Zc^{\,I}$ com a base canónica $\{\delta_t\}$, $\dim W=\abs{I}$ --- e $\Zc$ é
o \textbf{segundo degrau}, que a §\ref{fis:escada} constrói como $X_1^{\,2}/{\sim}$
e não um conjunto importado de fora: os coeficientes que aqui se usam são os do
degrau que fecha a face do \emph{oposto}. Na base, $S$ é a matriz
\textbf{subdiagonal de uns}, $S_{ij}=[\,i=j+1\,]$. A álgebra de incidência de
$(I,\le)$ representa-se por matrizes triangulares inferiores, e nela vivem os
\emph{mesmos} $\zeta$ e $\mu$:
\[
\zeta=\sum_{j\ge0}S^{j}\quad(\text{a triangular de uns}),\qquad \mu=1-S ,
\]
com o produto $(f*g)(u,t)=\sum_{u\le v\le t}f(u,v)\,g(v,t)$.
\end{definicao}

\begin{lema}[$S$ é nilpotente, e é um bloco de Jordan]\label{fis:lem:nilp}
$S^{\dim W}=0$, e nenhum expoente menor a anula. O único valor próprio de $S$ é
$0$ e $\dim\ker S=1$: $S$ \textbf{não} é diagonalizável.
\end{lema}

\begin{proof}
$S^{j}$ desloca $j$ posições, logo $S^{j}_{ik}=[\,i=k+j\,]$, que é não nula
enquanto $j<\dim W$ e nula depois. Sendo triangular com diagonal nula, o
polinómio característico é $\lambda^{\dim W}$, e o núcleo é gerado por $\delta_0$
apenas.
\end{proof}

\medskip\noindent\textbf{É daqui que sai tudo o resto}, e sem nenhuma noção de
convergência: uma matriz nilpotente torna $1-S$ invertível, com inversa a
\emph{série de Neumann finita} $\sum_{j<\dim W}S^{j}$ --- que é exatamente
$\zeta$. Compare-se com a §\ref{fis:transf}: lá os operadores comutam e
\emph{diagonalizam}, e a inversão existe se e só se nenhum valor próprio se
anula; aqui o operador é um bloco de Jordan e não diagonaliza, mas a inversão
existe \textbf{sempre}. Duas álgebras lineares, dois motivos de invertibilidade
--- e as duas operações da aranha, uma em cada.

\begin{teorema}[$\mu=\zeta^{-1}$ na ordem total]\label{fis:thm:mu}
Na ordem total, a inversa de $\zeta$ é
\[
\mu(t,t)=1,\qquad \mu(t-1,t)=-1,\qquad \mu(u,t)=0\ \text{nos restantes},
\]
isto é, a \emph{diferença finita}: $(b*\mu)(t)=b(t)-b(t-1)$.
\end{teorema}

\begin{proof}
$(\zeta*\mu)(u,t)=\sum_{u\le v\le t}\mu(v,t)$. Se $u=t$ a soma tem um termo,
$\mu(t,t)=1$. Se $u<t$, os únicos termos não nulos são $v=t-1$ e $v=t$, ambos no
intervalo, e somam $-1+1=0$. Logo $\zeta*\mu=\delta_0$, e $\mu=\zeta^{-1}$.
\end{proof}

\begin{teorema}[o ciclo é $\zeta$; a recuperação é $\mu$]\label{fis:thm:zetamu}
Fixe-se $x\in X$ e ponha-se $a_x(t)=\mathbf{1}_{\{\pi(t)=x\}}$. Então:
\begin{enumerate}\itemsep3pt
\item[(1)] \textbf{A escrita é a convolução com $\zeta$:} $G_t(x)=(a_x*\zeta)(t)$.
A cláusula~(3) do Teor.~\ref{fis:thm:mult} \emph{é} esta identidade.
\item[(2)] \textbf{A recuperação é a deconvolução com $\mu$:}
$a_x(t)=G_t(x)-G_{t-1}(x)$. Conhecida a família $\{G_t\}$, a trajetória inteira
volta --- sem que nada dela tenha sido guardado.
\item[(3)] \textbf{O levantamento é a mesma convolução, na fibra:}
$k(i)=G_i(\pi(i))$ é a acumulação de $a_{\pi(i)}$ ao longo da fibra, e a sua
diferença finita vale $1$ em cada visita. Isto \emph{é} $\tilde G\equiv1$ do
Teor.~\ref{fis:thm:levant}, dito na outra linguagem.
\end{enumerate}
\end{teorema}

\begin{proof}
(1) $G_t(x)=\lvert\{u\le t:\pi(u)=x\}\rvert=\sum_{u\le t}a_x(u)$, que é
$(a_x*\zeta)(t)$ pela Def.~\ref{fis:def:incid}.

(2) Aplique-se $\mu=\zeta^{-1}$ (Teor.~\ref{fis:thm:mu}) a (1):
$(G_\bullet(x)*\mu)=(a_x*\zeta*\mu)=a_x$. Explicitamente é a diferença
$G_t(x)-G_{t-1}(x)$, que vale $1$ exatamente quando $\pi(t)=x$.

(3) Sejam $i_1<\dots<i_g$ a fibra de $x$. Por (1) com $t=i_r$,
$k(i_r)=G_{i_r}(x)=\sum_{u\le i_r}a_x(u)=r$. A diferença finita da sucessão
$(r)_{r=1..g}$ é constante igual a $1$; e um levantamento com valor $1$ em cada
ponto é um levantamento com $\tilde G\equiv1$.
\end{proof}

\begin{corolario}[é a mesma inversão da árvore dos divisores]\label{fis:cor:mobius}
Troque-se a ordem $(I,\le)$ pela ordem por \emph{divisibilidade}. A mesma
Def.~\ref{fis:def:incid} dá $g=f*\mathbf{1}$, isto é $g(n)=\sum_{d\mid n}f(d)$, e
a inversa de $\mathbf{1}$ é a função de Möbius: $f=g*\mu$. A álgebra não muda ---
$\zeta$ acumula, $\mu$ desacumula --- e só muda a ordem sobre a qual se acumula.
\end{corolario}

\begin{proof}
A Def.~\ref{fis:def:incid} não usa que a ordem seja total: usa apenas que os
intervalos $[u,t]$ sejam finitos, o que a divisibilidade cumpre. Com essa ordem,
$\mathbf{1}$ é a função $\zeta$ e a sua inversa na álgebra é, por definição, a
função de Möbius.
\end{proof}

\medskip\noindent\textbf{Duas convoluções, dois papéis.} A desta secção é a da
\emph{ordem} $(I,\le)$, e serve para o directo e o inverso da própria aranha. A
da §\ref{fis:transf} é a do \emph{grupo}, e serve para compor e decompor
\emph{campos} de percursos diferentes. Uma acumula ao longo do \emph{domínio}
$I$, a outra soma sobre o \emph{contradomínio} $X$.

\medido{\code{tests/aranha\_n.c} §AN27: $S^{12}=0$ com $\mathbf{12}$ mínimo,
$\sum S^j$ a ser a triangular de uns, $\dim\ker S=1$. §AN15: $\mu=\zeta^{-1}$ em
$\mathbf{4096}$ pares; o directo e o inverso com resíduo $\mathbf{0}$ sobre a
fibra mais funda que a realização tem; e a mesma inversão na árvore dos
divisores. §AN18: $(1-S)\zeta=\mathrm{id}$.}

\section{A estrutura linear: o espaço, o dual, a transformada}\label{fis:transf}

\begin{definicao}[os dois andares]\label{fis:def:doisandares}
\textbf{Em baixo}, o espaço onde a trajetória anda: $X$ como espaço vetorial
sobre $B$, de dimensão $m$, com a base $\{e_k=2^{k}\}$ do
Teor.~\ref{fis:thm:base} --- e as direções de aresta do hipercubo \emph{são} essa
base. \textbf{Em cima}, o espaço das funções sobre ele, que é um \emph{módulo
livre} sobre o \textbf{segundo degrau} $\Zc$, e não um espaço vetorial, porque
$\Zc$ não é corpo. E $\Zc$ não é importado: é o que a §\ref{fis:escada} constrói,
o degrau que fecha a face do \textbf{oposto}. Tem soma, tem opostos e não tem
inversos, e é exatamente isso que aqui se usa:
\[
V=\Zc^{X},\qquad \operatorname{rank}_{\Zc}V=\abs{X}=2^{m},
\qquad \text{base } \{\delta_x\}_{x\in X}.
\]
O campo $G$ é um \textbf{vetor} de $V$, e a sua escrita na base canónica
$G=\sum_x G(x)\,\delta_x$ diz que os valores de $G$ \emph{são} as suas
coordenadas.
\end{definicao}

\begin{definicao}[convolução]\label{fis:def:conv}
A operação de $X$ estende-se a $V$ por bilinearidade a partir da base:
\[
\delta_a*\delta_b=\delta_{a\oplus b}
\qquad\Longrightarrow\qquad
(f*g)(y)=\sum_{x}f(x)\,g(x\oplus y).
\]
Com este produto, $V$ é o \textbf{anel de grupo}, comutativo e com unidade
$\delta_0$.
\end{definicao}

\begin{lema}[os operadores de convolução comutam]\label{fis:lem:comuta}
Para cada $f\in V$ o operador $C_f(g)=f*g$ é linear, e
$C_fC_g=C_gC_f=C_{f*g}$.
\end{lema}

\begin{proof}
Linearidade e comutatividade herdam-se do produto, que é comutativo e
associativo por $X$ o ser. Explicitamente,
$C_fC_g(h)=f*(g*h)=(f*g)*h=C_{f*g}(h)$.
\end{proof}

\medskip\noindent É este lema que conduz ao resto: os operadores de convolução,
por comutarem, admitem diagonalização simultânea. A pergunta deixa de ser «que
transformada escolher» e passa a ser «qual é a base que diagonaliza a álgebra de
convolução».

\begin{definicao}[o dual e o caractere]\label{fis:def:dual}
$V^{*}=\operatorname{Hom}(V,\Zc)$, com o emparelhamento canónico. Para cada
$k\in X$, o \textbf{caractere} é
\[
\chi_k\colon X\to\{\pm1\},\qquad \chi_k(j)=(-1)^{\langle k,j\rangle},
\]
com $\langle k,j\rangle=$ paridade$(k\wedge j)$ --- a mesma forma bilinear que
torna $\{e_k\}$ ortonormal (Teor.~\ref{fis:thm:base}).
\end{definicao}

\begin{lema}[ortogonalidade]\label{fis:lem:orto}
$\chi_k(j\oplus j')=\chi_k(j)\chi_k(j')$, e
$\sum_k \chi_k(j)\chi_k(j')=2^{m}\,\delta_{j,j'}$.
\end{lema}

\begin{proof}
A primeira: $\langle k,j\oplus j'\rangle\equiv\langle k,j\rangle+\langle k,j'\rangle
\pmod 2$, porque um bit comum a $k$ e $j\oplus j'$ está em exatamente um dos dois
quando não está em ambos. Para a segunda, se $j=j'$ a soma é $2^{m}$; se
$j\neq j'$, ponha-se $d=j\oplus j'\neq0$ e fixe-se um bit $b$ com $d_b=1$:
emparelhando $k$ com $k\oplus e_b$ vem $\chi_{k\oplus e_b}(d)=-\chi_k(d)$, e a
soma cancela aos pares.
\end{proof}

\begin{teorema}[os caracteres são a base própria comum]\label{fis:thm:proprios}
Para todos $f\in V$ e $k\in X$,
\[
\boxed{\;C_f\,\chi_k \;=\; \widehat{f}(k)\,\chi_k,\qquad
\widehat{f}(k):=\sum_j f(j)\chi_k(j)\;}
\]
--- cada $\chi_k$ é vetor próprio de \emph{todos} os operadores de convolução ao
mesmo tempo, e o valor próprio de $C_f$ em $\chi_k$ é a coordenada
$\widehat{f}(k)$. Os $2^{m}$ caracteres são linearmente independentes, logo
formam base de $V$.
\end{teorema}

\begin{proof}
$(f*\chi_k)(y)=\sum_x f(x)\chi_k(x\oplus y)=\chi_k(y)\sum_x f(x)\chi_k(x)$ pelo
Lema~\ref{fis:lem:orto}, que é $\widehat{f}(k)\chi_k(y)$. A independência sai da
ortogonalidade: uma família ortogonal de vetores não nulos é livre.
\end{proof}

\begin{definicao}[a transformada]\label{fis:def:transf}
$\mathcal{F}(f)_k=\widehat{f}(k)=\sum_x f(x)\chi_k(x)$. Na base canónica a sua
matriz é $H_{kx}=\chi_k(x)$, com entradas $\pm1$: é a \textbf{avaliação nos
caracteres}, inteira. E ela é a mudança de base \emph{não normalizada}: as
coordenadas de $f$ na base $\{\chi_k\}$ são $c=2^{-m}\mathcal{F}(f)$.
\end{definicao}

\begin{teorema}[inversa, involutividade e Parseval]\label{fis:thm:H}
$H^{2}=2^{m}\,\mathrm{Id}$; logo $\mathcal{F}$ é invertível com
$\mathcal{F}^{-1}=2^{-m}\mathcal{F}$, e
\[
\mathcal{F}\circ \mathcal{F}=2^{m}\cdot\mathrm{id},\qquad
\lVert \mathcal{F}(f)\rVert^{2}=2^{m}\lVert f\rVert^{2}.
\]
\end{teorema}

\begin{proof}
$(H^2)_{kj}=\sum_x\chi_k(x)\chi_x(j)$. Como $\langle k,x\rangle=\langle
x,k\rangle$, isto é $\sum_x\chi_x(k)\chi_x(j)=2^{m}\delta_{kj}$ pelo
Lema~\ref{fis:lem:orto}. A norma sai da mesma soma.
\end{proof}

\begin{teorema}[a dobra na norma]\label{fis:thm:dobranorma}
Seja $P$ o número de pares $i<j$ com $\pi(i)=\pi(j)$. Então
\[
\sum_x G(x)^2=\abs{I}+2P,
\qquad
\lVert \mathcal{F}(G)\rVert^{2}=2^{m}\bigl(\abs{I}+2P\bigr),
\]
e $\pi$ é injetiva se e só se $\lVert \mathcal{F}(G)\rVert^{2}=2^{m}\abs{I}$.
\end{teorema}

\begin{proof}
$\sum_x G(x)^2=\lvert\{(i,j):\pi(i)=\pi(j)\}\rvert$ contando pares ordenados, que
se parte na diagonal ($\abs{I}$) e nos $2P$ restantes. A segunda igualdade é
Parseval (Teor.~\ref{fis:thm:H}), e a equivalência é $P=0$.
\end{proof}

\medskip\noindent A dobra --- definida sobre a trajetória --- lê-se num único
número calculado sobre o campo, \textbf{sem visitar a trajetória}. E cada
coordenada $\widehat{G}(k)$ tem leitura própria: por
$\mathcal{F}(G)_k=\sum_t\chi_k(\pi(t))$, é o \textbf{saldo} da trajetória no
\emph{corte} que $\chi_k$ define no grafo --- a bipartição dos vértices que
alterna nas direções de $k$, com $\abs{k}\,2^{m-1}$ arestas cortadas.

\begin{corolario}[a adjacência é uma convolução, e daí o seu espectro]\label{fis:cor:modos}
Seja $\alpha=\sum_{i}\delta_{e_i}$ a indicadora das direções de aresta. Então o
operador de adjacência da Def.~\ref{fis:def:ops} é $A=C_\alpha$, e portanto
\[
A\,\chi_k=\widehat{\alpha}(k)\,\chi_k=\bigl(m-2\abs{k}\bigr)\,\chi_k ,
\]
com $\abs{k}$ o peso de Hamming: o valor próprio é, em cada vértice, o número de
arestas em que $\chi_k$ \emph{concorda} menos aquele em que \emph{discorda}.
\end{corolario}

\begin{proof}
$(\alpha*f)(y)=\sum_i f(e_i\oplus y)$, que é a soma sobre os vizinhos de $y$:
logo $A=C_\alpha$. Pelo Teor.~\ref{fis:thm:proprios} o valor próprio é
$\widehat{\alpha}(k)=\sum_i\chi_k(e_i)=\sum_i(-1)^{k_i}$, que vale $-1$ nas
$\abs{k}$ direções de $k$ e $+1$ nas outras.
\end{proof}

\medskip\noindent Os caracteres não foram escolhidos, e agora vê-se duas vezes
porquê: são a base própria comum de \emph{todas} as convoluções, e a adjacência
do grafo é uma delas.

\begin{teorema}[convolução, deconvolução, e onde degenera]\label{fis:thm:conv}
$\mathcal{F}(f*g)=\mathcal{F}(f)\cdot \mathcal{F}(g)$ ponto a ponto. Se
$\widehat{f}(k)\neq0$ para todo $k$, então $C_f$ é invertível. Se
$\widehat{f}(k_0)=0$ para algum $k_0$, então $C_f$ é singular:
$\chi_{k_0}\in\ker C_f$, e dois $g$ que difiram por um múltiplo de $\chi_{k_0}$
dão o mesmo produto.
\end{teorema}

\begin{proof}
O produto é o Teor.~\ref{fis:thm:proprios} lido em coordenadas: na base própria,
um operador que age por multiplicação escalar compõe-se multiplicando escalares.
A inversão é diagonal, e existe exatamente quando nenhum valor próprio é nulo. O
núcleo é gerado pelos $\chi_k$ com $\widehat{f}(k)=0$.
\end{proof}

\medskip\noindent\textbf{E na forma polar isto lê-se de outra maneira.}
Escrevendo cada caractere como $r\,e^{\partial\theta}$, o módulo é
\emph{constante}:
\[
\abs{\chi_k(j)}=1\ \text{para todos os }k,j,
\qquad
\theta_k(j)\in\{0,\pi\},\qquad
\theta_k(j)=\pi\cdot\langle k,j\rangle .
\]
Ou seja: \emph{o módulo não carrega nada, e a fase carrega tudo} --- e a fase é
\textbf{um bit}, a paridade $\langle k,j\rangle$. Lido pelas duas faces: aqui a
face multiplicativa está \emph{degenerada} --- todos os módulos valem $1$ --- e
resta a aditiva sozinha, que é a razão de tudo neste andar ser ou-exclusivo. É o
caso extremo da divisão do Teor.~\ref{fis:thm:BI} entre a intensidade e o sinal:
aqui a intensidade é sempre a mesma, e o que distingue dois caracteres é só
\emph{para onde apontam}.

\medskip\noindent\textbf{E a forma polar é a separação das duas faces.} Não é uma
notação conveniente: é o Teor.~\ref{fis:thm:duas}(4) escrito num sítio onde as
duas faces convivem. O \emph{módulo} vive na face multiplicativa --- multiplica
--- e a \emph{fase} vive na aditiva --- soma:
\[
r_{1}e^{i\theta_{1}}\cdot r_{2}e^{i\theta_{2}}
=\underbrace{(r_{1}r_{2})}_{\text{face }\cdot}\;
 e^{\,i\overbrace{(\theta_{1}+\theta_{2})}^{\text{face }+}} ,
\]
e a exponencial é o \textbf{dicionário} entre elas: leva a soma no produto, que é
exatamente levar uma face na outra.

\medskip\noindent\textbf{E aqui $\mu$ coincide com $\zeta$.} Na §\ref{fis:zetamu}
eles são operações distintas, mas isso é na ordem $(I,\le)$. Nesta ordem a fase é
um bit e o $\oplus$ é \emph{involutivo}, pelo que desfazer e fazer são a
\emph{mesma} operação; e a transformada é o seu próprio inverso a menos da
escala. O par $(\zeta,\mu)$ não é propriedade da álgebra: é da \textbf{ordem}.
\[
\boxed{\;\text{o par }(\zeta,\mu)\text{ não é propriedade da álgebra: é da ORDEM.}\;}
\]
Na ordem total, $\zeta$ acumula e $\mu$ é a diferença finita, e são
\emph{distintos}. Numa ordem de \textbf{expoente $2$} --- onde $x\oplus x=0$ ---
o operador é a sua própria inversa e a involução \textbf{identifica} os dois.
Não é a álgebra a mudar de lei: é a mesma lei sobre uma ordem em que fazer e
desfazer coincidem.

\begin{center}\small
\begin{tabular}{@{}lll@{}}
\toprule
 & \textbf{o índice} (ordem linear) & \textbf{o espectro} (expoente $2$)\\
\midrule
$\zeta$ & acumular & a projeção espectral\\
$\mu$ & a diferença finita & a \emph{mesma} projeção\\
o par & $\zeta\neq\mu$ & $\boxed{\zeta=\mu}$\\
a trajetória & descer deixando rasto, subir a recolher & ir e vir pelo mesmo operador\\
a escala & unitária & $\abs{X}$\\
\bottomrule
\end{tabular}
\end{center}

\noindent A última linha não é detalhe: é onde se paga a identificação. Na ordem
linear a volta fecha \emph{exatamente} --- $\mu(\zeta(b))=b$, sem factor --- e na
de expoente $2$ ela fecha a menos de $\abs{X}$. \emph{Fazer e desfazer tornarem-se
a mesma operação não sai de graça: sai com uma norma pendurada.} E o que colapsou
foi o \textbf{módulo}: com a fase restrita a $\{0,\pi\}$ o vetor degenera num
vetor de \emph{sinais}, e a informação passa a ser orientação, e não tamanho.

\medskip\noindent A degenerescência não é rara e tem nome próprio: é $C_f$ ter
núcleo. O campo de um percurso que visita duas células uma vez cada tem
$\widehat{f}(k)=1+\chi_k(1)$, nulo em \emph{metade} dos $k$ --- o núcleo tem
dimensão $2^{m-1}$.

\medido{\code{tests/aranha\_n.c} §AN26: $\delta_a*\delta_b=\delta_{a\oplus b}$;
$C_fC_g=C_gC_f$; $C_f\chi_k=\widehat f(k)\chi_k$;
$\mathcal{F}\circ\mathcal{F}=2^{m}\cdot\mathrm{id}$; e
$\mathcal{F}(\delta_a)_k=\chi_k(a)$ com a volta exata. §AN12: as coordenadas do
campo na base canónica; a dobra na norma por \textbf{três} caminhos, com o
percurso injetivo a dar excesso zero; convolução diagonalizada e deconvolução com
resíduo $\mathbf{0}$; e a degenerescência exibida. §AN17:
$A\chi_k=(m-2\abs{k})\chi_k$ exato em inteiros, e o corte com $\abs{k}2^{m-1}$
arestas. §AN28: $c=2^{-m}\mathcal{F}(f)$ a recompor $f$ nos $256$, com o passo
sem a normalização a bater só nos zeros de $f$.}

\section{A escada, andar a andar}\label{fis:escada}

\noindent A escada não parte da lista --- parte da \emph{distinção}. O primeiro
andar é o do Teor.~\ref{fis:thm:BI}: as palavras em $B$, com a soma por
transporte e o encaixe por prefixos. Daí em diante, cada degrau tem por domínio o
que o anterior produziu, de modo que a escada inteira se escreve a partir de $B$
e de mais nada:
\[
\boxed{
\begin{array}{lll}
X_1 & = & I,\ \text{as palavras em } B\\[2pt]
X_2 & = & X_1^{\,2}/{\sim}, \qquad (a,b)\sim(c,d)\iff a+d=b+c\\[2pt]
X_3 & = & \bigl(X_2\times X_2^{\ast}\bigr)/{\sim}, \qquad (p,q)\sim(r,s)\iff ps=qr\\[2pt]
X_4 & = & \text{os cortes de } X_3
\end{array}}
\]
com $X_2^{\ast}$ os elementos invertíveis, e as três realizações a serem
$(a,b)\mapsto a-b$, $(p,q)\mapsto p/q$ reduzido, e $q\mapsto$ o cilindro que o
contém. Nenhuma precisa de um símbolo importado: o domínio de cada uma é o
contradomínio da anterior, e o único dado de toda a cadeia é a distinção da
Def.~\ref{fis:def:B}.

\medskip\noindent\textbf{E o primeiro andar não é um quociente.} $X_1$
constrói-se \emph{sem identificar nada} --- na dobra a segunda cópia não é colada
à primeira, é \emph{multiplicada} --- ao passo que $X_2$, $X_3$ e $X_4$ são
construídos \emph{por} quociente.

\begin{teorema}[os degraus por quociente são realizações]\label{fis:thm:escada}
Nos quatro casos abaixo, $\pi$ é uma realização e $\sum_x G(x)=\abs{I}$. Nos três
últimos --- os que se constroem \emph{por} quociente --- a fibra $\pi^{-1}(x)$
\emph{é} a classe de equivalência que define o andar. No primeiro não há
quociente a identificar, e a fibra é a da soma:
\begin{center}\footnotesize
\begin{tabular}{@{}lllll@{}}
\toprule
$\pi$ & a aplicação & a fibra é & o repr. $k=1$ & a \textbf{face} que fecha \\
\midrule
a soma em $X_1$ \emph{(sem quociente)} & $(x,y)\mapsto x\oplus y$ & a fibra da soma & $x\oplus z$ & o \textbf{neutro} $0$ \\
$\pi_2$: $X_1^{2}/{\sim}$ & $(a,b)\mapsto a-b$ & $a+d=b+c$ & $a=0$ ou $b=0$ & o \textbf{oposto}: $\partial x+x=0$ \\
$\pi_3$: as frações & $(p,q)\mapsto p/q$ reduzido & $ps=qr$ & o irredutível & o \textbf{inverso}: $\partial x\cdot x=1$ \\
$\pi_4$: os cortes & $q\mapsto$ o corte de $p$ bits & a bola (§\ref{fis:ultra}) & o menor do cilindro & o \textbf{neutro} $1$ \\
\bottomrule
\end{tabular}
\end{center}
\end{teorema}

\begin{proof}
Em cada linha, $\pi$ está definida em todo o domínio e o contradomínio tem
igualdade decidível, o que é a Def.~\ref{fis:def:objeto}; a conservação é o
Lema~\ref{fis:lem:conserva}. Nas três últimas linhas, $\pi(u)=\pi(v)$ é,
respetivamente, $a-b=c-d$ (isto é $a+d=b+c$), $p/q=r/s$ (isto é $ps=qr$), e ter o
mesmo prefixo de $p$ bits --- que são exatamente as relações que definem esses
andares. Na primeira, $\pi(u)=\pi(v)$ é $x\oplus y=x'\oplus y'$, relação legítima
mas não a que constrói o andar. O representante $k=1$ é o primeiro da fibra na
ordem do domínio, e a enumeração escolhida fá-lo coincidir com o canónico.
\end{proof}

\medskip\noindent\textbf{E a última coluna fecha o documento no seu princípio.}
As duas equações da Def.~\ref{fis:def:op} têm \textbf{quatro} peças, e não duas
--- duas \emph{operações} e dois \emph{neutros}:
\[
\boxed{\;
\underbrace{\partial x+x=0}_{+\ \text{e}\ 0}
\qquad\qquad
\underbrace{\partial x\cdot x=1}_{\times\ \text{e}\ 1}
\;}
\]
São essas as \textbf{faces da operação}, e a escada percorre-as. Não as resolve
ao mesmo tempo: resolve \emph{uma por andar}, e os dois duais saem da
\textbf{mesma troca} --- o par $(a,b)$ trocado por $\iota(a,b)=(b,a)$ induz o
\emph{oposto} em $X_2$ e o \emph{inverso} nas frações. E os dois neutros não
nascem na escada: nascem na Def.~\ref{fis:def:op}, porque é lá que as duas
equações os \emph{nomeiam}. Em uma frase:
\[
\boxed{\ \text{quocientar é esquecer a distinção; }G\text{ mede quantos elementos
foram esquecidos juntos.}\ }
\]

\medido{\code{tests/aranha\_n.c} §AN32: os quatro casos, com
$\sum G=\abs{I}$; nos três por quociente a fibra a ser a classe, e o $k=1$ a ser
o canónico --- no degrau das frações, literalmente a fração irredutível.}

\colunas{Teor.~\ref{fis:thm:duas},~\ref{fis:thm:B},~\ref{fis:thm:BI},~\ref{fis:thm:coord}; Lema~\ref{fis:lem:conserva}; Teor.~\ref{fis:thm:mult},~\ref{fis:thm:mu},~\ref{fis:thm:proprios},~\ref{fis:thm:H},~\ref{fis:thm:escada}}
        {$B$, a lista $I$, o par $(m,\delta)$, o anel de grupo, a escada}
        {\code{tests/meio.c} §M3--§M7; \code{tests/aranha\_n.c} §AN12, §AN15--§AN18, §AN26--§AN28, §AN32--§AN36}

% ═══════════════════════════════════════════════════════════════════════════
\part{Topologia}
% ═══════════════════════════════════════════════════════════════════════════

\section{Espaços ultramétricos: a célula é a bola}\label{fis:ultra}

\noindent O Teor.~\ref{fis:thm:mult} não pediu métrica. Mas quando $X$ tem uma, a
relação entre $G$ e as bolas diz o que a métrica é.

\begin{definicao}[a árvore do encaixe]\label{fis:def:arvore}
Para $a\neq b$ de $p$ bits, seja $\operatorname{prof}(a,b)$ --- e não $\delta$,
que aqui é a base canónica e o Kronecker --- a profundidade da primeira
divergência: o menor $q\ge0$ tal que os bits de peso $p-1,\dots,p-q$ de $a$ e $b$
coincidem e o de peso $p-1-q$ difere. Ponha-se
\[
d(a,b)=\begin{cases}0,&a=b,\\[2pt] 2^{-\operatorname{prof}(a,b)},&a\neq b.\end{cases}
\]
A separação --- $d(a,b)=0$ se e só se $a=b$ --- é parte da definição, e não uma
consequência a verificar: com $\operatorname{prof}(a,a)=p$ viria
$d(a,a)=2^{-p}\neq0$, e não seria métrica.
\end{definicao}

\begin{lema}[$d$ é ultramétrica]\label{fis:lem:ultra}
$d(a,c)\le\max\{d(a,b),d(b,c)\}$ para quaisquer $a,b,c$.
\end{lema}

\begin{proof}
Se $a$ e $b$ concordam nos $\operatorname{prof}(a,b)$ bits mais altos e $b$ e $c$
nos $\operatorname{prof}(b,c)$, então $a$ e $c$ concordam nos primeiros
$\min\{\operatorname{prof}(a,b),\operatorname{prof}(b,c)\}$, logo
$\operatorname{prof}(a,c)\ge\min\{\operatorname{prof}(a,b),\operatorname{prof}(b,c)\}$.
Aplicando $t\mapsto2^{-t}$, que inverte a ordem, vem
$d(a,c)\le\max\{d(a,b),d(b,c)\}$.
\end{proof}

\begin{teorema}[as bolas particionam]\label{fis:thm:bolas}
Seja $(X,d)$ ultramétrico e $\mathcal{B}(a,r)=\{x: d(a,x)\le r\}$ a bola ---
caligráfico, para não colidir com o binário $B$ da Def.~\ref{fis:def:B}. Então:
\begin{enumerate}\itemsep3pt
\item[(1)] \textbf{Todo ponto é centro:} se $b\in \mathcal{B}(a,r)$ então
$\mathcal{B}(b,r)=\mathcal{B}(a,r)$.
\item[(2)] \textbf{Partição:} duas bolas de raio $r$ ou coincidem ou são
disjuntas. Em particular, para cada $r$ as bolas de raio $r$ formam uma partição
de $X$.
\item[(3)] Fixado $r$, seja $\pi_r(i)$ a bola de raio $r$ que contém $i$. Então
$\pi_r$ é uma realização no sentido da Def.~\ref{fis:def:objeto}, a sua célula
\emph{é} a bola, e $G(x)$ é o número de índices dentro dela.
\end{enumerate}
\end{teorema}

\begin{proof}
(1) Seja $x\in \mathcal{B}(a,r)$. Então $d(b,x)\le\max\{d(b,a),d(a,x)\}\le r$,
logo $x\in \mathcal{B}(b,r)$ e $\mathcal{B}(a,r)\subseteq \mathcal{B}(b,r)$. Como
$b\in \mathcal{B}(a,r)$ implica $a\in \mathcal{B}(b,r)$, a inclusão contrária sai
do mesmo argumento com os papéis trocados.

(2) Se $c\in \mathcal{B}(a,r)\cap \mathcal{B}(b,r)$, por (1) tem-se
$\mathcal{B}(a,r)=\mathcal{B}(c,r)=\mathcal{B}(b,r)$. Que cobrem $X$ é imediato:
$a\in \mathcal{B}(a,r)$.

(3) A aplicação está bem definida por (2): há \emph{exatamente} uma bola de raio
$r$ contendo $i$. Quanto à igualdade no contradomínio, na árvore finita de $p$
bits cada bola é identificada pelo seu \textbf{prefixo}, e a igualdade das
células reduz-se à igualdade de prefixos. A afirmação sobre $G$ é a
Def.~\ref{fis:def:objeto}.
\end{proof}

\medskip\noindent\textbf{O que a ultramétrica é, e o que não é.} Não é uma
desigualdade triangular «mais apertada» de que se tire um número melhor: é a
propriedade de partição do Teor.~\ref{fis:thm:bolas}. Numa métrica euclidiana as
bolas de raio fixo cruzam-se parcialmente --- existem $a,b$ e pontos $x,y,z$ com
$x\in\mathcal{B}(a)\cap\mathcal{B}(b)$, $y\in\mathcal{B}(a)\setminus\mathcal{B}(b)$,
$z\in\mathcal{B}(b)\setminus\mathcal{B}(a)$ --- e nenhum ponto interior é centro
de todas as bolas que o contêm. É por isso que $\sum_x G(x)=\abs{I}$ não exige
cuidado nenhum aqui: \emph{não há sobreposição a haver}.

\medskip\noindent\textbf{E é a métrica do corte.} Dois pontos construídos como
cortes estão próximos nesta métrica quando \emph{concordam até fundo}: a
distância decai com a profundidade em que a expansão binária ainda coincide. É a
topologia que o encaixe induz.

\begin{teorema}[a vizinhança na árvore, e ela não é $2n$]\label{fis:thm:vizarvore}
As cláusulas (1)--(3) do Teor.~\ref{fis:thm:mult} correm sobre a árvore sem
alteração, porque só pedem igualdade. A cláusula~(4) pede uma \emph{vizinhança},
e a árvore dá uma: para $p\ge1$, cada bola de raio $2^{-p}$ tem \textbf{exatamente
um} irmão --- a bola com o mesmo prefixo de $p-1$ bits e o bit seguinte trocado ---
e as duas particionam a bola-mãe de raio $2^{-(p-1)}$. Logo
$\abs{\mathcal{V}(x)}=1$, e o custo dos passos de leitura e decisão é $\Theta(1)$
e não $\Theta(n)$.

\medskip\noindent\emph{A hipótese $p\ge1$ é necessária e não é decorativa}: em
$p=0$ a bola é a \textbf{raiz} --- o espaço inteiro --- e não tem irmão nenhum,
porque não há bola-mãe acima dela. A vizinhança da árvore é a de uma bola
\emph{não-raiz}, e o autómato nunca a pede na raiz porque aí não há para onde
descer.
\end{teorema}

\begin{proof}
Fixado $p$, uma bola de raio $2^{-p}$ é o conjunto dos pontos com um prefixo de
$p$ bits dado (Teor.~\ref{fis:thm:bolas}(2): as bolas particionam, e o prefixo é
a etiqueta da classe). Trocar o $p$-ésimo bit dá outro prefixo de $p$ bits com o
mesmo início de $p-1$; há um só, porque o bit é binário. A união das duas é
exatamente o conjunto dos pontos com o prefixo de $p-1$ bits, que é a bola-mãe.
\end{proof}

\medskip\noindent A vizinhança é o parâmetro: $2n$ na grade, $1$ na árvore. Os
outros três passos do ciclo são os mesmos nas duas.

\medido{\code{tests/aranha\_n.c} §AN6: as bolas particionam; todo ponto é centro;
e a euclidiana falha \emph{ambas} com o mesmo raio. §AN14: o custo contado ---
$2n$ leituras por passo, nem uma a mais, o autómato não varre.}

\section{A dobra vista por Euler}\label{fis:euler}

\begin{definicao}[percurso]\label{fis:def:percurso}
$\pi$ é um \textbf{percurso} no grafo quando $\pi(t+1)\in \mathcal{V}(\pi(t))$
para todo $t<N$, isto é, quando cada passo é uma \emph{aresta}. Nesse caso o
índice conta arestas: $\abs{I}=N+1$ vértices visitados e $N$ arestas percorridas.

\medskip\noindent Ser percurso é uma condição sobre o \emph{par} (realização,
grafo), e não uma propriedade da realização: na grade, a curva de Heighway é
percurso e a reta \emph{diagonal} não é --- o seu passo $(1,1)$ não é aresta ---
embora ambas sejam realizações perfeitamente legítimas.
\end{definicao}

\begin{teorema}[a dobra vista por Euler]\label{fis:thm:euler}
Seja $\pi$ um percurso conexo, $v=\abs{\operatorname{supp}G}$ o número de células
e $a$ o de arestas \emph{distintas} usadas. A característica do subgrafo
percorrido é $\chi=v-a$, e o número de ciclos independentes é $b_1=1-\chi$. Se
além disso \textbf{nenhuma aresta se repete} --- $a=N$ --- então
\[
b_1\;=\;D,\qquad D:=\abs{I}-v ,
\]
com $D$ o número de revisitas: \emph{cada dobra fecha um ciclo, e um só}.
\end{teorema}

\begin{proof}
Para um grafo conexo de dimensão $1$, $b_1=a-v+1$, donde $b_1=1-\chi$. Com
$a=N=\abs{I}-1$ vem $b_1=\abs{I}-1-v+1=\abs{I}-v=D$.
\end{proof}

\medskip\noindent A hipótese é necessária: a espiral quadrada de $1024$ passos
usa $\mathbf{4}$ arestas, e aí $b_1=1$ com $D=\abs{I}-v=1025-4=\mathbf{1021}$.
Sem ela, o que $G$ conta por \emph{vértice} deixa de ser o que a topologia conta
em \emph{ciclos}. E a reta axial dá $\chi=1$ com $b_1=0$: \emph{percorrer sem
dobrar é percorrer uma árvore}.

\medido{\code{tests/aranha\_n.c} §AN25: $b_1=D=\mathbf{335}$ no dragão; a espiral
a reusar arestas; e a reta a ser árvore.}

\section{Onde a dimensão entra}\label{fis:grau}

\noindent A cláusula~(4) --- \emph{sentir é ler $G$} --- pede uma
\emph{vizinhança}. É aqui, e só aqui, que a estrutura de $X$ é chamada.

\begin{definicao}[a grade e a sua vizinhança]\label{fis:def:grade}
Diz-se que $X$ é uma \textbf{grade de posto $n$} quando é gerado por $n$ direções
independentes $e_1,\dots,e_n$ --- a base canónica --- e cada uma se percorre nos
dois sentidos. A vizinhança de $x$ é então
$\mathcal{V}(x)=\{x\pm e_i : i=1,\dots,n\}$. É só disto que a construção precisa:
as direções e os dois sentidos.
\end{definicao}

\begin{teorema}[o $n$ entra num sítio só]\label{fis:thm:dim}
Numa grade de posto $n$ tem-se $\abs{\mathcal{V}(x)}=2n$ para todo $x$, e os $2n$
vizinhos são distintos. Além disso, as cláusulas (1)--(3) do
Teor.~\ref{fis:thm:mult} e os passos de escrita e decisão do ciclo valem palavra
por palavra em qualquer posto $n\ge1$, sem alteração nenhuma. A dimensão ocorre
num sítio único --- a cláusula~(4) --- e ocorre como $\abs{\mathcal{V}(x)}=2n$.
\end{teorema}

\begin{proof}
Os $2n$ pontos $x\pm e_i$ são distintos: $x+\varepsilon e_i=x+\varepsilon'e_j$
com $\varepsilon,\varepsilon'\in\{+1,-1\}$ dá $\varepsilon e_i=\varepsilon'e_j$,
e como os $e_i$ são livres isso força $i=j$ e $\varepsilon=\varepsilon'$. A
segunda afirmação é imediata da demonstração do Teor.~\ref{fis:thm:mult}, que
usou apenas a igualdade em $X$ --- e numa grade ela decide-se coordenada a
coordenada.
\end{proof}

\begin{teorema}[o fecho, e o único sítio onde a dimensão faz mais]\label{fis:thm:fecho}
O passo de fecho do ciclo é o rotor $R^{4}=\mathrm{id}$. Numa grade de posto $n$,
uma rotação de $90^{\circ}$ vive num \emph{plano} coordenado $(e_i,e_j)$, atuando
por $R_{ij}\colon(a,b)\mapsto(-b,a)$ nessas duas coordenadas e fixando as outras.
Então $R_{ij}^{4}=\mathrm{id}$ para todo par $i<j$, e o número de planos é
$\binom{n}{2}$: nenhum no posto $1$, um no posto $2$, três no $3$, seis no $4$.
\end{teorema}

\begin{proof}
Nas coordenadas $(i,j)$, $R_{ij}$ é a matriz
$\left(\begin{smallmatrix}0&-1\\1&0\end{smallmatrix}\right)$, cujo quadrado é
$-\mathrm{id}$ e cuja quarta potência é $\mathrm{id}$; as restantes coordenadas
não são tocadas. A contagem de planos é a de pares não ordenados de índices.
\end{proof}

\medskip\noindent Este é o único passo em que a dimensão faz mais do que contar
vizinhos. Nos outros ela entra só como $\abs{\mathcal{V}(x)}=2n$; no fecho, ela
decide \emph{quantos} rotores há.

\begin{teorema}[a cadeia não é isomorfa à grade \emph{como vizinhança}]\label{fis:thm:vizniso}
Sejam $X$ uma grade de posto $n$ e $I$ a régua, que é a grade de posto $1$. Se
$\varphi\colon X\to I$ é injetiva e leva vizinhos em vizinhos --- isto é,
$\varphi(\mathcal{V}(x))\subseteq\mathcal{V}(\varphi(x))$ para todo $x$ --- então
$n\le1$.
\end{teorema}

\begin{proof}
Pelo Teor.~\ref{fis:thm:dim}, $\abs{\mathcal{V}(x)}=2n$ em $X$ e os $2n$ vizinhos
são \emph{distintos}; em $I$, que é de posto $1$, vale
$\abs{\mathcal{V}(t)}=2$. Como $\varphi$ é injetiva, ela leva os $2n$ vizinhos
distintos de $x$ em $2n$ pontos distintos, todos dentro de
$\mathcal{V}(\varphi(x))$, que tem $2$ elementos. Logo $2n\le2$.
\end{proof}

\medskip\noindent\textbf{O que o teorema não diz: codificar não é linearizar.} A
objeção imediata é «mas eu serializo uma matriz num vetor». \emph{Sim, e sem
perder um único dado.} O Teor.~\ref{fis:thm:vizniso} não afirma que a régua seja
incapaz de \textbf{representar} os dados de uma matriz --- afirma que não os pode
representar \emph{preservando a vizinhança}:
\begin{center}\small
\begin{tabular}{@{}lll@{}}
\toprule
 & preserva & é possível em posto $n\ge2$?\\
\midrule
\textbf{codificação} (serialização) & a \emph{identidade} de cada célula & \textbf{sim} --- e é reversível\\
\textbf{linearização estrutural} & a identidade \emph{e} a vizinhança & \textbf{não} --- Teor.~\ref{fis:thm:vizniso}\\
\bottomrule
\end{tabular}
\end{center}
A serpentina é uma codificação perfeita: injetiva, e a sua inversa devolve o
ponto exato. O que ela não devolve é \emph{quem estava ao lado de quem}.
\[
\boxed{\ \text{codificar}\;\neq\;\text{preservar a dinâmica}.\ }
\]

\medskip\noindent\textbf{E convém dizer o que a curva de Peano é, para não a
chamar em vão.} A de Peano clássica é uma aplicação \emph{contínua}, sobrejetiva
e \textbf{não injetiva} --- o preenchimento paga com pontos do parâmetro que caem
no mesmo ponto do quadrado. Aqui não é esse o caso: o objeto é \emph{finito}, as
cardinalidades batem, e a serpentina é bijetiva. \emph{O obstáculo do contínuo
não é o obstáculo deste andar}; o deste é a vizinhança, e mede-se contando pares.

\begin{corolario}[o índice de uma matriz não se achata para $n\ge2$]\label{fis:cor:matriz}
Uma matriz é uma configuração sobre um índice discreto, e esse índice é a grade
de posto $2$ --- com as vizinhanças naturais $(i,j)\leftrightarrow(i\pm1,j)$ e
$(i,j)\leftrightarrow(i,j\pm1)$. Pelo Teor.~\ref{fis:thm:vizniso}, achatá-la numa
cadeia preservando essa vizinhança exige $2\cdot2\le2$, o que é falso. \emph{O
índice não se achata para $n\ge2$}; e no posto $1$ há um só sítio e nenhuma
direção transversal, pelo que a matriz $1\times1$ coincide com o escalar.
\end{corolario}

\begin{proof}
O índice $I_n\times I_n$ com aquelas vizinhanças é, pela
Def.~\ref{fis:def:grade}, a grade de posto $2$: as duas direções são
independentes e cada uma percorre-se nos dois sentidos, donde
$\abs{\mathcal{V}}=4$ num ponto interior. Aplicando o
Teor.~\ref{fis:thm:vizniso} com $n=2$ vem $2n\le2$, isto é $4\le2$: falso. Logo
não existe injeção que preserve a vizinhança. No posto $1$ o índice tem um
elemento, não há vizinho interior, e a hipótese do teorema não morde.
\end{proof}

\medskip\noindent E é o \emph{índice} que este corolário trata, não o conteúdo.
As duas camadas separam-se:
\[
\underbrace{I\times I}_{\text{\textbf{onde} --- a estante, discreta}}
\qquad\qquad
\underbrace{A\colon I\times I\to X}_{\text{\textbf{o que acontece} entre os onde}}
\]
e é a primeira que o teorema alcança. Nada disto depende do que $X$ seja.

\medido{\code{tests/aranha\_n.c} §AN7: $\abs{\mathcal{V}(x)}=2n$ e os vizinhos
distintos. §AN13: o ciclo em qualquer posto, e $0,1,3,6$ planos com
$R^4=\mathrm{id}$. §AN29: o grau $\mathbf{4}$ contra $\mathbf{2}$.}

\section{A realização contrária: dobra contra buraco}\label{fis:folga}

\noindent Pelo Teor.~\ref{fis:thm:palavra}, uma realização é uma \textbf{projeção
com perda}: leva a palavra na sua imagem e esquece o resto. \emph{Encher}
comprime o parâmetro no espaço, e o excesso de índices aparece como
\textbf{dobra}. A pergunta dual é o que acontece na direção contrária ---
\emph{contar}, do espaço para o parâmetro --- e a resposta não é «a inversa»,
porque não há uma: pelo Teor.~\ref{fis:thm:vizniso} o grau proíbe-a.

\begin{teorema}[a realização contrária, e o que ela paga]\label{fis:thm:contraria}
Seja $\Gamma_M$ o quadrado discreto e $q_0,\dots,q_{M^2-1}$ uma enumeração dos
seus pontos. A aplicação $\rho(q_t)=2t$ é \textbf{injetiva} e \textbf{não
sobrejetiva}: a sua imagem são os pares, e todos os ímpares ficam por usar.
Assim, as duas realizações trocam exatamente os dois defeitos:
\begin{center}\small
\begin{tabular}{@{}lccl@{}}
\toprule
 & injetiva & sobrejetiva & o que paga \\
\midrule
$\pi\colon I\to X$ \ (enche) & não & sim & \textbf{dobra}: $G>1$ \\
$\rho\colon X\to I$ \ (conta) & sim & não & \textbf{buraco}: índices por usar \\
\bottomrule
\end{tabular}
\end{center}
\end{teorema}

\begin{proof}
Injetividade: $t\mapsto2t$ é injetiva e a enumeração é bijetiva sobre
$\Gamma_M$. Não sobrejetividade: a imagem só contém índices pares.
\end{proof}

\begin{corolario}[a folga é uma só]\label{fis:cor:folga}
Com $\pi$ sobrejetiva sobre $X$ e $\rho$ injetiva sobre $X$, as duas perguntas
--- \emph{quantas partes da folha foram coladas nesta célula?} e \emph{quantos
pontos da régua ainda estão sem folha?} --- têm a mesma resposta:
\[
\sum_{x}\bigl(G(x)-1\bigr)\;=\;\abs{I}-\abs{X}\;=\;\abs{I}-\abs{\operatorname{im}\rho} .
\]
\end{corolario}

\begin{proof}
A soma da esquerda é $\sum_x G(x)-\abs{\operatorname{supp}G}$, que pelo
Lema~\ref{fis:lem:conserva} é $\abs{I}-\abs{X}$, já que $\pi$ é sobrejetiva. A da
direita é $\abs{I}-\abs{X}$ porque $\rho$ é injetiva com domínio $X$.
\end{proof}

\medskip\noindent Não é coincidência de números: é a \emph{mesma} folga vista
pelos dois mapas. Do lado de $\pi$ ela aparece como colagem em excesso; do lado
de $\rho$, como lugar por preencher. E não se podem pedir as duas coisas --- é o
que a §\ref{fis:grau} já disse: nenhuma injeção de células é injeção de
vizinhanças. Por isso $\rho$ não se chama aqui «curva do quadrado na reta»:
chama-se \textbf{realização injetiva do quadrado na régua}, e o nome carrega o que
ela preserva e o que não pode preservar.

\medskip\noindent E o par $(\pi,\rho)$ é a mesma \emph{retração} de sempre:
pode-se ter $\pi\circ\rho=\mathrm{id}_X$, mas nunca
$\rho\circ\pi=\mathrm{id}_I$ --- o que sobra na régua não tem de onde vir. Em uma
linha:
\[
\boxed{\ \text{perder uma coordenada}\iff\text{criar multiplicidade},\qquad
\text{recuperá-la}\iff\text{desdobrar}.\ }
\]

\begin{teorema}[a cadeia guarda o quadrado, sem dobra e sem buraco]\label{fis:thm:bijeccao}
Seja $\Gamma_M$ o quadrado discreto e $s$ a serpentina, $s(t)=(x,y)$ com
$y=\lfloor t/M\rfloor$ e $x=t\bmod M$ nas linhas pares, $x=M-1-(t\bmod M)$ nas
ímpares. Então $s$ é uma \textbf{bijeção}, e a sua inversa é a \emph{mesma regra}
lida ao contrário. Não há dobra ($G\equiv1$) nem índice por usar.
\end{teorema}

\begin{proof}
$s^{-1}\circ s=\mathrm{id}$ e $s\circ s^{-1}=\mathrm{id}$ por substituição
direta; como os dois conjuntos têm o mesmo cardinal e $s$ é injetiva, é bijetiva.
\end{proof}

\begin{teorema}[a enumeração da grade é a representação posicional]\label{fis:thm:enumera}
Seja $X=I_M^{\,n}$ a grade de lado $M$ e posto $n$, com $I_M=\{0,\dots,M-1\}$, e
$N=M^{n}$. A aplicação
\[
\boxed{\;\varphi(t)=\Bigl(\,t\bmod M,\ \bigl\lfloor\tfrac{t}{M}\bigr\rfloor\bmod M,\ \dots,\
\bigl\lfloor\tfrac{t}{M^{\,n-1}}\bigr\rfloor\bmod M\,\Bigr)\;}
\]
é uma \textbf{bijeção}, com inversa $\varphi^{-1}(x)=\sum_{k}x_k\,M^{\,k-1}$. E
cada componente é uma aplicação da \emph{mesma} divisão: $\varphi$ é o processo do
Teor.~\ref{fis:thm:coord} com $B$ substituído por $I_M$ --- a tupla das decisões,
lida por ordem.
\end{teorema}

\begin{proof}
Existência e unicidade do dígito são a divisão euclidiana: para cada $k$ há um só
par $(q,r)$ com $t=qM+r$ e $0\le r<M$, e iterar $n$ vezes esgota $t<M^{n}$.
$\varphi^{-1}\circ\varphi=\mathrm{id}$ é a reconstrução do número a partir dos
dígitos; $\varphi\circ\varphi^{-1}=\mathrm{id}$ é a unicidade da representação em
base $M$ com $n$ dígitos. Ambas as contagens dão $M^{n}$, e uma injeção entre
finitos do mesmo cardinal é bijeção.
\end{proof}

\medskip\noindent\textbf{E é aqui que a régua do operador fica construída, e não
apenas nomeada.} O Teor.~\ref{fis:thm:BI} deu a lista $I$ como a ordem das
palavras; o Teor.~\ref{fis:thm:coord} deu a tupla como a sucessão das decisões de
$\partial$; o Teor.~\ref{fis:thm:enumera} liga os dois --- \emph{a cadeia é a
leitura posicional da tupla}, e o dígito $k$ é a $k$-ésima decisão. A base
$B=\{0,1\}$ é o caso $M=2$, e nada muda quando $M$ cresce:
\[
\underbrace{t\in I_N}_{\text{onde o operador opera}}
\;\xrightarrow{\ \varphi\ }\;
\underbrace{(x_1,\dots,x_n)}_{\text{as decisões, por ordem}}
\;\longrightarrow\;
\underbrace{x\in X}_{\text{onde elas aterram}}
\]
--- e a régua não é um sítio para onde o objeto é empurrado: é a \emph{ordem em
que ele foi decidido}. É por isso que ela guarda tudo, e é por isso que ela não
guarda a vizinhança: \emph{a ordem das decisões não diz quais decisões são
próximas umas das outras}.

\begin{teorema}[a travessia é uma operação, e compor não acumula]\label{fis:thm:navega}
Sejam $\sigma,\tau,\rho$ representações de $X$ ligadas por bijeções efetivas, e
$d$ a ultramétrica da Def.~\ref{fis:def:arvore}. Defina-se
$d_x(\sigma,\tau)=d(\sigma(x),\tau(x))$ e
$D(\sigma,\tau)=\sup_{x} d_x(\sigma,\tau)$. Então $D$ é ultramétrica sobre as
representações:
\[
\boxed{\;D(\sigma,\rho)\;\le\;\max\bigl\{D(\sigma,\tau),\,D(\tau,\rho)\bigr\}\;}
\]
\emph{Logo uma cadeia de $L$ travessias tem o custo do seu \textbf{pior passo}},
qualquer que seja $L$.
\end{teorema}

\begin{proof}
Para cada $x$, o Lema~\ref{fis:lem:ultra} aplicado aos três endereços dá
$d_x(\sigma,\rho)\le\max\{d_x(\sigma,\tau),d_x(\tau,\rho)\}$. Tomando o máximo em
$x$ do lado esquerdo e majorando cada parcela do direito pelo seu próprio máximo
vem a desigualdade para $D$. A indução em $L$ é imediata, porque o máximo é
associativo.
\end{proof}

\medskip\noindent\textbf{E é a passagem de $d_x$ para $D$ que faz das
representações objetos geométricos.} Sem ela a desigualdade seria «para cada
ponto» e não diria nada sobre as representações como objetos; com ela, a
travessia passa a ser uma operação \emph{entre representações} e não uma conta
ponto a ponto. E daí a navegação é uma operação, e não uma escolha: como as bolas
\textbf{particionam} (Teor.~\ref{fis:thm:bolas}), fixar um prefixo \emph{é} fixar
uma bola, e descer a árvore bit a bit corta o espaço ao meio sem ambiguidade.
\[
\boxed{\;\text{navegar} \;=\; \text{descer por prefixo onde o prefixo é a bola}\;}
\]

\begin{proposicao}[a travessia é uma permutação de dígitos]\label{fis:prop:travessia}
Com $X=I_M^{\,n}$ e $\varphi$ do Teor.~\ref{fis:thm:enumera}, chame-se
\textbf{leitura} a $R_r(x)=\sum_{k}x_k\,M^{\,r(k)}$, com $r$ uma bijeção das
posições. Sejam $R=R_r$ e $S=R_s$ duas leituras. Então:
\begin{enumerate}\itemsep2pt
\item[(1)] \textbf{A travessia escreve-se.} $T=S\circ R^{-1}$ é a permutação
$\pi=s\circ r^{-1}$ \emph{aplicada às posições dos dígitos}. Não há nada a
procurar --- lê-se o dígito de uma posição e escreve-se noutra.
\item[(2)] \textbf{Conserva a medida.} Para todo $E$, $\mu(T(E))=\mu(E)$, com
$\mu$ a contagem --- e isso sai da \emph{bijetividade} e de mais nada.
\item[(3)] \textbf{O preço lê-se de $\pi$, não de $X$} --- \emph{e esta alínea
pede uma hipótese que as outras duas não pedem}. A ultramétrica da
Def.~\ref{fis:def:arvore} mede em \textbf{bits}, e aqui $q$ é índice de
\emph{dígito}; os dois só coincidem quando o dígito \emph{é} o bit, isto é
quando $M=2$. Sob essa hipótese,
\[
\boxed{\;D(R,S)=\begin{cases}0,&\pi=\mathrm{id},\\[2pt]
2^{-q},&q=\min\{i:\pi(i)\neq i\},\end{cases}\;}
\]
com $q$ o primeiro dígito em que as duas ordens divergem. Logo $D$ custa $O(n)$
comparações, e não $O(N)$.

\emph{Para $M=2^{d}$ cada dígito ocupa $d$ bits}, e a primeira divergência
binária ocorre em $q_{\mathrm{bit}}=d\,q_{\mathrm{dig}}$, donde
$D(R,S)=2^{-d\,q_{\mathrm{dig}}}$; reindexando por blocos de $d$ bits recupera-se
a forma acima. \emph{Para $M$ que não seja potência de dois} a fórmula não se
aplica como está: $M^{n}$ não é um número de palavras de bits, e é preciso ou
fixar um preenchimento binário dos endereços $0,\dots,M^{n}-1$, ou substituir o
bloco binário pela unidade de dígito da própria régua. As alíneas (1) e (2) não
dependem disto: valem para todo $M\ge2$.
\end{enumerate}
\end{proposicao}

\begin{proof}
(1) $R$ e $S$ são bijeções por composição, logo $T$ existe e é bijeção; e se
$a=R_r(x)=\sum_k x_k M^{r(k)}$, então $T(a)=R_s(x)=\sum_k x_k M^{\,\pi(r(k))}$,
porque $\pi(r(k))=s(k)$. \emph{Os dígitos ficam; mudam-lhes os destinos.}

(2) Uma bijeção de um conjunto finito em si mesmo leva um subconjunto de $k$
elementos num subconjunto de $k$ elementos --- e nada mais é preciso no discreto.

(3) Se $\pi=\mathrm{id}$ então $R=S$ e $D=0$. Caso contrário seja
$q=\min\{i:\pi(i)\neq i\}$. Para $i<q$ vale $\pi(i)=i$, donde a posição $i$
recebe em ambas as leituras o mesmo dígito de $x$, seja qual for $x$; logo as
representações coincidem nas posições $0,\dots,q-1$ e
$\operatorname{prof}(R(x),S(x))\ge q$, donde $D(R,S)\le2^{-q}$. E existe $x$ cujos
dígitos nas posições $q$ e $\pi^{-1}(q)$ diferem, o que é possível porque
$M\ge2$; para esse $x$ a divergência é exatamente em $q$, e o supremo é atingido.
\end{proof}

\medskip\noindent\textbf{E daqui sai o que a travessia \emph{não} conserva.} A
medida não distingue as réguas, mas a \textbf{localidade} sim. \emph{O que uma
mudança de régua move é o custo do salto, nunca o tamanho do conjunto}; e é por
isso que a escolha de representação é uma questão de operação e não de medida.

\medskip\noindent\textbf{E a álgebra não desce --- transporta-se.} Sendo $c$ a
codificação, é \emph{falso} que $c(A+B)=c(A)+c(B)$ --- e tem de ser, senão a soma
de matrizes seria a soma dos endereços e o corpo não teria estrutura própria
nenhuma. O que existe é a operação \textbf{conjugada}
\[
\boxed{\;A\boxplus B \;=\; c\bigl(c^{-1}(A)+c^{-1}(B)\bigr)\;}
\]
que está bem definida \emph{porque} $c$ é bijeção. \emph{O corpo cabe em $I$; a
aritmética de $I$ é que não é a dele.} Ficam então separadas as três coisas que
se confundem: \textbf{codificar} (objeto $\to$ endereço), \textbf{contar}
($I^{k}\leftrightarrow I$) e \textbf{transportar a álgebra} ($\boxplus$).

\medido{\code{tests/pgwire.c} §W123: o Lema~\ref{fis:lem:ultra} em
$\mathbf{262\,144}$ triplos; o prefixo a dar a linha em $1024/1024$ na ordem por
linhas alternadas contra $32/1024$ nas outras duas; a profundidade mínima entre
vizinhos horizontais, $5$ contra $0$; e a cadeia de três travessias com
$\mathbf{454}$ na soma contra $\mathbf{10}$ no máximo. §W128, com $M=16$ e $n=2$:
a travessia atinge $256/256$ endereços; $\mu(T(E))=\mu(E)$ em $12/12$ conjuntos;
os saltos máximos passam de $(16,1)$ para $(86,43)$ --- \emph{a medida a ficar e a
localidade a mudar}; a fórmula de (3) dá $q=1$ em $8$ comparações e o supremo
varrido sobre $256$ pontos dá o mesmo; e o gume: a leitura não bijetiva $a+b$ leva
$256$ pontos em $31$, e a medida \emph{cai}. §W122, em $M=32$: a serpentina dá
$(1,63)$ e a de Hilbert $(853,341)$ nos máximos, com $h+v$ a valer $32\,736$ e
$38\,936$ --- a troca de coordenadas conserva o total e redistribui-o.
\code{tests/cantor.c} §K8: $c(A+B)$ coincide com $c(A)+c(B)$ em \textbf{zero} de
$256$ pares, e $\boxplus$ fecha em $\mathbf{256}$ de $256$.
\code{tests/aranha\_n.c} §AN29 e §AN37: $\sum(G-1)=256$ e
$\abs{I}-\abs{\operatorname{im}\rho}=256$; contar injetiva com $256$ lugares por
usar; encher com $G\equiv2$; $\pi\circ\rho=\mathrm{id}$ e $\rho\circ\pi$ a falhar
em $256$ de $512$; e a serpentina com dobra $0$ e buraco $0$ em
$M=2,4,8,16,32$.}

\colunas{Lema~\ref{fis:lem:ultra}; Teor.~\ref{fis:thm:bolas},~\ref{fis:thm:euler},~\ref{fis:thm:dim},~\ref{fis:thm:vizniso},~\ref{fis:thm:navega}; Cor.~\ref{fis:cor:folga}}
        {a árvore, o dragão e a espiral, a serpentina, a travessia por dígitos}
        {\code{tests/aranha\_n.c} §AN6--§AN7, §AN13--§AN14, §AN25, §AN29, §AN37; \code{tests/pgwire.c} §W122--§W123, §W128}

% ═══════════════════════════════════════════════════════════════════════════
\part{Análise}
% ═══════════════════════════════════════════════════════════════════════════

\section{A outra face, e o seu meio}\label{fis:geo}

\noindent O Teor.~\ref{fis:thm:duas}(4) obriga a \emph{duas} composições e diz o
que $\partial$ é em cada uma. Até aqui desenvolveu-se sobretudo a primeira: o
meio $m=\tfrac{a+b}{2}$ e, por indução, os diádicos $k/2^{n}$. A segunda tem o
seu próprio meio, e ele não é matéria nova.

\begin{lema}[a face multiplicativa tem o seu meio]\label{fis:lem:geo}
Na face onde $\partial$ é o inverso, a dobra que troca os extremos é
\[
\boxed{\;\partial^{\times}x=\frac{ab}{x}\;}
\]
--- de facto $\partial^{\times}a=b$ e $\partial^{\times}b=a$, e o invariante é
$ab$. O seu único ponto fixo positivo é
\[
g_{\ast}=\sqrt{ab},
\]
a \textbf{média geométrica}. É a mesma frase do Teor.~\ref{fis:thm:duas}(2) lida
na face que o Teor.~\ref{fis:thm:duas}(4) obriga a existir: \emph{cada face tem
uma dobra, e cada dobra tem um meio}.

\medskip\noindent\textbf{A estrela não é decoração.} $g_{\ast}$ é o ponto fixo
\emph{exato}, e ele em geral não é da grade; o que o processo usa é o seu piso
$g$, definido a seguir. \emph{Os dois têm de andar com nomes diferentes}, porque
as identidades exatas valem para $g_{\ast}$ e o algoritmo corre com $g$.
\end{lema}

\begin{proof}
$\partial^{\times}$ troca porque $ab/a=b$; é involutiva porque $ab/(ab/x)=x$; e
$ab/x=x$ tem a única solução positiva $x=\sqrt{ab}$, que existe sempre que a raiz
se inverta --- a mesma condição que o $2$ era na face aditiva.
\end{proof}

\medskip\noindent\textbf{A grade, e porque ela tem de se dizer antes.} O
Teor.~\ref{fis:thm:duas}(6) já fixou o objeto onde tudo isto vive: ao fim de $n$
dobras os pontos são os $k/2^{n}$, uma \textbf{grade} de grão $u=2^{-n}$. As duas
dobras operam \emph{nela}: a aritmética cai-lhe dentro sozinha --- o meio de dois
pontos da grade é um ponto da grade seguinte --- mas a geométrica \textbf{não},
porque $g_{\ast}=\sqrt{ab}$ não é diádico.

\begin{definicao}[o piso, e o resíduo que ele deixa]\label{fis:def:piso}
$g$ é o \emph{maior elemento da grade cujo quadrado não passa} $ab$ --- o
\textbf{piso} de $g_{\ast}$ --- e é isso que faz o processo permanecer no objeto
que o constrói. Por construção
\[
\boxed{\;g\;\le\;g_{\ast}\;<\;g+u,\qquad\text{isto é}\qquad
0\;\le\;\underbrace{g_{\ast}-g}_{\text{o resíduo de grade}}\;<\;u\;}
\]
e daí $g^{2}\le ab$, com igualdade \textbf{exatamente} quando $g_{\ast}$ já
pertence à grade. Só com isto dito é que a largura $m-g$ é um \textbf{número
inteiro de grãos}, e é literalmente esse inteiro que vai decrescer.
\end{definicao}

\medskip\noindent\emph{E é aqui que as duas escalas se separam, e convém que se
vejam:} a face geométrica dá a \textbf{contração}, e a grade dá o
\textbf{quantum}. Nenhuma das duas sozinha termina o processo --- juntas,
terminam.

\begin{corolario}[a desigualdade das médias é a conservação da norma]\label{fis:cor:medias}
Escrevendo $a=m+\delta$ e $b=m-\delta$, o Teor.~\ref{fis:thm:BI} dá
$ab=m^{2}-\delta^{2}$. Logo, no ponto fixo \emph{exato},
\[
g_{\ast}^{2}=m^{2}-\delta^{2}\;\le\;m^{2},\qquad\text{isto é}\qquad
g_{\ast}\le m,
\]
com igualdade \emph{exatamente} quando $\delta=0$; e como o piso não sobe,
\[
g\;\le\;g_{\ast}\;\le\;m .
\]
Não é uma desigualdade nova --- é a norma conservada, lida como comparação.
\end{corolario}

\begin{proof}
$g_{\ast}^{2}=ab$ é o Lema~\ref{fis:lem:geo}, e $ab=m^{2}-\delta^{2}$ é o
Teor.~\ref{fis:thm:BI}; logo $g_{\ast}^{2}=m^{2}-\delta^{2}\le m^{2}$, porque
$\delta^{2}\ge0$, com igualdade se e só se $\delta=0$. Como $g_{\ast}$ e $m$ são
positivos, a raiz preserva a ordem e $g_{\ast}\le m$. E $g\le g_{\ast}$ é a
Def.~\ref{fis:def:piso}: o piso nunca excede o que arredonda.
\end{proof}

\section{As duas faces acabam numa igualdade: o corte}\label{fis:corte}

\begin{teorema}[o corte]\label{fis:thm:corte}
Batendo as duas faces alternadamente, $(a,b)\mapsto(m,g)$ com $m$ e $g$ os dois
meios:
\begin{enumerate}\itemsep3pt
\item[(1)] \textbf{encaixa}: $g\le g'\le m'\le m$, pelo Cor.~\ref{fis:cor:medias}
aplicado em cada passo --- os intervalos $[g_{n},m_{n}]$ são encaixantes;
\item[(2)] \textbf{colapsa}, e em duas escalas. No ponto fixo \emph{exato} a
identidade é limpa,
\[
m-g_{\ast}=\frac{\delta^{2}}{m+g_{\ast}},
\qquad\text{donde}\qquad
\delta'_{\ast}\le\frac{\delta^{2}}{2m}\le\frac{\delta}{2};
\]
e o que o processo mede é o piso, que acrescenta o resíduo da
Def.~\ref{fis:def:piso}:
\[
\boxed{\;m-g=\frac{\delta^{2}}{m+g_{\ast}}+\varepsilon,
\qquad 0\le\varepsilon<u\;}
\]
--- \emph{a geometria dá a contração quadrática; a grade dá o quantum $u$};
\item[(3)] \textbf{acaba}: sendo o objeto \emph{finito} --- a régua tem grão ---
a largura não se aproxima de zero, \textbf{chega} ao fundo. Ao fim de um número
\emph{contado} $N$ de batidas o encaixe pára, e pára de uma de duas maneiras:
\[
\begin{array}{ll}
\textbf{I. o ponto é da grade} & g_{\ast}\in\Gamma:\ \text{o processo produz}\ \boxed{m_{N}=g_{N}}\ \text{e o intervalo é um \textbf{ponto};}\\[3pt]
\textbf{II. o ponto não é da grade} & g_{\ast}\notin\Gamma:\ \text{o processo termina no intervalo mínimo } [g_N,\,g_N+u],\\
& \text{e o que sobra é exatamente um \emph{resíduo de grade}.}
\end{array}
\]
\emph{Nos dois casos a conta termina}; o que muda é se ela \textbf{produz} o
ponto ou o \textbf{cerca} (§\ref{fis:rastro});
\item[(4)] \textbf{um habitante, e somente um --- a menos do grão}: o que está em
todos os intervalos está no último, porque eles encaixam; e se $x$ e $y$ estão em
todos, então $\abs{x-y}\le m_{n}-g_{n}$ para todo $n$. No caso~I isso dá
$\abs{x-y}\le m_N-g_N=0$ e o habitante é \textbf{único}; no caso~II dá
$\abs{x-y}<u$, e a unicidade é \emph{à resolução da régua} --- que é tudo o que
uma régua de grão $u$ pode afirmar;
\item[(5)] \textbf{e no caso~I esse ponto é o ponto fixo comum às duas faces.}
$\partial^{+}$ fixa $m$ e $\partial^{\times}$ fixa $g_{\ast}$; os dois são
\emph{diferentes} enquanto $a\neq b$, e coincidem exatamente onde o processo
acaba. \emph{O corte não é um ponto que se ache entre as classes --- é o ponto
fixo que as duas faces passam a partilhar.}
\end{enumerate}
Isto é um \textbf{corte}, e é \emph{produzido} em vez de postulado --- e
produzido por uma conta que \emph{termina}, não por um limite:
\[
\boxed{\;\text{grade}\;\to\;\text{duas faces}\;\to\;\text{encaixe}\;\to\;
\text{descida bem-fundada}\;\to\;\text{colapso}\;\to\;\text{paragem}\;}
\]
\emph{Não converge para o ponto: termina.} E as duas fontes da terminação são
distintas e ambas necessárias:
\[
\boxed{\;\text{contração geométrica}\;+\;\text{quantização da grade}
\;\Longrightarrow\;\text{terminação discreta}\;}
\]
\end{teorema}

\begin{proof}
(1) $g\le m$ é o Cor.~\ref{fis:cor:medias}; $m'\le m$ porque
$m'=\tfrac{m+g}{2}$ com $g\le m$. E $g'\ge g$ mesmo com o piso: o ponto fixo
exato do passo seguinte é $g'_{\ast}=\sqrt{mg}\ge\sqrt{g\cdot g}=g$, e como $g$
\emph{já é da grade}, o piso de um valor $\ge g$ continua $\ge g$ --- o piso não
desce abaixo de um ponto da grade que ele já ultrapassou. Daí
$[g',m']\subseteq[g,m]$.

(2) A identidade exata é sobre $g_{\ast}$, e sai da conservação da norma:
\[
m-g_{\ast}=\frac{m^{2}-g_{\ast}^{2}}{m+g_{\ast}}=\frac{\delta^{2}}{m+g_{\ast}} ,
\]
com $g_{\ast}\ge0$ a dar o majorante $\delta^{2}/(2m)$, e $\delta\le m$ a dar
$\delta/2$. Para o piso, escreva-se $m-g=(m-g_{\ast})+(g_{\ast}-g)$: a primeira
parcela é a identidade acima e a segunda é o resíduo da Def.~\ref{fis:def:piso},
com $0\le g_{\ast}-g<u$. Donde a forma discreta do enunciado, com
$\varepsilon=g_{\ast}-g$. \emph{A contração é da face; o $\varepsilon$ é da
grade.}

(3) A descida é \textbf{bem-fundada}, e prova-se sem olhar para a raiz: se $g<m$
então $\tfrac{g+m}{2}<m$, e o maior ponto da grade abaixo de algo menor que $m$ é
no máximo $m-u$; logo
\[
g<m\;\Longrightarrow\;m'\le m-u .
\]
A largura desce pela ponta \emph{de cima}, um grão por batida no pior caso, faça
o piso da geométrica o que fizer. Sendo a largura um inteiro de grãos e
não-negativa, uma sucessão estritamente decrescente dela \textbf{pára} --- não
pode descer para sempre. \emph{Onde pára, a largura não pode ser maior que um
grão}, e isso deixa as duas saídas do enunciado: se a largura chegou a $0$, é o
caso~I e $m=g$; se ficou em $u$, é o caso~II, e o intervalo $[g,g+u]$ é o menor
que a régua distingue. Note-se o que a finitude faz aqui: ela é a
\emph{hipótese}, e não uma limitação --- é por o objeto ser finito que há
\textbf{terminação} em vez de limite, e por a descida ser estrita numa ordem
bem-fundada que isto é uma paragem e não uma convergência disfarçada.

(4) O último intervalo está contido em todos os anteriores por (1), logo o que
está em todos está nele. No caso~I a largura do último é $0$ e quaisquer $x,y$
comuns coincidem; no caso~II ela é $u$, e distam menos que um grão.
\emph{Nenhum axioma de completude é invocado} --- o argumento é o do pombal com a
folga já no fundo.

(5) Os dois pontos fixos são $\tfrac{a+b}{2}$ e $\sqrt{ab}$, e coincidem sse
$(a+b)^{2}=4ab$, isto é sse $a=b$ --- que é o caso~I de (3). Aí, com $a=b=p$, vem
$\partial^{+}p=p+p-p=p$ e $\partial^{\times}p=p^{2}/p=p$.
\end{proof}

\medskip\noindent\textbf{A cota, e o que a segunda face acrescenta.} De
$m'\le m-u$ sai imediatamente que, se a largura inicial tem $N$ grãos, não há
mais de $N$ descidas: a terminação é literalmente $N,\,N-1,\,\ldots,\,0$. É
\emph{essa} cota que prova que acaba, e ela não usa a raiz. O que a face
geométrica acrescenta é a \textbf{velocidade}: por (2) o colapso é quadrático.
\emph{Sem a primeira não se sabe que acaba, sem a segunda não se sabe que acaba
depressa} --- e só a primeira é uma demonstração de terminação.

\medskip\noindent\textbf{E é aqui que a segunda face se paga.} Com \emph{uma}
face só, o par colapsa num passo --- $(a,b)\mapsto(m,m)$ --- e a batida seguinte
não move nada: para continuar é preciso \textbf{escolher} um lado em cada nível,
que é a bissecção do Teor.~\ref{fis:thm:duas}(6). Aí o ponto fica
\emph{endereçado} por um caminho, a um bit por nível. Com as \emph{duas} faces
não se escolhe nada: o corte é \textbf{determinado} pelos dois extremos.
\[
\boxed{\;\text{uma face \textbf{endereça}; duas \textbf{produzem}}\;}
\]
É a frase do Teor.~\ref{fis:thm:duas}(4) --- \emph{uma composição não chega} ---
dita um andar acima, e o nome clássico do processo é a média
\textbf{aritmético-geométrica}.

\medido{\code{tests/meio.c} §M8: $\partial^{\times}$ troca e fixa $\sqrt{ab}$;
$g_{\ast}^{2}=ab=m^{2}-\delta^{2}\le m^{2}$ com igualdade só em $\delta=0$ e
desigualdade \emph{estrita} fora dele ($\mathbf{1890}$ pares); encaixa e halva em
todas as batidas; $\delta'\le\delta^{2}/(2m)$ na forma inteira ($\mathbf{975}$
pares); e o processo \textbf{acaba} --- $m=g$ em $\mathbf{5}$ a $6$ batidas, com o
valor a bater nos valores \emph{conhecidos} do processo, que são externos ao
ficheiro. A unicidade mede-se dos dois lados, porque um só não dizia nada: o
último intervalo tem \emph{um} habitante e o penúltimo tinha mais. E a
\textbf{boa fundação} varre-se: na grade $1..300$ inteira, $\mathbf{116\,241}$
passos, exige-se em cada um que $m'<m$ --- a cláusula que garante a paragem e que
não depende da raiz --- e com ela as duas que a sustentam (a geométrica nunca
desce, o encaixe nunca quebra). O CONTROLO diz porque a segunda face faz falta ---
com uma só, a segunda batida não move. Tudo em inteiros, com a raiz por piso. O
encaixotamento certificado está em \code{tests/agm.c} §A1--§A3, e a cota
bem-fundada dá $10^{9}$ onde as duas faces fazem \textbf{cinco}.}

\medskip\noindent\textbf{E prende-se uma face para obter a raiz.} Se em vez de
mover as duas se \textbf{fixar o produto}, o par $(r,\,n/r)$ tem $g=\sqrt{n}$
constante --- a face geométrica está presa --- e só a aritmética desce. O encontro
$m=g$ do Teor.~\ref{fis:thm:corte} é então exactamente $\sqrt{n}$, e o processo é
o \emph{mesmo}. \emph{Não é um algoritmo para a raiz: é este algoritmo com uma
face imobilizada.}

\medido{\code{lib/cifra.h} \code{raizi}, medido em \code{tests/pgwire.c} §W217.
Era um laço linear --- $\sqrt{n}$ voltas --- e passou a ser o par $(r,n/r)$. O
gume é o próprio laço linear, que é um terceiro caminho: em $-50..120000$ os dois
dão a \textbf{mesma} resposta, sem uma divergência, e as batidas caem de
$\mathbf{346}$ (o máximo do laço nessa gama, $27\,653\,185$ voltas ao todo) para
$\mathbf{5}$ --- \emph{o mesmo cinco} que o §M8 mede para a AGM de duas faces, o
que não é coincidência mas o mesmo processo.}

\section{O que o finito guarda de um ponto que não é seu}\label{fis:rastro}

\noindent O Teor.~\ref{fis:thm:corte} tem uma hipótese que é preciso ler: o ponto
onde o processo pára é um ponto \emph{da grade}. Há um segundo desfecho, e é ele
que dá a interpretação:
\[
\begin{array}{ll}
\text{o ponto \textbf{está} na grade} & \text{acaba em }m=g\text{; o ponto é produzido}\\
\text{o ponto \textbf{não está}} & \text{acaba com a largura no grão; as duas faces \emph{cercam-no}}
\end{array}
\]
No segundo caso não há nada a produzir, e nada se pede: um ponto pode não caber
em grade nenhuma. O que há é a \textbf{representação}, e ela é o par que o
Teor.~\ref{fis:thm:BI} já deu. De $L$ e $U$ --- as duas faces --- saem
\[
\boxed{\;
\underbrace{(m,\delta)}_{\text{o \textbf{rastro}}}
\qquad
\underbrace{m=\frac{L+U}{2}}_{\text{a representação \textbf{comprimida}}}
\qquad
\underbrace{\delta=\frac{U-L}{2}}_{\text{a \textbf{perda}, medida}}
\;}
\]
Três coisas distintas. O \textbf{rastro} é o \emph{par}: $(m,\delta)$ é $(L,U)$
escrito duas vezes, com resíduo zero --- nada se perdeu ao mudar de escrita. A
\textbf{compressão} é ficar só com $m$; e o que ela deita fora não é
desconhecido, é $\delta$ --- \emph{medido}. Dizer que $m$ é o que a representação
conserva quer dizer que $m$ é o seu \textbf{invariante}, e não que $m$ contenha
tudo: \emph{o que $m$ não contém tem nome, tamanho e sinal de aviso}.

\medskip\noindent\textbf{O centro é a leitura, e não uma média de conveniência.}
Se fosse arbitrário, valeria o mesmo que as bordas. Mede-se que não: no encaixe
da torre de duplicação o centro erra $\mathbf{6}$ a $\mathbf{13}$ vezes menos do
que qualquer das duas fronteiras, em três grades diferentes. \emph{As fronteiras
são o que se certifica; o centro é o que se sabe.}

\medskip\noindent\textbf{E é aqui que as quatro frases da porta de entrada se
lêem uma a uma:}
\[
\begin{array}{ll}
\text{o finito é \textbf{codificado}} & \text{o corpo guarda dois números da grade}\\
\text{o infinito é \textbf{caminhado}} & \text{cada batida é um passo, e nenhuma o alcança}\\
\text{a representação \textbf{comprime}} & \text{do ponto fica }m\\
\text{a dobra \textbf{mede o que foi perdido}} & \delta\ \text{é essa medida}
\end{array}
\]
$\delta$ não é um erro a lamentar: é o que a dobra devolve, e é por existir que se
sabe quanto ficou de fora. \emph{A memória existe onde a trajetória passou} --- e
o rastro é essa memória.

\medskip\noindent\textbf{E nada disto invoca o ponto.} O processo não recebe o
ponto como ingrediente e não o consulta: ele constrói o caminho, e o caminho
determina o ponto.
\[
\boxed{\;\text{não se postula o ponto; constrói-se o caminho que o determina}\;}
\]
Em uma linha, e cada seta é um enunciado deste documento:
\[
\boxed{\;
\text{dobra}\to\text{caminho}\to\text{rastro}\to\text{encaixe}\to\text{ponto}
\;}
\]
É isto que a dobra explica, e era a promessa da primeira página: \emph{como um
objeto finito carrega a assinatura de um ponto sem ter de o conter}. Ele não o
contém --- contém $(m,\delta)$: onde está, e quanto disso não cabe.

\medido{\code{tests/meio.c} §M8: os dois desfechos medidos lado a lado --- a AGM
a acabar em $m=g$ com o ponto na grade, e a torre de duplicação a acabar com o
ponto cercado --- e o centro contra as suas fronteiras em três grades ($6$ a $13$
vezes melhor). \code{tests/metronomo\_pi.js} para o segundo. O valor de
referência usado para medir o erro é \emph{externo} ao ficheiro e não entra em
conta nenhuma: serve para \emph{conferir}, e o processo corre sem ele.}

\section{A série do número $\pi$}\label{fis:pi}

\noindent Uma série é uma trajetória: a sucessão das suas somas parciais. Em
aritmética \emph{inteira} --- sem vírgula, com divisão inteira --- convergir tem
um significado exato e finito, e o campo $G$ lê-o.

\begin{definicao}[a soma parcial como realização]\label{fis:def:serie}
Dada uma sucessão de termos inteiros $(a_k)$, a realização \emph{é} a sucessão
das somas parciais: $\pi(0)=a_0$ e $\pi(t+1)=\pi(t)+a_{t+1}$, com $X$ o degrau
$\Zc$. (Não se introduz letra para a soma parcial: $S$ é o deslocamento da
Def.~\ref{fis:def:ops}, e aqui a soma parcial é o próprio $\pi$.)
\end{definicao}

\begin{teorema}[o campo lê o custo da série]\label{fis:thm:serie}
Suponha-se que existe $T$ com $a_k=0$ para todo $k>T$ --- que é o que a divisão
inteira faz --- e que $\pi(0),\dots,\pi(T)$ são distintos. Se $N\ge T+1$, então
\[
\abs{\{x:G(x)>1\}}=1,\qquad
\abs{\{x:G(x)=1\}}=T,
\]
isto é: o período é $1$ --- a trajetória entrou num ponto fixo --- e a cauda
conta exatamente as somas parciais distintas, que é o \textbf{custo} da série.
\end{teorema}

\begin{proof}
Para $t>T$ tem-se $\pi(t+1)=\pi(t)$, logo $\pi(t)=\pi(T)$ para todo $t\ge T$: a
trajetória é constante a partir de $T$, o que é ser periódica de período $p=1$.
Como $\pi(0),\dots,\pi(T)$ são distintos, o primeiro índice repetido é $\ell=T$,
e o Teor.~\ref{fis:thm:periodo} aplica-se (a hipótese $N\ge\ell+2p-1=T+1$ é a do
enunciado). Donde $\abs{\{G>1\}}=p=1$ e $\abs{\{G=1\}}=\ell=T$.
\end{proof}

\medskip\noindent\textbf{E a série de $\pi$ sai da cisão por paridade.} A única
hipótese é $J^{2}=-1$. Com isso as potências de $J$ ciclam com período quatro,
\[
J^{0}=1,\quad J^{1}=J,\quad J^{2}=-1,\quad J^{3}=-J,\quad J^{4}=1,\ \dots
\]
e a série da exponencial \textbf{parte-se pela paridade do índice}: os $n$
\emph{pares} caem em $\pm1$ e os \emph{ímpares} em $\pm J$. Recolhendo cada
metade:
\[
\boxed{\;
\exp(tJ)=c(t)\cdot 1+s(t)\cdot J,\qquad
c(t)=\sum_{k}\frac{(-1)^{k}t^{2k}}{(2k)!},\qquad
s(t)=\sum_{k}\frac{(-1)^{k}t^{2k+1}}{(2k+1)!}\;}
\]
--- e o índice $k$ é o que conta as \emph{voltas} do ciclo de quatro: o termo
$n=2k$ deu $k$ voltas completas mais nada, o $n=2k+1$ deu $k$ voltas mais um
quarto. \emph{O $(-1)^{k}$ é o número de voltas lido módulo dois.}

\medskip\noindent E essa cisão é a mesma do algoritmo da §\ref{fis:algoritmo}:
lá, duas máscaras complementares tomam as posições pares e as ímpares, e
\emph{particionam} o byte. Aqui as mesmas duas classes particionam a série, e por
isso as duas metades não se misturam. \textbf{A paridade é o mesmo corte nos dois
sítios.}

\begin{proposicao}[a conservação, e a definição de $\pi$]\label{fis:prop:pi}
$c^{2}+s^{2}\equiv1$, e
\[
\boxed{\;\pi=\min\{\,t>0:\ \exp(tJ)\cdot 1=-1\,\}\;}
\]
isto é, \emph{o tempo que o fluxo leva a realizar a primeira equação da
Def.~\ref{fis:def:op}}.
\end{proposicao}

\begin{proof}
Das duas séries vem $c'=-s$ e $s'=c$; logo
$(c^{2}+s^{2})'=2cc'+2ss'=-2cs+2sc=0$, e o valor em $0$ é $1$, pelo que o par
nunca sai do nível.
\end{proof}

\medskip\noindent Nenhuma palavra disto é geométrica: não há arco, não há área, e
o círculo aparece \emph{depois}, como a figura que este fluxo desenha ---
\emph{a figura não é o dado, é o rasto}. E é isto que faz de $\pi$ uma
\textbf{saída} do relógio e não uma constante importada: cada dimensão tem o seu,
e ele lê-se contando voltas.

\medskip\noindent\textbf{Onde ela cai nas duas dependências.} O termo $a_k$ é
função de $k$, não da soma corrente: a série, como trajetória, lê o
\emph{índice}. Ela dobra apenas no fim, ao entrar no ponto fixo --- e é aí, e só
aí, que passa a ser função do estado. É o Cor.~\ref{fis:cor:duais} no caso mais
simples possível: \emph{a convergência é a passagem de ler o índice para ler a
célula}.

\medido{\code{tests/aranha\_n.c} §AN11: período $\mathbf{1}$, cauda igual ao
custo, e três combinações independentes de $\arctan$ em escala inteira
($16\arctan\tfrac15-4\arctan\tfrac1{239}$, $4\arctan\tfrac12+4\arctan\tfrac13$ e
$8\arctan\tfrac12-4\arctan\tfrac17$) a coincidirem. \emph{Nenhum dígito é escrito
à mão: se fossem copiados, a igualdade não poderia falhar, e assim pode} --- basta
que qualquer uma das três somas esteja errada.}

\colunas{Lema~\ref{fis:lem:geo}; Cor.~\ref{fis:cor:medias}; Teor.~\ref{fis:thm:corte},~\ref{fis:thm:serie}; Prop.~\ref{fis:prop:pi}}
        {a AGM, o rastro $(m,\delta)$, a raiz com uma face presa, $\pi$ como tempo}
        {\code{tests/meio.c} §M8; \code{tests/agm.c} §A1--§A3; \code{tests/pgwire.c} §W217; \code{tests/aranha\_n.c} §AN11}

% ═══════════════════════════════════════════════════════════════════════════
\part{Geometria}
% ═══════════════════════════════════════════════════════════════════════════

\section{O salto da aritmética para a geometria, numa figura só}\label{fis:semicirculo}

\noindent No semicírculo de diâmetro $a+b$, seja $P$ o ponto do diâmetro que o
divide em $a$ e $b$. Então
\[
\text{o \textbf{raio} é }m=\frac{a+b}{2},
\qquad
\text{a \textbf{altura} sobre }P\text{ é }h=\sqrt{ab}=g_{\ast},
\]
e Pitágoras no raio dá $h^{2}+\delta^{2}=m^{2}$ --- que é o
Teor.~\ref{fis:thm:BI} outra vez, $ab=m^{2}-\delta^{2}$, dito com uma figura em
vez de com uma conta. A desigualdade $g\le m$ passa então a ler-se \emph{a altura
não passa o raio}, com igualdade só quando $P$ é o centro.
\[
\boxed{\;\text{A média aritmética é a \textbf{reta}; a média geométrica é o
\textbf{círculo}}\;}
\]
e as duas são a mesma conservação vista de cada lado --- é o primeiro sítio deste
documento onde uma face se lê como \emph{comprimento} e a outra como \emph{área}.

\section{Os grafos, e quem é o operador}\label{fis:hipercubo}

\noindent $\mathcal{V}(x)$ apareceu como «a vizinhança». Ela é o conjunto dos
vértices adjacentes a $x$ num \textbf{grafo}, e $\abs{\mathcal{V}(x)}$ é o
\textbf{grau}. O campo percorre três, e não são o mesmo:

\begin{center}\small
\begin{tabular}{@{}llll@{}}
\toprule
grafo & vértice & aresta & grau \\
\midrule
a grade de posto $n$ & um ponto & $x\sim x\pm e_i$ & $2n$ \\
a árvore do encaixe & uma bola & mãe\,--\,filha & $3$ (interno) \\
o hipercubo $Q_m$ & palavra de $m$ bits & um bit trocado & $m$ \\
\bottomrule
\end{tabular}
\end{center}

\begin{teorema}[o hipercubo constrói-se pela dobra]\label{fis:thm:hiper}
$Q_m$ tem $2^m$ vértices e $m\,2^{m-1}$ arestas, e $Q_m$ são duas cópias de
$Q_{m-1}$ ligadas por um emparelhamento perfeito:
$\mathrm{ar}(Q_m)=2\,\mathrm{ar}(Q_{m-1})+2^{m-1}$.
\end{teorema}

\begin{proof}
Separando os vértices pelo último bit obtêm-se duas cópias de $Q_{m-1}$; as
arestas entre elas são as que trocam esse bit, uma por vértice de cada cópia ---
um emparelhamento perfeito com $2^{m-1}$ arestas. A contagem $m\,2^{m-1}$ sai de
cada um dos $2^m$ vértices ter grau $m$.
\end{proof}

\begin{teorema}[o espaço não é escolhido: a dobra dá o hipercubo]\label{fis:thm:espaco}
Seja $F_w$ o andar de largura $w$ da escada --- o índice conta a \emph{largura},
de modo que $F_w$ tem $2^{w}$ elementos --- com $F_{2w}=F_w\oplus\sigma F_w$.
Então $(F_{2w},+)\cong B^{\,2w}$, e o grafo de $F_{2w}$ com os geradores canónicos
\textbf{é} $Q_{2w}$: duas cópias do grafo de $F_w$ ligadas pelo emparelhamento
$\sigma$.
\end{teorema}

\begin{proof}
Em característica $2$ a soma é o ou-exclusivo coordenada a coordenada, o que dá o
isomorfismo de grupos aditivos. Os geradores canónicos são os $e_k$, e trocar um
bit é somar um deles --- que são exatamente as arestas de $Q_{2w}$. Separando pelo
bit de $\sigma$ obtêm-se as duas cópias e o emparelhamento, como no
Teor.~\ref{fis:thm:hiper}.
\end{proof}

\medskip\noindent\textbf{E neste grafo o campo $G$ É o triângulo de Pascal.} Um
percurso de $n$ passos pelas direções de aresta chega à célula que a \emph{soma}
dos passos indica; e a soma é comutativa, logo esquece a ordem por que foram
dados. Quantos índices caem então na mesma célula? Os que usam os mesmos eixos por
outra ordem --- e isso é $\binom{n}{k}$:
\[
G \;=\; \Bigl(\tbinom{n}{0},\tbinom{n}{1},\dots,\tbinom{n}{n}\Bigr),
\qquad
\sum G \;=\; 2^{n} \;=\; \abs{I} .
\]
A segunda igualdade é o Lema~\ref{fis:lem:conserva} lido no triângulo: cada linha
soma o andar inteiro, porque nenhum índice se perde --- só se sobrepõe. E o
triângulo também não vem de fora: a recorrência
$\binom{n}{k}=\binom{n-1}{k-1}+\binom{n-1}{k}$ \emph{é} o passo da torre lido como
contagem --- ou o eixo novo entra, ou não entra, dois casos exaustivos e
disjuntos. A simetria $\binom{n}{k}=\binom{n}{n-k}$ é a \textbf{involução}, e o
\textbf{máximo de $G$ cai no seu ponto fixo}, $k=n/2$: \emph{a dobra é maior no
meio}, que é onde mais caminhos se sobrepõem.

\begin{teorema}[as direções de aresta são a base ortonormal]\label{fis:thm:base}
As $m$ direções de aresta de $Q_m$ são os $e_k=2^{k}$, e para $m=8$ elas são a
base de $F_8$ sobre $B$ obtida como \emph{produtos} dos três geradores das três
dobras --- ortonormal para $\langle a,b\rangle=$ paridade$(a\wedge b)$, com Gram
$=\mathrm{Id}$. Uma \textbf{régua} é um peso por direção, $g^{(k)}$ medindo só
$e_k$: cada uma é cega às outras sete, as oito compostas medem todo salto, e
tirando uma qualquer fica uma direção por medir.
\end{teorema}

\begin{proof}
$\langle e_i,e_j\rangle=$ paridade$(2^i\wedge2^j)$ vale $1$ se $i=j$ e $0$ caso
contrário, porque potências de dois distintas não partilham bits: a matriz de
Gram é a identidade, e verifica-se nos $64$ pares. Que $g^{(k)}$ vale $1$ em
$e_k$ e $0$ nas outras direções é a definição; a soma das oito é a identidade, e
retirar a $j$-ésima anula a medida de $e_j$.
\end{proof}

\medskip\noindent Como a base é ortonormal, \emph{medir} e \emph{ler o bit} são a
mesma operação --- $\langle b,e_k\rangle=(b\gg k)\wedge1$ --- e é por isso que a
forma bilinear que dá o caractere da §\ref{fis:transf} é a mesma que dá as
coordenadas. E o oito não foi escolhido: a dimensão cresce por \emph{duplicação},
$2^{1}\to2^{2}\to2^{4}\to2^{8}$, e $8$ é o que três dobras binárias dão.

\medido{\code{tests/aranha\_n.c} §AN19: $F_{2w}=F_w\oplus\sigma F_w$ como
$Q_{2w}$. §AN24: Gram $=\mathrm{Id}$ pelos caracteres; as oito réguas cegas e
compostas; e os $\mathbf{256}$ elementos a reconstruírem-se das coordenadas.
§AN17: o grau nos três grafos. \code{tests/pascal.c}: o máximo de $G$ no ponto
fixo da involução.}

\section{O dragão: a régua de bits e a dobra}\label{fis:dragao}

\noindent A curva de Heighway é a realização de posto $2$ que dobra mais
economicamente, e é onde $G>1$ se vê a olho. Ela admite duas construções sem uma
linha em comum, e é essa coincidência que a fixa.

\begin{definicao}[régua de bits]\label{fis:def:bits}
Para $k\ge1$ escreva-se $k=2^{a}m$ com $m$ ímpar. Ponha-se $\upsilon(k)=+1$ se
$m\equiv1\pmod 4$ e $-1$ se $m\equiv3\pmod 4$. A \textbf{curva de bits}
$\mathcal{C}_K$ é a sucessão $P_0=(0,0)$, $P_1=P_0+u_0$ com $u_0=(1,0)$, e
$P_{s+1}=P_s+u_s$, onde $u_{s}=R_{\upsilon(s)}\,u_{s-1}$, com $R_{+1}$ a rotação
de $+90^{\circ}$ e $R_{-1}$ a de $-90^{\circ}$.
\end{definicao}

\begin{definicao}[a dobra]\label{fis:def:dobrar}
Dada uma sucessão finita de pontos $\mathcal{D}=(P_0,\dots,P_M)$, o seu
\textbf{dual} é $\mathcal{D}^{*}=(P^{*}_0,\dots,P^{*}_M)$ com
$P^{*}_j=P_M+R\,(P_M-P_{M-j})$, e $R$ uma rotação fixa de $\pm90^{\circ}$. A
\textbf{soma} $\mathcal{D}+\mathcal{D}^{*}$ é a concatenação suprimindo o ponto
repetido. Define-se $\mathcal{D}_0=\bigl((0,0),(1,0)\bigr)$ e
$\mathcal{D}_{k+1}=\mathcal{D}_k+\mathcal{D}_k^{*}$.
\end{definicao}

\noindent Em palavras: $\mathcal{D}^{*}$ é a curva percorrida ao contrário e
rodada de $90^{\circ}$ em torno do ponto final. \emph{É a dobra do papel} ---
dobrar a tira ao meio é colar-lhe a sua própria imagem invertida --- e é o mesmo
passo $T+T^{*}$ do Teor.~\ref{fis:thm:espaco}, aqui numa curva.

\begin{lema}[antissimetria da régua]\label{fis:lem:anti}
Para $1\le j<2^{k}$ vale $\upsilon(2^{k}+j)=-\upsilon(2^{k}-j)$.
\end{lema}

\begin{proof}
Escreva-se $j=2^{a}m$ com $m$ ímpar; como $j<2^{k}$, tem-se $a<k$. Então
\[
2^{k}-j=2^{a}\bigl(2^{k-a}-m\bigr),\qquad
2^{k}+j=2^{a}\bigl(2^{k-a}+m\bigr),
\]
e como $a<k$ o fator $2^{k-a}$ é par, logo $2^{k-a}\mp m$ são \emph{ímpares} ---
são portanto as partes ímpares de $2^{k}\mp j$. Ora
$\bigl(2^{k-a}-m\bigr)+\bigl(2^{k-a}+m\bigr)=2^{\,k-a+1}\equiv0\pmod 4$, porque
$k-a+1\ge2$. Dois ímpares cuja soma é $\equiv0\pmod4$ são um $\equiv1$ e o outro
$\equiv3$ módulo $4$. Logo $\upsilon$ toma valores opostos neles.
\end{proof}

\begin{teorema}[os dois caminhos dão a mesma curva]\label{fis:thm:dragao}
Para todo $k\ge0$, a curva de bits $\mathcal{C}_k$ e a curva $\mathcal{D}_k$
coincidem ponto a ponto, para uma das duas escolhas do sinal de $R$ --- a mesma
para todo $k$.
\end{teorema}

\begin{proof}
Indução em $k$. Para $k=0$ ambas são $\bigl((0,0),(1,0)\bigr)$.

Suponha-se $\mathcal{C}_k=\mathcal{D}_k$, com $M=2^{k}$ passos, de extremo $P_M$.
A curva $\mathcal{C}_{k+1}$ tem $2^{k+1}$ passos e as suas primeiras $M$ direções
são as de $\mathcal{C}_k$ --- porque $\upsilon(s)$ não depende de $k$ --- logo
prolonga-a. Resta identificar a segunda metade.

Para $1\le j<M$, a direção do passo $M+j$ obtém-se da anterior pela rotação
$R_{\upsilon(M+j)}$, e pelo Lema~\ref{fis:lem:anti}
$\upsilon(M+j)=-\upsilon(M-j)$. Ou seja: \emph{a sucessão de viragens da segunda
metade é a da primeira metade lida de trás para a frente e com os sinais
trocados}. Ora percorrer $\mathcal{C}_k$ ao contrário troca cada viragem pela
oposta --- é a definição de percorrer ao contrário, pois uma viragem à esquerda
vista do outro sentido é uma viragem à direita. Logo a sucessão de viragens de
$\operatorname{rev}(\mathcal{C}_k)$ é exatamente a da segunda metade de
$\mathcal{C}_{k+1}$. Duas curvas com a mesma sucessão de viragens e o mesmo passo
unitário diferem por uma isometria direta do plano; como ambas começam em $P_M$,
a isometria é a rotação em torno de $P_M$, e ela é de $\pm90^{\circ}$ porque
$\upsilon$ toma valores em $\{\pm1\}$. Portanto a segunda metade é
$P_M+R\,(P_M-P_{M-j})$, que é $\mathcal{D}_k^{*}$.

Que o sinal de $R$ é o \emph{mesmo} para todo $k$: ele é determinado por
$\upsilon(2^{k}+1)$; os casos $k=0$ e $k=1$ verificam-se diretamente, e para
$k\ge2$ o número $2^{k}+1$ é ímpar com $2^{k}+1\equiv1\pmod4$, valor constante.
\end{proof}

\medskip\noindent\textbf{Duas construções, nenhuma linha em comum.} A
Def.~\ref{fis:def:bits} decide cada viragem por uma conta sobre os bits do
índice; a Def.~\ref{fis:def:dobrar} nunca olha para um índice --- copia, inverte
e roda a curva inteira. E a coincidência é sobre a \emph{curva}, e não sobre a
medida: $\abs{I}$ e $\sum_x G(x)$ coincidem para qualquer realização do mesmo
comprimento (Teor.~\ref{fis:thm:medida}). \emph{A concordância só tem conteúdo
quando é sobre o que a medida não vê.}

\medido{\code{tests/aranha\_n.c} §AN3: bits $=$ dobra
$\mathcal{D}_k+\mathcal{D}_k^{*}$, $k=1..12$, com os dois sinais de $R$ corridos e
\textbf{exatamente um} a bater. Dobra e controlos em §AN4.}

\section{A aresta é a memória da divisão}\label{fis:folha}

\noindent Uma aresta não é a ligação entre dois vértices: ela \textbf{é a memória
da divisão}. Um vértice interior da grade divide-se em $2n$ saídas, e a aresta é
o registo de \emph{por qual delas} se saiu. Daí que $\abs{\mathcal{V}(x)}=2n$ não
seja uma contagem de vizinhos, mas a capacidade local de abrir o estado em
direções distintas; e daí que os passos se codifiquem na própria base, sem
inventar codificação nenhuma.

\begin{definicao}[a palavra da trajetória]\label{fis:def:palavra}
Seja $\mathcal{A}$ o alfabeto das direções de aresta --- $2n$ letras numa grade de
posto $n$, porque $+e_i\neq-e_i$; $m$ letras no hipercubo, porque $\oplus e_i$ é
involutivo. Um percurso determina a \textbf{palavra} $s_0s_1\cdots s_{N-1}$ com
$\pi(t+1)=\pi(t)+T(e_{s_t})$, onde $T$ diz como cada letra se realiza.
\end{definicao}

\begin{teorema}[a palavra reconstrói, e a dobra é o núcleo]\label{fis:thm:palavra}
$\pi$ recupera-se de $\pi(0)$ e da palavra. Além disso, a posição é a \emph{soma}
das letras realizadas, e a soma é comutativa; logo o morfismo
$w\colon\mathcal{A}^{*}\to X$ tem \textbf{núcleo}, e
\[
G(x)>1
\iff
\text{dois prefixos da palavra têm a mesma imagem}
\iff
\text{um fator da palavra está em }\ker w .
\]
\end{teorema}

\begin{proof}
A reconstrução é a recorrência da Def.~\ref{fis:def:palavra}. Se $\pi(i)=\pi(j)$
com $i<j$, os prefixos de comprimentos $i$ e $j$ têm a mesma imagem, e o fator
entre eles tem imagem $0$; o recíproco é a mesma conta ao contrário.
\end{proof}

\medskip\noindent Um fator com imagem $0$ é um \emph{ciclo fechado}, o que liga
isto ao Teor.~\ref{fis:thm:euler}. Em três linhas, e cada uma guarda menos que a
anterior:
\[
\boxed{
\begin{array}{l}
\text{a \textbf{palavra} guarda a sequência das escolhas;}\\
\text{a \textbf{aresta} \emph{é} a memória da divisão que produziu a posição;}\\
\text{o \textbf{espaço} guarda só a posição --- e }G\text{ mede o que se colapsou.}
\end{array}}
\]
\[
\boxed{\ \text{aresta}=\text{memória da divisão}\ }
\qquad
\boxed{\ \text{dobra}=\text{a célula esquecer a ordem}\ }
\]

\medskip\noindent\textbf{E a estrutura não é a geometria.} A palavra vive no
alfabeto; a geometria entra por $T$. A imagem é literal: a \textbf{palavra é a
folha} --- plana, de ordem total, escrita num alfabeto de movimentos --- e $T$ é
o \textbf{ato de colar} essa folha no espaço. Quando a colagem faz a folha voltar
sobre si própria, uma célula recebe mais do que uma ordem, e é aí, e só aí, que
$G>1$. \emph{A dobra não é defeito da palavra: é o ato de colagem. O dragão não é
a folha nem o espaço --- é a folha já colada.}

\medskip\noindent E o que a produz não é a dimensão: é $T$ poder \emph{voltar}.
Todas as letras a avançar no mesmo sentido dá uma realização monótona, logo
injetiva, em qualquer dimensão. \emph{A dobra é propriedade da realização.}

\medido{\code{tests/aranha\_n.c} §AN30: as $\mathbf{1024}$ letras do dragão; a
reconstrução com divergência $\mathbf{0}$; a palavra permutada a chegar ao mesmo
fim por outro caminho; a paridade no hipercubo; e a \emph{mesma} palavra com três
$T$ diferentes a dar $\max G$ de $\mathbf{2}$ (posto $2$, com $\pm e_i$),
$\mathbf{1}$ (posto $4$, com quatro direções independentes, porque somas de letras
distintas nunca colidem) e $\mathbf{28}$ (na reta, porque há menos para onde ir).}

\colunas{Teor.~\ref{fis:thm:hiper},~\ref{fis:thm:espaco},~\ref{fis:thm:base},~\ref{fis:thm:dragao},~\ref{fis:thm:palavra}; Lema~\ref{fis:lem:anti}}
        {o semicírculo, $Q_m$, Pascal, as oito réguas, o dragão, a folha colada}
        {\code{tests/aranha\_n.c} §AN3--§AN4, §AN17, §AN19, §AN24, §AN30; \code{tests/pascal.c}}

% ═══════════════════════════════════════════════════════════════════════════
\part{Mecânica}
% ═══════════════════════════════════════════════════════════════════════════

\section{O ciclo da aranha}\label{fis:algoritmo}

\noindent Agora que a construção está feita --- objeto, multiplicidade,
vizinhança --- o autómato \emph{cai} dela: o Teor.~\ref{fis:thm:mult} já
\textbf{é} um autómato, com as quatro cláusulas a serem os quatro passos e a
ordem delas a ser a ordem em que ele corre. Não há aqui um algoritmo genérico
proposto de fora e depois justificado; há a construção lida como dinâmica.

\begin{enumerate}\itemsep3pt
\item[1.] \textbf{Escrita} --- $G(x)\mathrel{+}=1$ na célula corrente
$x=\pi(t)$. Realiza a cláusula~(3).
\item[2.] \textbf{Leitura} --- $\mathcal{P}_{t}(x)=\bigl(G_t(y)\bigr)_{y\in
\mathcal{V}(x)}$ sobre a vizinhança; a fronteira lê-se como obstáculo. Realiza a
cláusula~(4).
\item[3.] \textbf{Decisão} --- gradiente de exploração: o passo vai ao vizinho de
\emph{menor} $G$. Lê a dobra; não pergunta «por onde passei?».
\item[4.] \textbf{Fecho} --- ao fim dos passos de exploração, o rotor
$R^{4}=\mathrm{id}$ (Teor.~\ref{fis:thm:fecho}).
\end{enumerate}

\medskip\noindent\textbf{Não há deliberação em passo nenhum.} A decisão do
ponto~3 não é uma escolha entre planos: é o mínimo de $\mathcal{P}_t(x)$, e
$\mathcal{P}_t(x)$ é o que está escrito no chão. O agente não compara percursos
nem avalia futuros --- \emph{reage à marca}. É por isso que a cláusula~(3) tem o
peso que tem: \emph{tirada a marca do espaço, não sobra estado de onde decidir}.
Um agente que guardasse $\pi(0),\dots,\pi(t)$ teria a informação toda e o campo
seria supérfluo; a aranha é o caso em que ele não é.

\medskip\noindent\textbf{Estado global, e memória do agente.} O \emph{sistema}
mantém dois dados: o campo $G$, que está no espaço, e a célula corrente. O
\emph{agente} não mantém nenhum dos dois como história. Dizer «não há estado»
seria falso --- o que não há é estado \emph{no agente}.

\medskip\noindent\textbf{Representação de $G$.} Armazená-lo como \emph{tabela
sobre $X$} é uma escolha de representação, não o algoritmo: custa $\abs{X}$ e não
$\abs{I}$. A representação fiel à cláusula~(3) é \textbf{esparsa}: \emph{a memória
existe onde a trajetória passou}. Com endereçamento aberto num arranjo fixo de
capacidade $C\ge2\abs{I}$, uma escrita ou leitura custa $O(1)$ esperado e a
memória total é $O(\abs{I})$, decidida antes de correr --- não há alocação
dinâmica em parte nenhuma.

\medskip\noindent\textbf{Complexidade, e onde o $n$ aparece nela.} São $N+1$
passos. O passo~1 escreve num ponto: $O(1)$. Os passos~2 e~3 leem
$\abs{\mathcal{V}(x)}$ valores e tomam o mínimo: $\Theta(n)$ no posto $n$. Total:
\[
\Theta\bigl(n\,\abs{I}\bigr)\ \text{operações},\qquad
\Theta\bigl(\abs{I}\bigr)\ \text{memória},
\]
e \textbf{nenhuma dependência de $\abs{X}$} --- que é a afirmação que importa,
porque é a que separa o algoritmo da tabela. \emph{Não é varredura do espaço: é
leitura de $2n$ células vizinhas, o mesmo em qualquer ponto e em qualquer
instante.}

\begin{teorema}[$G$ não vê o construtor]\label{fis:thm:agentes}
Atribua-se a cada escrita um \emph{agente}, e considerem-se duas execuções que
produzam a mesma multiplicidade por célula --- por exemplo, uma em que \emph{um}
agente visita cada célula $K$ vezes, e outra em que $K$ agentes a visitam uma vez
cada. Então os campos coincidem célula a célula, e $\abs{I}$ é o mesmo. Isto é:
$G$ é \textbf{invariante sob qualquer reatribuição de agentes} que preserve a
ocupação.
\end{teorema}

\begin{proof}
Pela cláusula~(3), cada escrita incrementa $G$ na célula corrente e nada mais
regista; em particular nenhuma escrita carrega a identidade de quem a fez. Como
as duas execuções produzem a mesma multiplicidade de escritas por célula, os
campos coincidem.
\end{proof}

\medskip\noindent Convém dizer o enunciado ao contrário para não o ler ao
contrário: não é que $G$ \emph{identifique} coisa nenhuma --- é que $G$ não
\emph{muda} quando se trocam os construtores. Ele regista \textbf{ocupação}.
Nenhuma etapa da cadeia
\[
\text{base}\longrightarrow\text{palavra de arestas}\longrightarrow T
\longrightarrow\text{posição}\longrightarrow G
\]
carrega quem construiu --- $T$ é uma soma comutativa, e ela descarta o agente
\emph{antes} de a posição existir --- e por isso o campo não o pode devolver. Em
uma linha:
\[
\boxed{\ \text{o campo é memória da \emph{construção}, e não do \emph{construtor}}\ }
\]
Isto não é uma limitação, é a cláusula~(3) a fazer o seu trabalho, e tem duas
consequências. A primeira é que o mecanismo \textbf{escala}: acrescentar
construtores não muda nada no ciclo, e uma ocupação por \emph{multiplicidade de
agentes} e uma expansão pela \emph{memória de um só} são indistinguíveis do lado
do espaço. A segunda é sobre o que cada coordenada repõe: a folha do
levantamento repõe a \emph{ordem} da visita, dentro de um agente; separar
\emph{agentes} pede outra coordenada --- o identificador --- e essa não sai do
campo.

\medido{\code{tests/aranha\_n.c} §AN31: um agente com $4$ voltas e $4$ agentes
com uma volta --- campos iguais nas $\mathbf{64}$ células, $G=4$ nas $16$, com o
controlo negativo a ver as $16$ a diferir contra $3$ voltas. E a separação: a
ordem separa \textbf{todos} os pares que caem na mesma célula ($\mathbf{96}$ de
$96$), e o identificador separa só os que atravessam a fronteira entre agentes
($\mathbf{64}$).}

\section{A lei local: vértice e aresta}\label{fis:handshake}

\begin{teorema}[a lei local]\label{fis:thm:handshake}
Seja $\pi$ um percurso e $G_{\mathrm{ar}}(e)$ o número de vezes que ele usa a
aresta $e$. Então, para cada vértice $x$,
\[
\sum_{e\,\ni\,x}G_{\mathrm{ar}}(e)\;=\;2\,G(x)\;-\;[\,x=\pi(0)\,]\;-\;[\,x=\pi(N)\,],
\qquad\text{e}\qquad \sum_{e}G_{\mathrm{ar}}(e)=N .
\]
\end{teorema}

\begin{proof}
Cada visita a $x$ num índice interior $0<t<N$ usa duas arestas incidentes, a de
$t-1$ para $t$ e a de $t$ para $t+1$; a visita em $t=0$ usa uma só, e a em $t=N$
também. Somando sobre as $G(x)$ visitas vem o enunciado. A segunda igualdade é
contar as $N$ arestas do percurso.
\end{proof}

\medskip\noindent A multiplicidade do vértice determina, assim, a das arestas
incidentes: \emph{$G$ não é uma contagem entre outras --- é a que fecha a
contabilidade local do percurso}. É o \emph{handshaking} dito sobre a trajetória,
e os dois termos de correcção são as pontas: o que entra iguala o que sai excepto
onde a trajetória começa e acaba.

\medido{\code{tests/aranha\_n.c} §AN18: a lei local com $\mathbf{0}$ violações em
$513$ vértices; e quais realizações são percursos e quais não.}

\section{O levantamento, e a volta}\label{fis:levantamento}

\noindent A projeção $\pi$ perde exatamente a dobra. A pergunta dual: pode a dobra
ser \emph{desfeita} sem inventar geometria? Sim, e ao preço de uma coordenada.

\begin{definicao}[número de visita]\label{fis:def:k}
Para $i\in I$, $k(i)=\bigl\lvert\{j\in I: j\le i,\ \pi(j)=\pi(i)\}\bigr\rvert$.
\end{definicao}

\begin{teorema}[levantamento em folhas]\label{fis:thm:levant}
Seja $\tilde\pi(i)=\bigl(\pi(i),k(i)\bigr)$ e $\tilde G$ a multiplicidade de
$\tilde\pi$. Então:
\begin{enumerate}\itemsep3pt
\item[(1)] \textbf{Injetividade:} $\tilde\pi$ é injetiva, e portanto
$\tilde G\equiv1$ na imagem.
\item[(2)] \textbf{Projeção:} $\operatorname{pr}_1\circ\tilde\pi=\pi$.
\item[(3)] \textbf{Fibra vertical:} para cada $x\in\operatorname{im}\pi$,
$\{k(i):\pi(i)=x\}=\{1,2,\dots,G(x)\}$, exatamente.
\item[(4)] \textbf{Conservação:}
$\abs{\operatorname{im}\tilde\pi}=\abs{I}=\sum_x G(x)$.
\item[(5)] \textbf{Volta:} $G$ recupera-se de $\tilde\pi$ sem qualquer outro
dado.
\end{enumerate}
Além disso, numa grade de posto $n$ o par $(x,k)$ ocupa $n+1$ direções: \emph{o
levantamento sobe \textbf{um} andar, seja qual for $n$}.
\end{teorema}

\begin{proof}
(3) primeiro, porque as outras se apoiam nela. Fixe-se $x$ e sejam
$i_1<i_2<\dots<i_g$ os índices da fibra $\pi^{-1}(x)$, $g=G(x)$. Por definição
$k(i_r)=\lvert\{j\le i_r:\pi(j)=x\}\rvert=r$. Logo $k$ restrita à fibra é a
enumeração $i_r\mapsto r$, que é uma bijeção sobre $\{1,\dots,g\}$.

(1) Sejam $\tilde\pi(i)=\tilde\pi(j)$. Então $\pi(i)=\pi(j)=:x$, logo $i$ e $j$
estão na mesma fibra, e $k(i)=k(j)$. Por (3) $k$ é injetiva nessa fibra, donde
$i=j$. Que $\tilde G\equiv1$ é a cláusula~(2) do Teor.~\ref{fis:thm:mult}
aplicada a $\tilde\pi$.

(2) É a definição da projeção. (4) Por (1), $\tilde\pi$ é uma bijeção de $I$ sobre
a sua imagem; a segunda igualdade é o Lema~\ref{fis:lem:conserva}. (5) Por (2), a
pré-imagem da primeira coordenada é $\pi^{-1}(x)$, cuja cardinalidade é $G(x)$.

A última afirmação é que a folha acrescenta uma direção às $n$ da grade.
\end{proof}

\medskip\noindent\textbf{O algoritmo inverso.} Inicializar um campo $c\equiv0$;
para $i=0,\dots,N$, ler $x\leftarrow\pi(i)$, escrever $c(x)\leftarrow c(x)+1$ e
emitir $\bigl(x,\,c(x)\bigr)$. O campo $c$ é o próprio $G$ parcial: \emph{o
inverso é o directo com o valor lido a sair}. Não inventa geometria --- desfaz a
identificação que a realização tinha feito.

\medskip\noindent\textbf{O par é uma retração, não uma involução.}
$\operatorname{pr}_1\circ\tilde\pi=\pi$ fecha de um lado só: $\tilde\pi$ sobe,
$\operatorname{pr}_1$ desce, e a composição na outra ordem não é a identidade nem
está definida fora da imagem. É a forma que separa \emph{reconstrução} de
\emph{extração}, e distingue-se de uma involução, onde os dois lados fecham no
mesmo espaço.

\medskip\noindent E a injetividade verifica-se \emph{sem comparar pares}: basta
correr o próprio ciclo \emph{sobre $\tilde\pi$}, em dimensão $n+1$, e ler
$\max\tilde G$. Isto é uma verificação por outro caminho, não a mesma conta duas
vezes --- a prova de (1) usa (3) e a ordem dos índices; a corrida usa só a
igualdade de vetores. E a volta (5) é o que distingue um \emph{inverso} de um
\emph{levantamento}: um levantamento que não devolvesse $G$ teria acrescentado
informação em vez de reorganizar a que havia.

\medido{\code{tests/aranha\_n.c} §AN5 e §AN9: $\tilde G\equiv1$ pela própria
aranha em dimensão $n+1$; $\operatorname{pr}_1\circ\tilde\pi=\pi$ com resíduo
$\mathbf{0}$; fibra vertical $\{1..G\}$; e a volta a devolver $G$ célula a
célula. \code{tests/aranha\_inversa.c} §AI1--§AI7.}

\section{O espaço e o tempo}\label{fis:tempo}

\noindent O \textbf{espaço} está definido na §\ref{fis:hipercubo}: é o que a dobra
constrói. Falta o \textbf{tempo}, e ele não se importa de fora --- deriva-se.
Opera-se um sistema \textbf{variando-o}: se nada varia, nada se vê. A variação
mais barata é um parâmetro numa \emph{reta}, e o operador dessa variação é o
deslocamento $S$ da Def.~\ref{fis:def:ops}.

\begin{definicao}[o tempo, e as duas grandezas que o acompanham]\label{fis:def:tempo}
Três coisas distintas, cada uma com o seu nome:
\begin{enumerate}\itemsep3pt
\item[(i)] O \textbf{tempo} é o parâmetro: a \emph{contagem das aplicações do
operador} $S$ ao longo da tupla. Contar essas aplicações é precisamente contar a
lista --- de modo que dizer «o tempo» e dizer «a contagem dos símbolos
percorridos» é dizer a mesma coisa.
\item[(ii)] O \textbf{tempo próprio} $\tau$ é o mesmo percurso \emph{medido por
uma régua}: $(\partial\tau)^{2}=g_{uv}\,\partial x^{u}\,\partial x^{v}$.
\item[(iii)] A \textbf{dobra acumulada} $D=\zeta\theta$, com
$\theta(t)=[\,G_{t}(\pi(t))>1\,]$, e $\mu D=\theta$. Esta \emph{não} é um tempo:
é memória. No fim do percurso $D(N)=\abs{I}-v$ --- o mesmo $D$ do
Teor.~\ref{fis:thm:euler}, e não outra grandeza com a mesma letra.
\end{enumerate}
(O \emph{incremento} escreve-se $\partial$, que é o operador da
Def.~\ref{fis:def:op} --- não há aqui um diferencial importado do contínuo.)
\end{definicao}

\medskip\noindent A distinção não é de estilo: chamar «tempo» a (iii) confunde o
que \emph{conta os passos} com o que \emph{mede o que os passos deixaram}. E é
(i) que responde à pergunta de onde vem o tempo: \emph{dos passos do operador, e
de mais nada}.

\medskip\noindent\textbf{E o número depende da régua.} O comprimento de um salto
não é dado pelo percurso --- é dado pela métrica com que se mede. O mesmo
percurso, medido na grade, tem todos os saltos iguais a $1$; medido na
ultramétrica da árvore sobre o parâmetro, tem saltos de comprimento
$2^{-\operatorname{prof}(t,t+1)}$, que varia com a profundidade em que os índices
divergem. As duas são réguas do mesmo objeto: dão a mesma ordem de índices e
nenhuma colapsa um salto a zero. \emph{O isomorfismo é entre as réguas, e não
entre os números que elas produzem.}

\subsection*{A métrica}

\noindent O salto tem \textbf{duas} grandezas, e elas não são a mesma. A primeira
é a \textbf{forma quadrática} no incremento; a segunda é o \textbf{tempo próprio},
que é a soma da sua raiz:
\[
\boxed{\;q(\partial x):=g_{uv}\,\partial x^{u}\,\partial x^{v},
\qquad\qquad
\tau=\sum_{t}\sqrt{q(\partial x_t)}\;}
\]
com $g$ simétrica e, no seu domínio, não degenerada. \emph{Sobre uma aresta $q$ é
inteiro}; só $\tau$ pede raiz, e é por isso que a construção corre em inteiros
enquanto se ficar por $q$. Confundir as duas é o erro fácil: com $g_{ii}=4$ um
salto tem $q=4$ e $\partial\tau=2$.

\begin{teorema}[a régua canónica]\label{fis:thm:gcanon}
Uma aresta na direção $i$ dá $q=g_{ii}$: \emph{o percurso só sonda a diagonal de
$g$}. Logo
\[
\tau=\sum_{t}\sqrt{g_{i(t)i(t)}} ,
\]
e $\tau=N$ --- o tempo próprio coincide com o parâmetro --- se e só se
$g_{ii}=1$ para toda direção usada; e isso \emph{não} exige $g=\mathrm{id}$: os
termos fora da diagonal permanecem invisíveis a um percurso de arestas.
\end{teorema}

\begin{proof}
Uma aresta é $\partial x=\pm e_i$, com uma componente só, logo
$\partial x^{u}\partial x^{v}=0$ fora da diagonal e $q=g_{ii}$. Somando as raízes
sobre os $N$ saltos vem $\tau=\sum_t\sqrt{g_{i(t)i(t)}}$. Cada parcela é $\ge0$ e
vale $1$ exatamente quando $g_{ii}=1$; logo a soma vale $N$ exatamente quando
todas valem, que é a condição do enunciado. Para medir os termos cruzados seria
preciso um salto diagonal --- que, pela Def.~\ref{fis:def:percurso}, não é aresta.
\end{proof}

\medskip\noindent É esta a razão de a grade ser a régua barata: nela, variar o
parâmetro e andar no espaço são a mesma conta.

\medskip\noindent\textbf{E $g$ não tem de ser constante.} Generalizando-a para
depender do ponto, e deixando que o que a faça variar seja o próprio campo,
\[
\boxed{\;g_{uv}(x)=G(x)\,h_{uv}
\qquad\Longrightarrow\qquad
q=G(x)\,h_{ii}\ \text{ numa aresta},\qquad
\partial\tau=\sqrt{G(x)\,h_{ii}}\;}
\]
isto é: \textbf{o quadrado do intervalo de um salto é a multiplicidade da célula
de chegada}, e o tempo próprio é a sua raiz. Onde a trajetória já passou, o
tempo \emph{próprio} passa outro --- e é a memória, ao ficar no espaço, que o
curva. O parâmetro, esse, continua a contar os passos do sucessor, um por um. Com
$G\equiv1$ a métrica é plana e $\tau=N$; onde há dobra, $\tau$ afasta-se do
parâmetro.

\medskip\noindent\textbf{E o domínio tem de se dizer.} Como
$\det g(x)=G(x)^{n}\det h$, a forma degenera exatamente onde $G(x)=0$. \emph{A
métrica é considerada sobre o suporte} $\operatorname{supp}G=\operatorname{im}\pi$,
onde $G\ge1$ pela Def.~\ref{fis:def:objeto} --- e aí $g=G\,h$ permanece não
degenerada. Fora do suporte não há trajetória para medir, e por isso não há o que
pedir à régua: \emph{a métrica existe onde a trajetória passou}, que é a mesma
frase da cláusula~(3) do Teor.~\ref{fis:thm:mult} dita sobre $g$ em vez de sobre
$G$.

\begin{teorema}[o passo é o mínimo métrico local]\label{fis:thm:geodesica}
Se $h=\mathrm{id}$, o vizinho de menor $G$ é o vizinho de menor $q$ --- e, porque
a raiz é \emph{monótona}, também o de menor incremento $\partial\tau$. Portanto o
passo~3 do ciclo move-se, a cada salto, por onde a métrica é mínima. (É
minimização \emph{local}, salto a salto.) A hipótese é necessária: para $h$
anisotrópica as duas escolhas divergem.
\end{teorema}

\begin{proof}
Com $h=\mathrm{id}$, $q=G(y)$ para todo vizinho $y$, e minimizar um é minimizar o
outro; e como $t\mapsto\sqrt t$ é estritamente crescente em $t\ge0$, o mínimo de
$q$ é também o mínimo de $\sqrt q=\partial\tau$ --- \emph{as três escolhas
coincidem}. Com $h_{ii}$ não constante, $q=G(y)h_{ii}$ e um vizinho de $G$ maior
pode ter $q$ menor; exibe-se um.
\end{proof}

\medskip\noindent\emph{E convém dizer qual das duas grandezas o passo minimiza:}
ele lê $G$, logo minimiza $q$. Que isso seja também o mínimo de $\partial\tau$ é
consequência da monotonia da raiz, e não uma segunda hipótese --- mas é uma
consequência, e não a definição.

\medskip\noindent Fica assim dito o que o passo~3 é: não uma heurística de
exploração nem uma escolha, mas \emph{o mínimo que a métrica impõe em cada
salto}.
\[
\boxed{\;\text{A aranha não decide --- \textbf{cai}.}\;}
\]

\medido{\code{tests/aranha\_n.c} §AN21--§AN22: $(\partial\tau)^2$ inteiro por
aresta com três $g$; $\tau=N$ com
$g=\bigl(\begin{smallmatrix}1&3\\3&1\end{smallmatrix}\bigr)$, que tem a diagonal
a $1$ e o cruzado a $3$ --- \emph{a testemunha de que $g=\mathrm{id}$ não é
preciso} --- e uma $g$ diagonal com uma entrada a $4$ já a afastar, que é o
controlo sem o qual a primeira não diria nada; $g_{uv}(x)=G(x)\delta_{uv}$ curva
onde há dobra e é plana com $G\equiv1$; e a geodésica a divergir do gradiente
quando $h$ é anisotrópica. §AN20: $\mu D=\theta$ com resíduo $\mathbf{0}$; e a
mesma trajetória com salto uniforme $1$ na grade e $\mathbf{11}$ comprimentos
distintos na árvore, de $2$ a $2048$.}

\section{Órbitas deterministas: $G$ conta o período}\label{fis:periodo}

\begin{definicao}[realização determinista]\label{fis:def:determinista}
$\pi$ diz-se \textbf{determinista} se existe $f\colon X\to X$ com
$\pi(t+1)=f(\pi(t))$ para todo $t<N$.
\end{definicao}

\begin{teorema}[o período lê-se no campo]\label{fis:thm:periodo}
Seja $\pi$ determinista, e suponha-se que existem $i<j$ com $\pi(i)=\pi(j)$.
Definam-se
\[
\boxed{\;
\ell:=\min\{\,i:\exists\,j>i,\ \pi(i)=\pi(j)\,\},
\qquad
p:=\min\{\,j-\ell:\ j>\ell,\ \pi(j)=\pi(\ell)\,\}\;}
\]
--- a \textbf{cauda} e o \textbf{período}. \emph{As duas definições são
explícitas de propósito}: «primeiro índice repetido» e «início da parte
periódica» são a mesma coisa aqui, mas só depois de se provar a estrutura de
órbita, e enquanto isso não está provado convém que os símbolos não dependam
dessa identificação. Se $N\ge\ell+2p-1$, então
\[
\bigl\lvert\{x: G(x)>1\}\bigr\rvert=p
\qquad\text{e}\qquad
\bigl\lvert\{x: G(x)=1\}\bigr\rvert=\ell .
\]
Se, pelo contrário, $\pi$ é injetiva, então $G=1$ na imagem e $0$ fora dela, e os
dois conjuntos são $\varnothing$ e $\operatorname{im}\pi$.
\end{teorema}

\begin{proof}
Comece-se pela propriedade que o determinismo dá: se $\pi(i)=\pi(j)$ então
$\pi(i+m)=\pi(j+m)$ para todo $m$ com $j+m\le N$, por indução aplicando $f$. Logo
a sucessão é periódica de período $p$ a partir de $\ell$.

\emph{As células do ciclo são $p$ e são distintas.} Se fosse
$\pi(\ell+r)=\pi(\ell+s)$ com $0\le r<s<p$, aplicando $f$ $(p-s)$ vezes viria
$\pi(\ell+r+p-s)=\pi(\ell)$ com $0<r+p-s<p$, contra a minimalidade de $p$.

\emph{As células da cauda são $\ell$, distintas entre si e disjuntas do ciclo.} Se
$\pi(a)=\pi(b)$ com $a<b<\ell$, então $a$ seria um índice repetido menor que
$\ell$, contra a minimalidade de $\ell$. Se $\pi(a)=\pi(\ell+r)$ com $a<\ell$ e
$0\le r<p$, aplicando $f$ $(p-r)$ vezes viria $\pi(a+p-r)=\pi(\ell)$ --- e como
$a<\ell$, o índice $a$ repete-se, contra a minimalidade de $\ell$ novamente.

\emph{Contagem.} Como $N\ge\ell+2p-1$, a órbita completa pelo menos duas voltas:
cada célula do ciclo é atingida em dois índices distintos, donde $G\ge2$; cada
célula da cauda é atingida uma só vez, donde $G=1$. E a imagem é a união das
duas. O caso injetivo é a cláusula~(2) do Teor.~\ref{fis:thm:mult}.
\end{proof}

\medskip\noindent\textbf{A hipótese $N\ge\ell+2p-1$ é necessária.} Sem ela o
enunciado é falso, e falso do modo mais banal: se a trajetória parar depois de
entrar no ciclo mas antes de o fechar duas vezes, as células do ciclo ainda
visitadas uma só vez contam para $\{G=1\}$. Exemplo mínimo, um passo abaixo do
limiar: $f(z)=z^2-1$ com $\pi(0)=0$ e $N=2$ dá a órbita $0,-1,0$, com $\ell=0$ e
$p=2$, logo o limiar é $N\ge3$; e $\abs{\{G>1\}}=1\neq2=p$. Com $N=3$ a órbita é
$0,-1,0,-1$ e o enunciado vale. \emph{E o contra-exemplo não pode ser $N=1$}: aí
não há repetição nenhuma, a hipótese do próprio teorema falha, e $\ell$ e $p$ nem
estão definidos --- um caso onde a hipótese principal não se verifica não
testemunha nada sobre a auxiliar.

\medskip\noindent\textbf{A deteção de ciclo sem o segundo ponteiro.} O algoritmo
clássico mantém dois cursores a velocidades diferentes precisamente porque o
agente não tem memória do percurso. O Teor.~\ref{fis:thm:periodo} diz que, se a
memória estiver \emph{no espaço}, um cursor basta: $\ell$ e $p$ lêem-se no campo,
sem segunda passagem e sem comparar posições. \emph{É a cláusula~(3) a
pagar-se.}

\medskip\noindent\textbf{E o exemplo acima tem nome.} A $f(z)=z^{2}-1$ é uma
órbita de Julia, e o par que ela forma lê-se inteiro na dobra deste documento.
\textbf{Julia} é o lado \emph{multiplicativo}, $z\mapsto z^{2}$; \textbf{Cantor} é
o \emph{aditivo}, o deslocamento $\theta\mapsto2\theta$; e o que os conjuga é a
\textbf{forma polar} $z=e^{2\pi i\theta}$, que precisa apenas de $J^{2}=-1$: as
potências de $J$ ciclam, a exponencial parte-se pela paridade do índice no par
$(\cos,\sin)$, e $\cos^{2}+\sin^{2}\equiv1$ é a conservação nessa base. A
conjugação é então a identidade $\bigl(e^{2\pi i\theta}\bigr)^{2}=e^{2\pi
i(2\theta)}$ e nada mais. \emph{Elevar ao quadrado no multiplicativo é dobrar o
ângulo no aditivo}, e é só isso que a conjugação diz.

\medskip\noindent\textbf{Daí que o período seja legível, e de dois lados.} Do
lado do ângulo, $\theta=p/q$ com $q$ ímpar tem órbita de comprimento
$\operatorname{ord}_q(2)$, porque duplicar é multiplicar por $2$ no denominador
que não se gasta. Do lado do campo, o Teor.~\ref{fis:thm:periodo} dá o mesmo
número sem seguir a trajetória. \emph{São dois caminhos para a mesma contagem}, e
o segundo não precisa do ângulo.

\medskip\noindent\textbf{E o que decide se há período é a conservação, não o
operador.} Com $q$ ímpar o $2$ é invertível módulo $q$: o deslocamento não tem o
que gastar, e a palavra \textbf{cicla}. Com $q$ par ele \emph{divide} o
denominador a cada passo --- $1/12\to1/6\to1/3$ --- e só passa a ciclar quando
chega ao ímpar. Donde a formulação que fica: \emph{ciclar não é uma propriedade
do operador, é o que sobra quando o orçamento deixa de poder ser gasto}. É o
mesmo que aqui se diz da cauda e do período --- $\ell$ é o que se gastou, $p$ é o
que sobrou.

\medido{\code{tests/aranha\_n.c} §AN8: $\abs{\{G>1\}}=p$ e $\abs{\{G=1\}}=\ell$
sobre $z\mapsto z^2+c$, com o período obtido por \emph{busca independente} --- e
não pelo campo --- e a órbita que escapa como controlo.}

\section{O que o passo lê}\label{fis:trilema}

\begin{teorema}[dobrar sem período obriga a ler o índice]\label{fis:thm:leitura}
Seja $\pi$ com $G(x)>1$ para algum $x$ (dobra) e tal que $\pi$ \emph{não} é
eventualmente periódica. Então não existe $f\colon X\to X$ com
$\pi(t+1)=f(\pi(t))$: o passo depende necessariamente do índice.
\end{teorema}

\begin{proof}
Suponha-se que $f$ existe. Como $\pi$ dobra, há $i<j$ com $\pi(i)=\pi(j)$.
Aplicando $f$ repetidamente, $\pi(i+m)=\pi(j+m)$ para todo $m$ admissível, logo
$\pi(t+(j-i))=\pi(t)$ para todo $t\ge i$: $\pi$ é eventualmente periódica de
período $j-i$ a partir de $i$, contra a hipótese.
\end{proof}

\begin{corolario}[a aranha e o dragão são duais]\label{fis:cor:duais}
Sejam as duas maneiras de uma trajetória guardar a sua própria história:
\begin{center}\small
\begin{tabular}{@{}lll@{}}
\toprule
 & \textbf{a aranha} & \textbf{o dragão} \\
\midrule
a marca fica & no \textbf{espaço}: é $G$ & no \textbf{tempo}: é o índice \\
o passo lê & a célula: $\pi(t{+}1)=f(\pi(t))$ & o parâmetro: $\pi(t{+}1)=g(t,\pi(t))$ \\
o campo $G$ lê & período e cauda & a dobra, e só ela \\
\bottomrule
\end{tabular}
\end{center}
O levantamento $\tilde\pi=(\pi,k)$ é a ponte entre os dois --- e cobra o preço:
junta espaço e tempo num só estado, e a dobra desaparece, $\tilde G\equiv1$. As
três coisas --- \emph{dobrar}, \emph{ler só o estado} e \emph{não ter período} ---
não podem valer as três.
\end{corolario}

\begin{proof}
A incompatibilidade das três é o Teor.~\ref{fis:thm:leitura} lido ao contrário. A
coluna do campo é o Teor.~\ref{fis:thm:periodo} de um lado e a
Def.~\ref{fis:def:objeto} do outro; e que o levantamento apaga a dobra é o
Teor.~\ref{fis:thm:levant}(1).
\end{proof}

\medskip\noindent\textbf{E verifica-se.} Que o passo lê o parâmetro não se
argumenta: procura-se na trajetória a mesma célula com sucessores
\emph{diferentes}. Na curva de Heighway ela existe e exibe-se com os quatro
índices; no Julia, no relógio $J$ e na reta não existe. E o relógio
$J\colon(a,b)\mapsto(-b,a)$, com $J^2=-\mathrm{id}$ e $J^4=\mathrm{id}$, é função
da célula: o campo lê-lhe o período sem aparato nenhum --- quatro células com
$G>1$, cauda zero.

\medido{\code{tests/aranha\_n.c} §AN10: a mesma célula do dragão com sucessores
diferentes, exibida; Julia, relógio e reta a lerem o estado; e o período $4$ do
relógio no campo.}

\section{Realizações}\label{fis:realizacoes}

\noindent O ciclo corre em todas elas sem uma linha diferente.

\begin{center}\footnotesize
\begin{tabular}{@{}llrrrccp{0.20\textwidth}@{}}
\toprule
$n$ & realização & $\abs{I}$ & $\max G$ & células & dobra & percurso & \\
\midrule
$1$ & Cantor & $512$ & $1$ & $512$ & não & não & o endereço ternário do terço médio \\
$2$ & Heighway & $4097$ & $2$ & $2558$ & sim & sim & a curva da §\ref{fis:dragao} \\
$3$ & dragão no espaço & $2049$ & $3$ & $1750$ & sim & sim & a mesma dobra com o plano da viragem a rodar \\
$2$ & espiral quadrada & $4097$ & $1025$ & $4$ & sim & sim & sempre à esquerda \\
$2$ & reta axial & $4097$ & $1$ & $4097$ & não & sim & $x\mapsto x+e_0$ \\
$2$ & reta diagonal & $4097$ & $1$ & $4097$ & não & não & passo $(1,1)$: não é aresta \\
$2$ & relógio $J$ & $61$ & $16$ & $4$ & sim & não & $(a,b)\mapsto(-b,a)$, período $4$ \\
$2$ & Julia, $c=-1$ & $61$ & $31$ & $2$ & sim & sim & $z\mapsto z^{2}+c$ \\
$4$ & o passo da máquina & $512$ & $12$ & $48$ & sim & não & (registo, slot, banco, andar) \\
\bottomrule
\end{tabular}
\end{center}

\medskip\noindent Os números dependem da realização \emph{e} do comprimento
$\abs{I}$, e por isso ambos aparecem na tabela: $\max G$ não é propriedade
puramente geométrica da imagem da curva --- depende de quantas vezes a
parametrização visita cada célula. Duas realizações ficam de fora por serem de
outro tipo: as \emph{somas parciais} da §\ref{fis:pi}, com período $1$ e a cauda
a ser o custo; e a \emph{árvore do encaixe} da §\ref{fis:ultra}, onde a célula é
a bola e $\abs{\mathcal{V}}=1$.

\medskip\noindent\textbf{Dobrar e ser percurso são independentes}, e as quatro
combinações estão na tabela: a espiral dobra e é percurso; o relógio dobra e não
é; a reta axial não dobra e é; o Cantor e a reta diagonal não fazem nem uma coisa
nem outra. Note-se ainda que o relógio tem $\max G=16$ e período $4$: são números
diferentes --- \emph{o máximo conta visitas, o período conta células}.

\medskip\noindent\textbf{Controlos.} Uma afirmação «$G>1$ aqui» só tem conteúdo
com uma realização onde $G\equiv1$ ao lado; senão passaria com um campo que
contasse mal. Servem a reta axial e o Cantor. E note-se que são \emph{dois}: o
Cantor não dobra, apesar de ser a mais intrincada das figuras da lista.

\medskip\noindent\textbf{E $\max G$ é uma medida da curva.} Medido, o dragão de
posto $3$ tem $\max G=3$ e o do plano $2$: dobra \emph{mais}, e não menos. A
grandeza compara as duas curvas, que são construções distintas --- a de posto $3$
roda o plano da viragem --- e não os dois espaços. A dimensão, essa, já apareceu
onde tinha de aparecer: no Teor.~\ref{fis:thm:dim}.

\medido{\code{tests/aranha\_n.c} §AN23 --- a tabela inteira é \textbf{medida} e
não escrita: os nove $\max G$, as células e a coluna \emph{percurso}.}

\section{O algoritmo, e as quatro operações}\label{fis:neuronio}

\noindent O documento fecha com o algoritmo que executa tudo o que ficou dito, e
ele tem \textbf{quatro} operações --- nem uma a mais. Nenhuma é nova: cada uma é
um enunciado das secções anteriores, escrito como passo. E ele cabe em três
linhas. Sobre um ficheiro lido byte a byte, com $b$ o byte corrente:
\[
\boxed{\;
e=\operatorname{pop}(b\wedge \mathtt{0xAA}),
\qquad
o=\operatorname{pop}(b\wedge \mathtt{0x55}),
\qquad
\text{sai } [\,e+o,\ e\,].
\;}
\]
As duas máscaras são \textbf{complementares} --- $\mathtt{0xAA}\vee\mathtt{0x55}
=\mathtt{0xFF}$ e $\mathtt{0xAA}\wedge\mathtt{0x55}=0$ --- logo particionam o
byte: é a cisão em duas fases, as posições de índice par e as de índice ímpar. O
$\operatorname{pop}$ é o $\sum$. E o par que sai \emph{é o gato aplicado uma
vez}. Não há torre nenhuma a correr --- \textbf{uma} aplicação. E a saída é um
\emph{levantamento}: guardado $(e+o,\,e)$, a fase que não aparece recupera-se,
$o=(e+o)-e$. O total $e+o=\operatorname{pop}(b)$ é a conservação
(Lema~\ref{fis:lem:conserva}); a segunda coordenada é a coordenada de folha que
desdobra. \emph{Três linhas, e as quatro operações estão lá.}

\begin{center}\small
\begin{tabular}{@{}llll@{}}
\toprule
 & a operação & o que faz & de onde vem \\
\midrule
$\oplus$ & \textbf{cisão} & agrupar em $n$ fases & Def.~\ref{fis:def:objeto}: é a realização $\pi$ \\
$\sum$ & \textbf{soma} & \emph{popcount} --- o Kirchhoff & Lema~\ref{fis:lem:conserva}: $\sum_x G(x)=\abs{I}$ \\
$\otimes$ & \textbf{o gato} & \textbf{sobe} --- a convolução & Teor.~\ref{fis:thm:shift}: $\zeta$ acumula \\
$\oslash$ & \textbf{o esquilo} & \textbf{desce} --- a deconvolução & Teor.~\ref{fis:thm:mu}: $\mu=\zeta^{-1}$ \\
\bottomrule
\end{tabular}
\end{center}

\begin{definicao}[o gato e o esquilo]\label{fis:def:gato}
Fixado um inteiro $m\ge1$ --- o \textbf{metal} --- o \textbf{gato} é
$\otimes\colon\Zc^{n}\to\Zc^{n}$,
\[
(\otimes c)_0=m\,c_0+c_{n-1},\qquad (\otimes c)_i=c_{i-1}\ \ (1\le i\le n-1),
\]
e o \textbf{esquilo} é $\oslash\colon\Zc^{n}\to\Zc^{n}$,
\[
(\oslash d)_j=d_{j+1}\ \ (0\le j\le n-2),\qquad
(\oslash d)_{n-1}=d_0-m\,d_1 .
\]
Na base canónica $\otimes$ é a matriz \emph{companheira} de
$p(x)=x^{n}-m\,x^{n-1}-1$.
\end{definicao}

\begin{teorema}[a volta é exata, e é inteira]\label{fis:thm:gato}
Com a Def.~\ref{fis:def:gato}:
\begin{enumerate}\itemsep3pt
\item[(1)] $\oslash\circ\otimes=\mathrm{id}$ em $\Zc^{n}$, para todo $n\ge2$ e
todo $m$ --- \textbf{sem hipótese nenhuma sobre o vetor}.
\item[(2)] $\abs{\det\otimes}=1$: a inversa tem entradas \emph{inteiras}, e a
volta não sai do anel.
\item[(3)] Em $n=2$ o polinómio é $x^{2}=mx+1$, cujas raízes cumprem
\[
\sigma+\sigma'=m\ (\text{o traço}),\qquad \sigma\sigma'=-1\ (\text{o determinante}),
\]
com $\abs{\sigma}>1$ e $\abs{\sigma'}<1$: um \textbf{sorvedouro} e uma
\textbf{fonte}.
\end{enumerate}
\end{teorema}

\begin{proof}
(1) Escreva-se $d=\otimes c$. Para $1\le i\le n-1$ tem-se $d_i=c_{i-1}$, logo
para $0\le j\le n-2$ vem $(\oslash d)_j=d_{j+1}=c_j$. Falta a última:
$(\oslash d)_{n-1}=d_0-m\,d_1=(m\,c_0+c_{n-1})-m\,c_0=c_{n-1}$. Todas as
coordenadas coincidem, e \emph{em nenhum passo se dividiu}.

(2) Desenvolvendo o determinante pela primeira linha, o único termo não nulo é o
da entrada $c_{n-1}$, cujo menor é a identidade de ordem $n-1$; donde
$\det=\pm1$. Uma matriz inteira de determinante $\pm1$ tem inversa inteira, que é
a fórmula de (1).

(3) É o polinómio característico em $n=2$: traço $m$ e determinante $-1$ dão
$x^{2}-mx-1$. O produto das raízes é $-1$ e a soma é $m$; de $\sigma\sigma'=-1$
vem $\sigma'=-1/\sigma$, e como $\sigma>1$ para $m\ge1$, o outro tem módulo menor
que $1$.
\end{proof}

\begin{teorema}[um bit desenha a estrutura, e o dual é a ausência]\label{fis:thm:bitunico}
Seja $n=2$ e $c=(1,0)$ --- \textbf{um} bit ligado, e o resto no neutro. Então
$\otimes^{k}c$ é o par de termos consecutivos da recorrência metálica
$u_{k+1}=m\,u_k+u_{k-1}$ com $u_0=1$, $u_{-1}=0$; em particular, para $m=1$ são
os números de Fibonacci, e a razão de termos consecutivos tende para $\sigma$. E
$\oslash^{k}\otimes^{k}c=c$: a deconvolução devolve o bit único, com resíduo $0$.
\end{teorema}

\begin{proof}
Por (1) do Teor.~\ref{fis:thm:gato}, $\oslash^{k}\otimes^{k}=\mathrm{id}$. Para a
primeira parte, em $n=2$ o gato é $(a,b)\mapsto(ma+b,\,a)$, que é exatamente o
passo da recorrência com o par a guardar dois termos consecutivos --- a
\emph{memória de ordem dois}. A razão de termos consecutivos de uma recorrência
linear de raízes $\sigma,\sigma'$ com $\abs{\sigma'}<1$ tende para $\sigma$,
porque a componente em $\sigma'$ decai.
\end{proof}

\medskip\noindent\textbf{O discreto e o contínuo, de um bit só.} O dado é
$b\in B$. A \emph{presença} $b=1$ é o único operacional; a \emph{ausência} $b=0$
é o suporte neutro, e é ela o dual. Sobre esse dado:
\[
\boxed{
\begin{array}{l}
\text{o \textbf{discreto} é o bit;}\\
\text{o \textbf{contínuo} é o rasto da convolução --- a razão que tende para }\sigma;\\
\text{o \textbf{dual} é o bit ausente.}
\end{array}}
\]
Não há aqui dois mundos a conciliar: há uma operação a correr sobre um símbolo, e
$\sigma$ aparece \textbf{construído}, como o número $\pi$ da §\ref{fis:pi}
aparece --- pela mesma aritmética inteira que constrói o resto.

\medskip\noindent\textbf{E as quatro operações fecham as quatro faces.} O
Teor.~\ref{fis:thm:escada} mostrou que a escada percorre as quatro peças das duas
equações --- $0$, $+$, $\times$, $1$ --- uma por andar. O algoritmo percorre-as
num passo:
\[
\begin{array}{llll}
\oplus\ \text{cisão} & \text{a realização} & \pi\colon I\to X & \text{o \textbf{neutro} }0\ \text{é o suporte}\\
\sum\ \text{soma} & \text{a conservação} & \sum_x G(x)=\abs{I} & \text{o }+\\
\otimes\ \text{gato} & \text{a convolução} & \zeta & \text{o }\times\\
\oslash\ \text{esquilo} & \text{a deconvolução} & \mu=\zeta^{-1} & \text{o \textbf{neutro} }1:\ \oslash\otimes=\mathrm{id}
\end{array}
\]
--- e a última linha é a \textbf{realização operacional} da segunda equação da
Def.~\ref{fis:def:op}, $\partial x\cdot x=1$, com $\otimes$ a desempenhar o papel
de $x$ e $\oslash$ o de $\partial x$. \emph{Realização, e não identidade}: a
equação vive nas faces do operador e $\oslash\circ\otimes=\mathrm{id}$ vive nos
operadores que as executam --- é o mesmo enunciado um andar acima, e escrever
«literalmente» colaria dois níveis que convém manter separados.

\medskip\noindent\textbf{E em oito fases as coordenadas são as réguas.} A cisão
agrupa em $n$ fases, e $n$ é parâmetro. O caso que interessa ao byte é $n=8$, uma
fase por \textbf{direção}, e as máscaras deixam de ser duas escolhidas para serem
as oito que já existem: pelo Teor.~\ref{fis:thm:base} as direções de aresta são os
$e_k=2^{k}$, base ortonormal, e as oito máscaras \emph{são} os $e_k$. As duas do
caso $n=2$ são o que sobra ao \emph{agrupar} estas oito duas a duas. O algoritmo
expandido é
\[
\boxed{\;
f_k \;\mathrel{+}=\; (b\gg k)\wedge 1 \quad (k=0,\dots,7),
\qquad
\text{sai } \bigl[\,\textstyle\sum_k f_k,\ f_0,\dots,f_7\,\bigr],
\;}
\]
oito somas \textbf{fixas}, sem uma única decisão: nenhum ramo, nenhum laço cujo
número de voltas dependa do valor. E a conservação é a mesma, agora uma partição
de oito: $\sum_k f_k=\sum_{\text{bytes}}\operatorname{pop}(b)$, que é o
Lema~\ref{fis:lem:conserva} lido nesta cisão.

\medido{\code{papers/aranha.c}: a forma de três linhas --- as duas máscaras
complementares, o \emph{popcount} e uma aplicação do gato.
\code{tests/neuronio.c}: a volta $\oslash\otimes=\mathrm{id}$ em toda dimensão de
$2$ a $8$; o bit único a desenhar o metal; e a razão a tender para $\sigma$.
\code{tests/aranha\_g.c} §AG1--§AG11.}

\colunas{Teor.~\ref{fis:thm:agentes},~\ref{fis:thm:handshake},~\ref{fis:thm:levant},~\ref{fis:thm:gcanon},~\ref{fis:thm:geodesica},~\ref{fis:thm:periodo},~\ref{fis:thm:leitura},~\ref{fis:thm:gato}}
        {o ciclo, o campo esparso, $g=G\,h$, o gato e o esquilo}
        {\code{tests/aranha\_n.c} §AN1--§AN2, §AN5, §AN8--§AN11, §AN18, §AN20--§AN23, §AN31; \code{tests/aranha\_inversa.c} §AI1--§AI7; \code{tests/neuronio.c}}

% ═══════════════════════════════════════════════════════════════════════════
\part{Óptica}
% ═══════════════════════════════════════════════════════════════════════════

\section{A caixa: os espelhos, e a distância entre eles}\label{fis:caixa}

\noindent \emph{A caixa não é um conjunto de números} --- é o conjunto das réguas
cuja cifra tem determinante inversível, e a cifra é o par (traço, determinante).

\begin{definicao}[a caixa e a distância]\label{fis:def:caixa}
Sobre um anel comutativo $A$, chame-se \textbf{caixa} a
\[
\mathcal{P}(A)=\bigl\{\,(B,C)\;:\;B\in A,\ C\in A^{\times}\,\bigr\},
\qquad
\Delta(B,C)=B^{2}-4C ,
\]
--- as réguas quadráticas binárias cujo determinante é \textbf{invertível}. Cada
elemento é um \textbf{espelho}. Para $A$ ordenado, a distância é
$d_{\mathcal{P}}(P,Q)=\abs{\Delta(P)-\Delta(Q)}$; o seu \emph{núcleo}
$d_{\mathcal{P}}=0\iff\Delta(P)=\Delta(Q)$ não usa a ordem e vale em qualquer
$A$.

\medskip\noindent\textbf{Duas condições com o mesmo apelido, e convém separá-las
aqui e não adiante:}
\[
\begin{array}{lll}
\textbf{determinante unidade} & C\in A^{\times} & \text{o que a caixa pede para existir}\\[2pt]
\textbf{determinante de sinal} & C=\pm1 & \text{o caso }A^{\times}=\{\pm1\}
\end{array}
\]
A primeira é a condição de \emph{pertença}, e é a que a conservação exprime: um
sistema sem perda é um cujo determinante se inverte. A segunda é mais forte, e é
sob ela que as leituras deste capítulo correm --- é ela que faz «impróprio»
querer dizer $C=-1$, e é ela que o Cor.~\ref{fis:cor:serie} usa. \emph{«Unidade»
aqui é sempre elemento invertível, e nunca o número $1$.}
\end{definicao}

\medskip\noindent\textbf{O que a caixa pede a $A$, e o que não pede.} Pede
\emph{comutatividade} (para $\Delta$ ser invariante de conjugação) e nada mais
para existir. A tricotomia das classes precisa que $A$ seja \textbf{ordenado} ---
sem ordem não há três classes, há uma. \emph{Nenhuma delas precisa de um conjunto
numérico concreto.}

\medskip\noindent\textbf{E uma matriz $2\times2$ que age sobre o par (posição,
ângulo) é a matriz de um sistema óptico, e $\abs{\det}=1$ é a
\textbf{conservação}: o sistema não perde nem ganha.} A caixa é, portanto, o
conjunto dos sistemas \emph{sem perda} --- e não por analogia: é a mesma condição
escrita duas vezes. Ela tem duas coordenadas, que se leem:
\[
\begin{array}{lll}
C=\det & \text{a \textbf{orientação}} & +1\ \text{próprio};\ -1\ \text{\textbf{impróprio}, reflete}\\[2pt]
\operatorname{sgn}\Delta & \text{o \textbf{regime}} & <0\ \text{roda};\ =0\ \text{cisalha};\ >0\ \text{expande}
\end{array}
\]

\medskip\noindent\textbf{A caixa tem duas réguas, e elas não são intermutáveis.} A
$d_{\mathcal{P}}$ é métrica mas \textbf{não é ultramétrica}: com
$\Delta=20,13,8$ vem $d_{\mathcal{P}}(20,8)=12$ contra $\max\{7,5\}=7$. A
ultramétrica da caixa é a da Def.~\ref{fis:def:arvore} lida sobre os mesmos
$\Delta$, e essa cumpre o Lema~\ref{fis:lem:ultra}. A diferença não é de precisão:
é de \textbf{alcance}.
\begin{center}\footnotesize
\begin{tabular}{@{}lll@{}}
\toprule
 & $d_{\mathcal{P}}=\abs{\Delta_P-\Delta_Q}$ & $d_{\mathrm{u}}=2^{-\operatorname{prof}}$ \\
\midrule
alcance & \textbf{local}: entre dois corpos & \textbf{global}: sobre a caixa toda \\
num percurso & \textbf{acumula} (triangular) & \textbf{absorve} (Teor.~\ref{fis:thm:navega}) \\
o que dá & \emph{quanto} falta --- a recusa & a \textbf{partição} (Teor.~\ref{fis:thm:bolas}) \\
onde se usa & dentro de uma vizinhança & para navegar entre classes \\
\bottomrule
\end{tabular}
\end{center}
\noindent \emph{Uma vê o caminho; a outra só o pior degrau dele.}

\medido{\code{tests/pgwire.c} §W211: sobre os $26^3=\mathbf{17\,576}$ triplos dos
espelhos de $\abs{m},\abs{t}\le6$, a aritmética viola
$d(x,z)\le\max\{d(x,y),d(y,z)\}$ em $\mathbf{4448}$ deles e a ultramétrica em
$\mathbf{0}$. E no percurso $\Delta=-4\to5\to0\to40\to-3$ a local soma
$\mathbf{97}$ para uma distância directa de $\mathbf{1}$, enquanto a global mede
$2^{-31}$ entre as pontas contra um pior passo de $2^{0}$. Exige-se o \emph{par},
porque só uma das duas metades não distinguiria as réguas.}

\begin{teorema}[o que a caixa separa, e o que ela liga]\label{fis:thm:pandora}
\begin{enumerate}\itemsep3pt
\item[(1)] \textbf{A classe é grosseira, o espelho é fino.} A aplicação
$P\mapsto\operatorname{sgn}\Delta(P)$ tem \emph{três} valores, e cada uma tem
infinitos espelhos: $\Delta=5$ e $\Delta=13$ são \emph{corpos distintos da mesma
classe hiperbólica}. Logo $d_{\mathcal{P}}=0$ é estritamente mais forte do que ter
a mesma classe.
\item[(2)] \textbf{Distância zero é mudança de gerador, e ela escreve-se.} Se
$\Delta(P)=\Delta(Q)$, as raízes vivem na mesma extensão e diferem por
$2t=B_P-B_Q$. Onde esse $t$ está em $A$, o cisalhamento
$\varphi_t=\left(\begin{smallmatrix}1&t\\0&1\end{smallmatrix}\right)$ é a
\textbf{mudança de base}, e nela a multiplicação por $\alpha_P$ é $M_Q+tI$, de
traço $B_P$ e determinante $C_P$. \emph{A travessia não é dada por existência: é
emitida.}
\item[(3)] \textbf{Distância não nula é ausência de travessia.} Com
$\Delta(P)\neq\Delta(Q)$ não há isomorfismo de réguas, e uma comparação que
atravesse as duas não é uma pergunta difícil: é uma pergunta \emph{mal posta}. O
que se deve fazer é \textbf{recusar}, e a distância diz \emph{quão} longe se
estava de poder responder.
\end{enumerate}
\end{teorema}

\begin{proof}
(1) $\operatorname{sgn}$ tem três valores e a tricotomia da ordem esgota. Que
cada classe é infinita vê-se na família $\Delta=m^{2}+4$, estritamente crescente
em $m\ge0$. A recíproca falha em $\Delta=5,13$.

(2) De $\Delta=B^{2}-4C$ vem $B_P^{2}-B_Q^{2}\in(4)$, logo $B_P-B_Q\in(2)$ e $t$
está em $A$ --- \emph{esta} é a condição, e ela é sobre como o $2$ se comporta no
anel. O cisalhamento é unimodular, e nele a multiplicação por $\alpha_P$
escreve-se $M_Q+tI$, cujo traço é $B_Q+2t=B_P$ e cujo determinante é $C_P$. \emph{E
não é uma conjugação}: traço e determinante são invariantes de conjugação, pelo
que $B_P\neq B_Q$ exclui qualquer conjugação --- o que $\varphi_t$ faz é
transladar o gerador, e é por isso que o traço muda.

(3) Se houvesse um isomorfismo de réguas ele conjugaria as companheiras, e o
$\Delta$ --- polinómio simétrico em traço e determinante, ambos invariantes por
conjugação --- teria de coincidir.
\end{proof}

\medskip\noindent\textbf{E ela leva o corpo, não o sistema.} O $\Delta$ igual diz
que as raízes vivem na \textbf{mesma extensão}, e é isso que a travessia
transporta. A \emph{orientação} não: $M\mapsto M+tI$ é uma translação e não uma
composição, pelo que o determinante \textbf{muda}. Existem espelhos de $\Delta$
igual com $C$ oposto --- $\Delta=5$ tem os áureos $(\pm1,-1)$ e os cristalinos
$(\pm3,+1)$ --- e entre eles a travessia existe enquanto um \emph{reflete} e o
outro não. \emph{Opticamente não são o mesmo sistema, ainda que sejam o mesmo
corpo.}

\medido{\code{tests/pgwire.c} §W213: sobre \emph{todos} os pares de corpos com
$\abs{m},\abs{t}\le40$ dos dois tipos, há $\mathbf{168}$ com $\Delta$ igual e $B$
diferente. Em $\mathbf{0}$ deles a diferença é ímpar --- \emph{a condição do~(2)
não morde nesta casa, e isso varre-se em vez de se supor} --- e a testemunha
$M_Q+tI=(B_P,C_P)$ fecha nos $\mathbf{168}$. Oito são \emph{cruzados}, áureo
contra cristalino, que é o que impede a varredura de ser vazia. O controlo move o
$t$ de um e exige que a cifra deixe de bater.}

\section{A lei do discriminante, e as ordens que uma rede admite}\label{fis:leidisc}

\begin{lema}[a rede só admite cinco ordens]\label{fis:lem:cristal}
Se $A$ é uma matriz $2\times2$ de entradas \emph{inteiras} com $\abs{\det A}=1$ e
$A^{k}=I$ para algum $k\ge1$, então $k\in\{1,2,3,4,6\}$. Em particular
\textbf{não há ordem $5$}.
\end{lema}

\begin{proof}
O polinómio característico é $\lambda^{2}-t\lambda+d$ com $t=\operatorname{tr}A$
e $d=\det A=\pm1$, ambos \textbf{inteiros} por $A$ o ser. De $A^{k}=I$ segue
$\lambda^{k}=1$ para cada valor próprio, logo $\abs{\lambda}=1$ e
$\lambda_{1}\lambda_{2}=d$; e o polinómio mínimo divide $\lambda^{k}-1$, que não
tem raiz múltipla, portanto $A$ \textbf{diagonaliza}.

Se $d=-1$, duas raízes de módulo $1$ com produto $-1$ não podem ser conjugadas
--- o produto daria $\abs{\lambda}^{2}=+1$ --- logo são reais, $\{+1,-1\}$, e
$A^{2}=I$: ordem $2$.

Se $d=+1$, as raízes são conjugadas ou ambas reais, e em qualquer caso
$t=\lambda_{1}+\lambda_{2}$ é real com $\abs{t}\le2$. Um inteiro nesse intervalo
só tem cinco valores, e cada um fixa o polinómio:
\[
\begin{array}{cclr}
t=\phantom{-}2 & (\lambda-1)^{2} & A=I, & k=1\\
t=\phantom{-}1 & \lambda^{2}-\lambda+1 & \lambda=e^{\pm i\pi/3}, & k=6\\
t=\phantom{-}0 & \lambda^{2}+1 & \lambda=\pm i, & k=4\\
t=-1 & \lambda^{2}+\lambda+1 & \lambda=e^{\pm2i\pi/3}, & k=3\\
t=-2 & (\lambda+1)^{2} & A=-I, & k=2
\end{array}
\]
--- nos três do meio as raízes são primitivas das ordens $6,4,3$; nos dois
extremos a raiz é dupla e a diagonalização força $A=\pm I$. A ordem $5$ exigiria
$t=2\cos(2\pi/5)=\tfrac{\sqrt5-1}{2}$, que \textbf{não é inteiro}.
\end{proof}

\medskip\noindent \emph{É o traço que fecha a lista}: a inteireza da rede passa
para $t$, e um intervalo de comprimento $4$ só tem cinco inteiros. E vale marcar
onde a hipótese é \textbf{essencial} e onde é \textbf{acidental}: no
Teor.~\ref{fis:thm:leidisc} a inteireza é acidental --- o que a tricotomia usa é a
\emph{ordem}, e a prova corre igual em qualquer anel ordenado; aqui é essencial
--- o argumento é a \emph{discretude}, e tirá-la não generaliza o lema,
destrói-o.

\begin{teorema}[a lei do discriminante: três classes]\label{fis:thm:leidisc}
Seja $A$ com $\abs{\det A}=1$ --- \emph{e esta hipótese é a que transforma
«hiperbólico» em expande/contrai} --- sobre um anel ordenado, e
$\Delta=\operatorname{tr}^{2}A-4\det A$. As três possibilidades esgotam:
\[
\begin{array}{lll}
\Delta<0 & \textbf{elíptico} & \text{as raízes ficam no círculo: a órbita é limitada, \emph{roda}}\\[2pt]
\Delta=0 & \textbf{parabólico} & \text{raiz \textbf{dupla}: o corte do Teor.~\ref{fis:thm:corte} \emph{terminado}}\\[2pt]
\Delta>0 & \textbf{hiperbólico} & \text{duas direções reais: uma expande, a outra contrai}
\end{array}
\]
\end{teorema}

\begin{proof}
A tricotomia \emph{da ordem} esgota, e é ela --- e não uma realização numérica ---
que dá as três possibilidades. Se $\Delta<0$ as raízes são conjugadas com produto
$\det$ e módulo $\sqrt{\abs{\det}}=1$: ficam no círculo, e a órbita não cresce. Se
$\Delta=0$ há raiz dupla, isto é $\delta=0$ no Teor.~\ref{fis:thm:corte}(3), que é
onde o processo \textbf{chega} --- e aí as duas faces partilham o ponto fixo. Se
$\Delta>0$ as raízes são reais e distintas, e como o produto tem módulo $1$, uma
tem módulo maior que $1$ e a outra menor.
\end{proof}

\begin{corolario}[e sobre uma rede inteira o elíptico tem ordem finita]\label{fis:cor:ordem}
Se além disso $A$ tem \textbf{entradas inteiras}, então
\[
\boxed{\;\Delta<0\;\Longrightarrow\;\operatorname{ord}(A)\in\{3,4,6\}\;}
\]
\end{corolario}

\begin{proof}
Por $\Delta<0$ as raízes estão no círculo e são conjugadas; sendo $A$ inteira, o
traço é inteiro e o Lema~\ref{fis:lem:cristal} aplica-se, dando a ordem em
$\{1,2,3,4,6\}$. As ordens $1$ e $2$ ficam de fora: $A=\pm I$ tem $\Delta=0$, e
uma involução com $A\neq\pm I$ tem $\operatorname{tr}=0$ e $\det=-1$, donde
$\Delta=4>0$.
\end{proof}

\medskip\noindent\textbf{E as duas camadas ficam separadas, que é o que importa.}
O teorema é sobre a \emph{ordem}: dá os três \textbf{regimes}, e corre em qualquer
anel ordenado. O corolário é sobre a \emph{discretude}: dá a \textbf{ordem
finita}, e sem a inteireza não há lista nenhuma --- num anel ordenado geral,
$\Delta<0$ põe as raízes no círculo mas não as obriga a serem raízes da unidade.
\[
\boxed{\;\text{o \textbf{regime} é da ordem; a \textbf{ordem finita} é da discretude}\;}
\]
É a mesma distinção que o Lema~\ref{fis:lem:cristal} já fazia ao dizer que ali a
inteireza é \emph{essencial} --- agora dita também no sítio onde ela era usada
sem estar na hipótese.

\medskip\noindent\textbf{E o traço sozinho não classifica.} Onde $\det=+1$, a
condição $\Delta<0$ é $\abs{\operatorname{tr}}<2$; mas com $\det=-1$ vale
$\Delta=\operatorname{tr}^{2}+4>0$ \emph{sempre}, e a matriz é hiperbólica
qualquer que seja o traço. \emph{Quem classifica é $\Delta$.}

\medskip\noindent\textbf{E o $\Delta$ é o critério de estabilidade de um
ressoador --- não um análogo dele.} Com $\det=+1$ a órbita é limitada exactamente
quando $\abs{\operatorname{tr}}<2$, isto é $\operatorname{tr}^{2}-4\det<0$, isto é
$\Delta<0$. \emph{As três classes são os três regimes: estável, marginal,
instável.} E o que cada regime faz a um raio distingue-se sem ambiguidade:
seguindo a órbita de $(1,0)$ sob a companheira, a norma \emph{fica} no elíptico,
cresce \emph{linearmente} no parabólico e \emph{geometricamente} no hiperbólico
--- e no hiperbólico os números são os de Fibonacci, que é o ouro a aparecer onde
a expansão é.

\medskip\noindent\textbf{E a tríade atravessa duas classes.} Os operadores
separam-se pela \emph{ordem} --- o menor $k$ com $f^{k}=\mathrm{id}$:
\[
\begin{array}{llcl}
\text{Cantor / Julia} & \text{o espelho da dobra} & 2 & \nu^{2}=\mathrm{id}\\
\text{o trial} & \tau\in\{-1,0,+1\} & 3 & \tau^{3}=\mathrm{id}\\
\text{Viviani} & i,\ \text{o meio-ângulo} & 4 & i^{4}=\mathrm{id}
\end{array}
\]
A dobra da Def.~\ref{fis:def:op} é a \textbf{primeira linha}: $\partial\partial=
\mathrm{id}$, e é por isso que ela espelha. \emph{Viviani é a sua raiz} --- o que
ao quadrado dá o espelho --- e por isso tem o dobro da ordem: girar duas vezes um
quarto de volta é espelhar. O trial (ordem $3$) e Viviani (ordem $4$)
\emph{rodam}, e são elípticos. O espelho $\nu$ (ordem $2$) \textbf{não roda:
reflete} --- a sua testemunha é
$\left(\begin{smallmatrix}0&1\\1&0\end{smallmatrix}\right)$, de traço $0$ e
$\det=-1$, logo $\Delta=4>0$, hiperbólica. \emph{A diferença entre
$\{1,2,3,4,6\}$ e $\{3,4,6\}$ é exatamente o que reflete em vez de rodar.} E o
$6$ é onde as duas primeiras voltam juntas, $\operatorname{lcm}(2,3)$ --- o toro é
o sítio onde elas vivem juntas.

\medido{\code{tests/pgwire.c} §W116: as três classes esgotam $\mathbf{360}$
matrizes de entradas em $-4..4$; as elípticas têm todas ordem finita e dão
$\{3,4,6\}$, com a $5$ ausente e a $2$ também; $\abs{\operatorname{tr}}=2$ e
$\Delta=0$ coincidem nas $\mathbf{180}$ de $\det+1$; e a família metálica é toda
hiperbólica, com $\det A_m=-1$ e $\operatorname{tr}A_m^{2}=m^{2}+2$ pelo motor.
\emph{A primeira escrita classificou pelo traço sem separar por determinante e
produziu $16$ «elípticas» sem ordem finita --- a asserção apanhou-as, e a correção
é o próprio resultado.} §W219: varridos $B\in[-300,300]$, a órbita de $(1,0)$ dá
$1,1,1,1,1,1$ no elíptico, $1,2,3,4,5,6$ no parabólico e $1,1,2,3,5,8$ no
hiperbólico; e as ordens do elíptico saem $3$, $4$ e $6$, \emph{distintas entre
si}, o que impede a asserção de passar por acaso.}

\section{O impróprio é sempre instável, e a orientação compõe-se}\label{fis:improprio}

\begin{corolario}[o impróprio é sempre instável]\label{fis:cor:optica}
Se $C=-1$ então $\Delta=B^{2}+4\ge4>0$, qualquer que seja $B$. Logo um espelho que
\emph{reflete} é \textbf{sempre} hiperbólico: nunca roda nem cisalha. Das seis
combinações de (orientação, regime) existem \textbf{quatro} --- o impróprio ocupa
uma das três, e o próprio ocupa as três.
\end{corolario}

\begin{proof}
Com $C=-1$ vem $\Delta=B^{2}-4C=B^{2}+4$, e $B^{2}\ge0$ em qualquer anel
ordenado, donde $\Delta\ge4>0$ para todo $B$: o Teor.~\ref{fis:thm:leidisc}
põe-no na classe hiperbólica, sem excepção. Logo das três combinações que
caberiam ao impróprio ele realiza \emph{uma}. Já com $C=+1$ vem
$\Delta=B^{2}-4$, que é negativo em $\abs B<2$, nulo em $\abs B=2$ e positivo em
$\abs B>2$: o próprio realiza as \emph{três}. Total $1+3=4$.
\end{proof}

\medskip\noindent O conteúdo é a leitura: \emph{a inversão de orientação e o
confinamento excluem-se}. É a mesma frase que separa o que roda do que não roda,
vista agora como \textbf{impossibilidade}.

\medskip\noindent\textbf{E há duas fontes de órbita finita, não uma.} O
Cor.~\ref{fis:cor:ordem} diz que $\Delta<0$ dá ordem finita em $\{3,4,6\}$; o
Lema~\ref{fis:lem:cristal} diz que ordem finita está em $\{1,2,3,4,6\}$. \emph{As
duas listas diferem no $\{1,2\}$}, e é fácil ler a primeira como se fosse a
caracterização --- «ordem finita $\iff$ elíptico». Não é: o espelho $\nu$ tem
$\Delta=4>0$ e $\nu^{2}=\mathrm{id}$. \textbf{Reflectir duas vezes é a
identidade}, e isso fecha a órbita sem rodar nada. As duas fontes são
\textbf{disjuntas}, e quem o garante é este corolário: o que reflete é impróprio,
logo hiperbólico, logo não é elíptico. \emph{Rodar fecha a órbita em $\{3,4,6\}$;
reflectir fecha-a em $2$, e nunca há um mesmo espelho a fazer as duas coisas.}

\begin{corolario}[a orientação compõe-se, o regime não]\label{fis:cor:serie}
A aplicação $\det\colon\mathcal{P}\to A^{\times}$ é um \textbf{homomorfismo}, e a
caixa é fechada para a composição porque unidade vezes unidade é unidade; em
particular, com $A^{\times}=\{\pm1\}$, o produto de dois espelhos impróprios é
\textbf{próprio} --- \emph{duas reflexões são uma rotação}. Já
$\operatorname{sgn}\Delta$ \textbf{não é função} das classes dos factores: existem
$X,Y$ elípticos com $XY$ hiperbólico.
\end{corolario}

\begin{proof}
A primeira metade é a multiplicatividade do determinante. A segunda é um
contra-exemplo, e ele está medido abaixo.
\end{proof}

\medskip\noindent \emph{A razão está na natureza dos dois invariantes.} O
determinante é invariante de \textbf{composição}; o $\Delta$ é invariante de
\textbf{conjugação} --- classifica o \emph{elemento}, não o produto. Por isso a
caixa não é graduada pelo $\Delta$, e por isso \textbf{compor dois sistemas
estáveis pode dar um instável}: a estabilidade não é uma propriedade que se
transporte ao longo de uma série.

\medido{\code{tests/pgwire.c} §W220: nas $\mathbf{2500}$ composições dos espelhos
de $\abs{B}\le12$ dos dois tipos, $\det(XY)=\det X\cdot\det Y$ em \textbf{todas}.
E das nove combinações de classes, \textbf{sete} não determinam a do produto ---
incluindo elíptico$\,\circ\,$elíptico, que produz as três. Só
parabólico$\,\circ\,$hiperbólico e o seu simétrico saem sempre hiperbólicos. O
controlo compõe um impróprio com um próprio e exige $\det=-1$: sem ele, «é
homomorfismo» não se distinguiria de «$\det$ é sempre $+1$». §W219: dos
impróprios, \textbf{zero} são elípticos, \textbf{zero} são parabólicos e os
$\mathbf{601}$ restantes hiperbólicos; das seis combinações existem
\textbf{quatro}, e o número é asserido --- \emph{escrevê-lo sem o contar foi o que
quase deixou passar um «três»}. §W225: as elípticas dão ordens em $\{3,4,6\}$ e as
três realizam-se; a única não-elíptica que fecha dá $\mathbf{2}$, e o controlo
varre os elípticos de $\abs{B}\le60$ à procura de um de ordem $2$ e exige zero.}

\section{A borda: onde a régua colapsa numa forma linear}\label{fis:borda}

\noindent A régua $q(a,b)=a^{2}+Bab+Cb^{2}$ é \textbf{multiplicativa} ---
$q(xy)=q(x)q(y)$ --- e uma máquina que só saiba somar não a calcula: somar dá
progressão aritmética, multiplicar dá geométrica, e a ponte entre as duas é
precisamente o que falta. \emph{Esta é a razão pela qual uma realização pode
ordenar por $q$ e não conseguir comparar por $q$}, e ela é estrutural, não uma
falha de esforço.

\begin{corolario}[a régua factoriza exactamente na borda]\label{fis:cor:borda}
Sobre um anel comutativo $A$, a forma $q$ do espelho $(B,C)$ escreve-se como
produto de duas formas \emph{lineares} sobre $A$,
\[
q(a,b)=(a+p\,b)(a+q'b),\qquad p+q'=B,\quad p\,q'=C ,
\]
se e só se $\Delta$ é um quadrado em $A$ e as metades $(B\pm\sqrt{\Delta})/2$
estão em $A$. Em particular: $\Delta=0$ --- a classe \textbf{parabólica} --- dá
$p=q'=B/2$ e $q$ é o \textbf{quadrado de uma única linear}; e $\Delta$ quadrado
não nulo dá duas lineares distintas, com o \emph{sinal} de $q$ a ser o produto
dos sinais delas.
\end{corolario}

\begin{proof}
As raízes de $x^{2}-Bx+C$ são $(B\pm\sqrt{\Delta})/2$, e
$q(a,b)=b^{2}\,p(a/b)$ com $p$ o polinómio característico; a factorização existe
exactamente quando as raízes estão em $A$. A condição sobre as metades é a mesma
do Teor.~\ref{fis:thm:pandora}(2) --- o comportamento do $2$ no anel --- e
reaparece aqui sem ser convidada.
\end{proof}

\medskip\noindent\textbf{E é isso que uma máquina aditiva alcança.} No caso
parabólico a comparação $q\ \mathrm{OP}\ k$ é
$\abs{a+\tfrac{B}{2}b}\ \mathrm{OP}\ \sqrt{k}$, com $\sqrt{k}$ calculado
\emph{antes} de correr: cabe inteira, sem uma multiplicação em execução. No outro
caso o valor ainda pede produto, mas o sinal não --- as duas rectas $L_1=0$ e
$L_2=0$ \textbf{cortam} o plano em quatro partes e $q$ tem sinal constante em cada
uma. \emph{E a borda é estreita}: as classes são três e a borda não é uma delas.

\medskip\noindent\textbf{E a ponte entre as duas progressões já está construída:
é a AGM.} A §\ref{fis:corte} formaliza o par $(a,b)\mapsto(m,g)$, e diz o que ele
custa: \emph{uma face endereça, duas produzem}. É literalmente a ponte que falta
ao avaliador, e a borda lê-se nas suas coordenadas. Com $\delta=\tfrac{a-b}{2}$, a
identidade $g_{\ast}^{2}=ab=m^{2}-\delta^{2}=(m+\delta)(m-\delta)$ \emph{é} a diferença
de quadrados, e daí as duas bordas saem em forma fechada:
\[
\Delta=0:\ q=(2m)^{2}
\qquad\qquad
\Delta=4:\ q=4\,m\,\delta .
\]
\emph{A classe parabólica usa só a face aritmética; a borda hiperbólica usa as
duas.} É a dicotomia do Teor.~\ref{fis:thm:duas}(4) lida na caixa.

\medskip\noindent\textbf{E a borda $\Delta=4$ já tinha nome: é o espelho $\nu$.} A
testemunha da reflexão é
$\left(\begin{smallmatrix}0&1\\1&0\end{smallmatrix}\right)$, de traço $0$ e
$\det=-1$ --- cifra $(0,-1)$. Com $m=0$ a borda dá $\sigma^{2}=1$: uma
\textbf{involução}, e os seus vectores próprios exibem-se sem conta nenhuma:
\[
\nu(1,1)=(1,1),\qquad \nu(1,-1)=-(1,-1) .
\]
As duas rectas $a+b=0$ e $a-b=0$ não são um artefacto da conta: são os
\textbf{eixos próprios} de $\nu$, e é sobre eles que a forma factoriza ---
diagonalizar uma forma \textbf{é} factorizá-la.

\medskip\noindent\textbf{E o $2$ volta a aparecer, sem ter sido convidado.} A
matriz de mudança para essa base é
$\left(\begin{smallmatrix}1&1\\1&-1\end{smallmatrix}\right)$, de determinante
$-2$: ela \emph{não é unimodular} se o $2$ não for invertível em $A$. Portanto a
afirmação certa não é «uma involução diagonaliza sobre qualquer anel» --- é esta:
\emph{os eixos próprios de $\nu$ estão em $A$ sempre, e a factorização passa a
correr sobre $A$ quando a divisão por $2$ é permitida}. É a mesma condição que o
Cor.~\ref{fis:cor:borda} pede às metades $(B\pm\sqrt\Delta)/2$ e que o
Teor.~\ref{fis:thm:pandora}(2) pede ao $t$ --- \emph{o comportamento do $2$ no
anel}, a aparecer pela terceira vez e sempre no mesmo sítio.
\emph{Reflectir é ser a própria inversa, e é isso que a factorização exprime.}

\medskip\noindent\textbf{No interior e na fronteira.} Dentro de uma classe de
$\Delta$ igual, o Teor.~\ref{fis:thm:pandora}(2) emite a mudança de gerador: é o
\emph{mesmo} sistema visto noutro referencial. Entre $\Delta$ diferentes não há
travessia, e a recusa diz o que isso é opticamente: \emph{não são o mesmo
aparelho}. Já na \textbf{borda} os eixos próprios estão no próprio anel --- com a
ressalva sobre o $2$ dita acima --- e fora dela vivem na extensão: real
quadrática se $\Delta>0$, \emph{complexa conjugada} se $\Delta<0$. É por isso que a classe $\Delta<0$ se chama elíptica, e não é
coincidência de nomenclatura: \emph{os eixos são os mesmos das cónicas}.

\begin{corolario}[o primeiro corte da árvore é o elíptico, e só ele]\label{fis:cor:pandbola}
Na ultramétrica da caixa, a bola de raio $1$ --- os espelhos que \emph{não}
divergem no bit mais alto --- é exatamente $\{\Delta<0\}$ e o seu complemento.
Logo o primeiro nível da árvore separa \textbf{elíptico} de
\textbf{não-elíptico}, e mais nada: a distinção entre \emph{parabólico} e
\emph{hiperbólico} vive \textbf{mais fundo}.
\end{corolario}

\begin{proof}
O bit mais alto da representação de $\Delta$ é o do sinal, pelo que
$\operatorname{prof}(\Delta_P,\Delta_Q)=0$ se e só se os sinais diferem. Que
$\Delta=0$ não forma bola vê-se por não haver raio que o isole dos positivos.
\end{proof}

\medskip\noindent \emph{Das três classes, só a elíptica é uma bola.} E o corte
que a árvore faz primeiro é justamente o que separa \emph{o que roda} de \emph{o
que não roda} --- a régua ultramétrica encontra essa fronteira no
\textbf{primeiro} bit, sem que ninguém lha tenha dito.

\medskip\noindent\textbf{E o catálogo é uma realização, com o seu campo $G$.} A
aplicação $\pi\colon\text{nome}\mapsto\text{cifra}$ é uma realização no sentido da
Def.~\ref{fis:def:objeto}, e onde dois nomes dão a mesma cifra eles são \emph{o
mesmo espelho} --- que é o Teor.~\ref{fis:thm:mult}(2) lido no catálogo. As duas
leis que a obra prova em abstracto ganham aqui um objecto com nomes próprios: a
escada $\sum_x G(x)=\abs{I}$ e a folga $\sum_x(G(x)-1)=\abs{I}-\abs{X}$. E os
nomes batem com a leitura sem terem sido escolhidos por ela: o corpo a que a obra
chamou \code{optico} é $(0,1)$ com $\Delta=-4$ --- \textbf{elíptico}, o que roda,
e é isso que um sistema óptico redondo faz; o \code{criativo} é $(0,-1)$ com
$\Delta=4$, que é o espelho $\nu$, e está na borda; o \code{racional} é $(2,1)$
com $\Delta=0$: o \textbf{corte} parabólico. \emph{Não foram baptizados por esta
secção; a secção leu-os.}

\medskip\noindent\textbf{E é daqui que vem o nome.} Abrir a caixa é misturar os
espelhos: comparar através de dois com $\Delta$ diferente devolve sempre
\emph{algum} número, e esse número não é sobre objeto nenhum. \emph{O que sai da
caixa não é um erro que se veja --- é uma resposta com a cara de certa.} Por isso
a regra é a do Teor.~\ref{fis:thm:pandora}(3), e a distância existe para que a
recusa seja \textbf{quantitativa}: diz-se quanto falta, e não apenas que falta.

\medido{\code{tests/pgwire.c} §W214: varridos os corpos de $\abs{m},\abs{t}\le12$
dos dois tipos --- $50$ --- apenas $\mathbf{3}$ factorizam (os dois cristalinos de
$\Delta=0$ e o áureo de $\Delta=4$, onde o ouro colapsa na unidade) e os outros
$\mathbf{47}$ são irredutíveis. §W224: varrido $\abs{B}\le3000$ nos dois sinais de
$C$, e depois $C\in[-40,40]$ com $\abs{B}\le200$: em todos os casos de $\Delta$
quadrado --- $\mathbf{723}$ no segundo varrimento --- as metades estão no anel.
§W218: dos $26$ espelhos, o critério da cifra e o do Cor.~\ref{fis:cor:borda}
coincidem em $\mathbf{25}$ e divergem em $\mathbf{1}$ --- o $(0,1)$ de $\Delta=-4$,
onde a forma factoriza sobre a \emph{extensão} e não sobre a base. \emph{É a
hipótese «sobre um anel comutativo $A$» a fazer trabalho: uma testemunha que o
ignore acerta quase sempre --- que é a forma de testemunha mais perigosa, a que
quase serve.} §W211: as três condições do Cor.~\ref{fis:cor:pandbola} verificam-se
nos $26$ espelhos, e a recíproca ingénua falha em $\mathbf{84}$ dos $676$ pares,
\emph{todas} com $\Delta=0$. §W222--§W223: os $\mathbf{44}$ corpos do catálogo,
\textbf{zero} fora da caixa e \textbf{zero} violações do Cor.~\ref{fis:cor:optica};
e os $44$ nomes caem em $\mathbf{18}$ cifras, com $\sum G=44=\abs{I}$ e
$\sum(G-1)=26=\abs{I}-\abs{X}$, a maior fibra com \textbf{nove} nomes numa cifra
só. §W215: com o segundo andar da célula nomeado, a régua escreve-se e responde
certo nos quatro limiares. §W212: os bits davam a \textbf{inversa} do valor, e com
a régua a mesma descida devolve a ordem certa --- \emph{é a régua que ordena, e não
os bits}. \code{banco/sql.c} \code{checa\_corpos}: a recusa com $\Delta=5$ contra
$\Delta=13$, distância $\mathbf{8}$, \emph{sem} invocar classe --- os dois são
hiperbólicos; e com distância $\mathbf{0}$ o motor não recusa: \textbf{emite} a
travessia.}

\colunas{Teor.~\ref{fis:thm:pandora},~\ref{fis:thm:leidisc}; Lema~\ref{fis:lem:cristal}; Cor.~\ref{fis:cor:optica},~\ref{fis:cor:serie},~\ref{fis:cor:borda},~\ref{fis:cor:pandbola}}
        {a caixa, as duas coordenadas, as três órbitas de $(1,0)$, a borda, o catálogo}
        {\code{tests/pgwire.c} §W116, §W211--§W215, §W218--§W220, §W222--§W225; \code{banco/sql.c}}

% ═══════════════════════════════════════════════════════════════════════════
\part{Partículas e Núcleo}
% ═══════════════════════════════════════════════════════════════════════════

\section*{A tese desta parte}\label{fis:tese}

\noindent \textbf{Uma partícula não é um objeto fundamental: é um modo espectral
de uma cadeia.} O que se chama \emph{massa} é a distância logarítmica de uma raiz
ao círculo unitário; o que se chama \emph{carga} é uma ocupação; o que se chama
\emph{ligação} é a multiplicidade que a realização produziu; o que se chama
\emph{decaimento} é uma órbita; e o que se chama \emph{estabilidade} é um ponto
fixo.

\medskip\noindent Nada disto é acrescentado à construção. É a construção lida:

\begin{center}\small
\begin{tabular}{@{}lll@{}}
\toprule
& \textbf{a estrutura} & \textbf{o que a obriga} \\
\midrule
partícula & uma raiz do espectro & Def.~\ref{fis:def:atomo} \\
massa & $\abs{\log\rho}$ & Def.~\ref{fis:def:atomo} \\
energia de ligação & $\log\beta$, a raiz dominante & Def.~\ref{fis:def:atomo} \\
luz & $\rho=1$: o modo roda; a órbita \emph{fecha} & Teor.~\ref{fis:thm:massa},~\ref{fis:thm:orbita} \\
matéria & há raiz \emph{fora}; a órbita \emph{escapa} & Teor.~\ref{fis:thm:massa},~\ref{fis:thm:orbita} \\
modo ligado & $\rho<1$: contrai --- e é obrigatório & Prop.~\ref{fis:prop:neutro} \\
neutralidade & $\prod\rho=\pm1$, o determinante unidade & Prop.~\ref{fis:prop:neutro} \\
defeito de massa & $\abs{I}-\abs{X}=\sum(G-1)$ & Teor.~\ref{fis:thm:defeito} \\
o vale do ferro & o ponto fixo comum às duas faces & Teor.~\ref{fis:thm:ferro} \\
reação em cadeia & a iteração do operador de reprodução & Def.~\ref{fis:def:reproducao} \\
subcrítico / crítico / supercrítico & o raio espectral $k=\varrho(R)$ & Teor.~\ref{fis:thm:kappa} \\
decaimento & uma órbita: cauda e período & Teor.~\ref{fis:thm:cadeia} \\
estabilidade & um ponto fixo & Teor.~\ref{fis:thm:cadeia} \\
carga & a fibra: $Q=N/3$ & Teor.~\ref{fis:thm:carga} \\
antipartícula & a conjugação, com ponto fixo $Q=0$ & §\ref{fis:carga} \\
\bottomrule
\end{tabular}
\end{center}

\medskip\noindent\textbf{E a regra da parte é a da coluna da direita.} Não se
postulam partículas: constroem-se cadeias. Não se postulam massas: mede-se a
distância ao círculo. Não se postula a carga: contam-se ocupações. Não se
postula a estabilidade: procuram-se pontos fixos. Não se postula a ligação:
mede-se a dobra. Não se postula a luz: acha-se a fronteira espectral onde a massa
desaparece.

\medskip\noindent Cada uma dessas frases tem, à sua direita, o enunciado que a
obriga --- e é essa a diferença entre uma ontologia e um manifesto.

\section{A cadeia é o átomo}\label{fis:atomo}

\noindent A caixa da §\ref{fis:caixa} classificou espelhos de grau $2$. Subindo o
grau, o mesmo objeto tem um nome próprio: \textbf{uma cadeia de grau $N$ é um
átomo com $N$ níveis}, e os níveis são as raízes.

\begin{definicao}[o átomo]\label{fis:def:atomo}
Seja $c$ uma cadeia de grau $N$ com $\abs{c_0}=1$ --- a mesma condição de
\emph{unidade} que define a caixa --- e sejam $\rho_1,\dots,\rho_N$ as raízes da
sua companheira, ordenadas por módulo decrescente. Chama-se
\begin{center}\small
\begin{tabular}{@{}lll@{}}
\toprule
peça & o que é & a régua \\
\midrule
\textbf{núcleo} & a raiz dominante $\beta=\rho_1$ & energia de ligação $=\log\beta$ \\
\textbf{modo ligado} & uma raiz com $\rho<1$ & \emph{elétron}: cai para dentro \\
\textbf{modo livre} & uma raiz com $\rho=1$ & massa $\abs{\log\rho}=0$ \\
\bottomrule
\end{tabular}
\end{center}
e a \textbf{massa} de um modo é $\abs{\log\rho}$.
\end{definicao}

\begin{proposicao}[o átomo é neutro]\label{fis:prop:neutro}
O produto das raízes é $\prod_i\rho_i=\pm1$, e portanto
\[
\boxed{\;\sum_i \log\abs{\rho_i}=0\;}
\]
--- \emph{a soma das massas com sinal é nula}. O átomo é neutro, e a neutralidade
não é imposta: é o determinante unidade.
\end{proposicao}

\begin{proof}
O produto das raízes de uma companheira é $\pm c_0$, e $\abs{c_0}=1$ por
hipótese; logo $\abs{\prod_i\rho_i}=1$ e o logaritmo da norma anula-se. É a mesma
condição $\abs{\det}=1$ que a Def.~\ref{fis:def:caixa} pede à caixa --- e é o
Teor.~\ref{fis:thm:gato}(2), $\abs{\det\otimes}=1$, lido no espectro em vez de na
matriz.
\end{proof}

\medskip\noindent\textbf{E daqui sai a tabela dos átomos.} O que está
\emph{demonstrado} é a estrutura: grau $2$ --- a família metálica do
Teor.~\ref{fis:thm:gato}(3), com $\sigma$ o \emph{sorvedouro} e
$\sigma'=-1/\sigma$ a \emph{fonte} --- tem um modo dominante e um modo ligado
\textbf{real}; grau $3$ tem um par conjugado a espiralar, $\theta\neq0$.
\emph{O grau é o número de níveis, e a conservação da norma mantém o átomo
neutro.}

\medskip\noindent\textbf{Hipótese de identificação.} A primeira correspondência
candidata é o grau $2$ como a estrutura mínima de um átomo de um modo ligado ---
um núcleo e um elétron, sem fase --- e o grau $3$ como o primeiro que traz fase.
\emph{Estas identificações não são teoremas: são alvos de varrimento.} O que a
construção tem de produzir, sem ajuste posterior, são invariantes comparáveis aos
dos objetos candidatos --- e a §\ref{fis:mapa} é onde o varrimento começa a
responder.

\medskip\noindent E o par sorvedouro/fonte não é figura de estilo: é o
Teor.~\ref{fis:thm:leidisc} na classe hiperbólica --- \emph{duas direções reais,
uma expande e a outra contrai}. Num átomo, a que contrai é o elétron ligado e a
que expande é o núcleo.

\medido{\code{tests/neuronio.c}: a volta $\oslash\otimes=\mathrm{id}$ em toda
dimensão de $2$ a $8$, e a razão de termos consecutivos a tender para $\sigma$.
\code{tests/pgwire.c} §W116: a família metálica é toda hiperbólica, com
$\det A_m=-1$ e $\operatorname{tr}A_m^{2}=m^{2}+2$ pelo motor. A cadeia de grau
$7$ com $\abs{c_0}=1$ tem sete modos e conserva $\prod=\pm1$; o associador mede
$\approx10^{-15}$ em grau $3$ (associativo) e $\approx13$ em grau $7$.}

\section{A massa é o logaritmo do raio, e a luz é a borda}\label{fis:massa}

\begin{teorema}[a fronteira entre matéria e luz é o círculo]\label{fis:thm:massa}
Com a massa $\abs{\log\rho}$ da Def.~\ref{fis:def:atomo}:
\begin{enumerate}\itemsep3pt
\item[(1)] \textbf{Um modo tem massa nula se e só se $\rho=1$}, isto é, se e só se
a raiz está \emph{no círculo unitário}.
\item[(2)] \textbf{Matéria é haver raiz FORA.} Uma cadeia tem algum modo com
massa se e só se alguma raiz tem $\rho\neq1$ --- e, pela
Prop.~\ref{fis:prop:neutro}, isso obriga a haver raiz \emph{fora} \textbf{e}
raiz \emph{dentro}, porque $\sum\log\rho_i=0$.
\item[(3)] \textbf{Luz é estarem todas SOBRE.} Uma cadeia tem \emph{todos} os
modos de massa nula se e só se todas as raízes estão no círculo.
\end{enumerate}
\emph{O círculo unitário é a fronteira onde a matéria vira luz.}
\end{teorema}

\begin{proof}
(1) $\abs{\log\rho}=0$ se e só se $\log\rho=0$, isto é $\rho=1$. (3) é (1)
aplicada a cada raiz. Para (2): se alguma raiz tem $\rho>1$ então
$\log\rho>0$, e como a soma dos logaritmos é \emph{nula}
(Prop.~\ref{fis:prop:neutro}), alguma outra tem $\log\rho<0$, isto é $\rho<1$.
\emph{Não pode haver cadeia com todas as raízes dentro do círculo}: o produto dos
módulos seria menor que $1$, e ele vale exactamente $1$. \emph{A massa vem
sempre aos pares} --- um modo que expande obriga a um que contrai.
\end{proof}

\medskip\noindent\textbf{E daí os três regimes de um modo, sem ambiguidade:}
\[
\boxed{\;
\rho<1\ \longleftrightarrow\ \text{contrai}
\qquad
\rho=1\ \longleftrightarrow\ \text{roda}
\qquad
\rho>1\ \longleftrightarrow\ \text{expande}
\;}
\]
--- e é por isso que \emph{núcleo} e \emph{modo ligado} não são dois nomes para o
mesmo sinal: o núcleo é o que expande, o elétron é o que contrai, e a
neutralidade obriga-os a coexistir.

\medskip\noindent\textbf{E é a mesma borda da caixa.} O
Teor.~\ref{fis:thm:leidisc} classifica pelo módulo das raízes: com determinante
unidade, $\Delta<0$ põe as duas raízes \emph{sobre} o círculo --- é o
\textbf{elíptico}, o que roda, e é o regime sem massa; $\Delta>0$ tira-as de lá
--- é o \textbf{hiperbólico}, uma expande e a outra contrai, e é o regime com
massa. A fronteira $\Delta=0$ é a \emph{raiz dupla}, que o
Teor.~\ref{fis:thm:corte}(3) já tinha identificado como o sítio onde o encaixe
das duas faces \textbf{chega}.
\[
\boxed{\;\text{massa}=0\ \iff\ \rho=1\ \iff\ \text{o modo roda em vez de expandir}\;}
\]

\medskip\noindent\textbf{E o átomo de grau $10$ lê-se peça a peça.} A cadeia de
grau $10$ com $\abs{c_0}=1$ e raiz dominante mínima tem \textbf{um} núcleo,
\textbf{um} elétron e \textbf{oito} modos no círculo: núcleo e elétron com massa
$0{,}1624$ e $\theta=0$, e os oito fotões em quatro pares conjugados, com
$\theta$ a valer $160{,}6^{\circ}$, $137{,}3^{\circ}$, $62{,}8^{\circ}$ e
$107{,}0^{\circ}$. \emph{É radiação pura na borda}, e os pares conjugados são a
involução do Teor.~\ref{fis:thm:duas}(2) a agir no espectro: $\theta$ e $-\theta$
são o que a dobra troca, e o par sobrevive porque o módulo --- que é o que dá a
massa --- entra ao quadrado (Teor.~\ref{fis:thm:BI}).

\section{O mapa espectral: que átomos a álgebra permite}\label{fis:mapa}

\noindent O Teor.~\ref{fis:thm:massa} separa matéria de luz pelo módulo da raiz.
Mas \emph{avaliar} um módulo é sair dos inteiros, e não é preciso: a mesma
separação lê-se na \textbf{órbita da companheira}, que é inteira do primeiro
passo ao último.

\begin{teorema}[a dicotomia da órbita]\label{fis:thm:orbita}
Seja $c$ uma cadeia de grau $N$ com $\abs{a_0}=1$ e $C$ a sua companheira. A
órbita de $e_1$ sob $C$ --- que é uma sucessão de vectores \emph{inteiros} ---
faz uma de duas coisas, e nunca outra:
\[
\begin{array}{lll}
\textbf{fecha} & \text{volta a }e_1\text{ ao fim de }p\text{ passos}
  & \text{todas as raízes na unidade: } \textbf{luz}\\[2pt]
\textbf{escapa} & \text{as componentes crescem sem limite}
  & \text{alguma raiz fora do círculo: } \textbf{matéria}
\end{array}
\]
E quando fecha, fecha \textbf{sem cauda}: $\ell=0$.
\end{teorema}

\begin{proof}
$\det C=\pm a_0$, logo $\abs{\det C}=1$ e $C$ é invertível sobre os inteiros. Se
a órbita for limitada, ela vive num conjunto \emph{finito} de vectores inteiros,
portanto repete um estado --- e por $C$ ser invertível a repetição só pode ser em
$e_1$, donde $\ell=0$ e $C^{p}e_1=e_1$. Como $e_1$ é vector cíclico da
companheira, $C^{p}=\mathrm{id}$: a ordem é finita e todos os valores próprios
são raízes da unidade, isto é, todos os modos têm $\rho=1$ e massa nula
(Teor.~\ref{fis:thm:massa}). Se a órbita não for limitada, algum valor próprio
tem módulo maior que $1$, e esse modo tem massa. As duas alternativas esgotam, e
a decisão faz-se \emph{por igualdade de vectores inteiros} --- que é a
cláusula~(1) do Teor.~\ref{fis:thm:mult}: duplicidade é dobra, e a dobra
decide-se por igualdade.
\end{proof}

\medskip\noindent\textbf{E daí o varrimento.} Tomadas \emph{todas} as cadeias
mónicas de grau $2$ a $10$ com coeficientes no trial $\{-1,0,+1\}$ e
$\abs{a_0}=1$ --- $\mathbf{59\,046}$ delas --- a dicotomia reparte-as assim:

\begin{center}\small
\begin{tabular}{@{}lrrrrrrrrrr@{}}
\toprule
grau & $2$ & $3$ & $4$ & $5$ & $6$ & $7$ & $8$ & $9$ & $10$ & total \\
\midrule
cadeias & $6$ & $18$ & $54$ & $162$ & $486$ & $1458$ & $4374$ & $13122$ & $39366$ & $59\,046$ \\
\textbf{fecham} & $4$ & $4$ & $8$ & $10$ & $18$ & $20$ & $32$ & $32$ & $56$ & $\mathbf{184}$ \\
\textbf{escapam} & $2$ & $14$ & $46$ & $152$ & $468$ & $1438$ & $4342$ & $13090$ & $39310$ & $58\,862$ \\
\bottomrule
\end{tabular}
\end{center}

\noindent\emph{A luz é rara.} De cinquenta e nove mil cadeias, cento e oitenta e
quatro fecham --- e as que fecham não têm um único modo com massa.

\subsection*{Os períodos, e onde o cinco entra}

\noindent Quem fecha tem um \textbf{período}, e os períodos que existem em cada
grau são estes:
\begin{center}\small
\begin{tabular}{@{}rl@{}}
\toprule
grau & os períodos que a álgebra permite \\
\midrule
$2$ & $2,\ 3,\ 4,\ 6$ \\
$3$ & $3,\ 4,\ 6$ \\
$4$ & $4,\ \mathbf{5},\ 6,\ 8,\ 10,\ 12$ \\
$5$ & $5,\ 6,\ 8,\ 10,\ 12$ \\
$6$ & $6,\ 7,\ 8,\ 9,\ 10,\ 12,\ 14,\ 18,\ 24,\ 30$ \\
$7$ & $7,\ 8,\ 9,\ 12,\ 14,\ 15,\ 18,\ 20,\ 24,\ 30$ \\
$8$ & $8,\ 9,\ 10,\ 12,\ 14,\ 15,\ 16,\ 18,\ 20,\ 24,\ 30,\ 36,\ 42,\ 60$ \\
$9$ & $9,\ 10,\ 12,\ 16,\ 18,\ 20,\ 21,\ 24,\ 28,\ 30,\ 36,\ 40,\ 42,\ 60$ \\
$10$ & $10,\ 11,\ 12,\ 15,\ 16,\ 18,\ 20,\ 22,\ 24,\ 28,\ 30,\ 36,\ 40,\ 42,\ 48,\ 60$ \\
\bottomrule
\end{tabular}
\end{center}

\medskip\noindent\textbf{Em grau $2$ o cinco não está lá.} É o
Lema~\ref{fis:lem:cristal} --- a restrição cristalográfica --- obtido por um
caminho que não é o dele: aqui não há traço, não há intervalo de comprimento $4$,
não há raiz avaliada. Há uma órbita inteira que fecha ou não fecha, e o $5$
simplesmente não aparece entre os períodos.

\medskip\noindent\textbf{E em grau $4$ o cinco aparece.} Este é o resultado do
varrimento, e ele corrige a leitura fácil da restrição:
\[
\boxed{\;\text{a proibição do cinco não é da REDE --- é do POSTO}\;}
\]
Em posto $2$ não há sítio para uma rotação de ordem cinco; em posto $4$ há, e ela
existe. O Lema~\ref{fis:lem:cristal} é um enunciado sobre matrizes
$2\times2$, e o que ele proíbe é a ordem cinco \emph{ali} --- não em toda a
parte. Subir o grau liberta o período, e o varrimento mostra onde: o $5$ entra ao
grau $4$, o $7$ ao grau $6$, o $11$ ao grau $10$.

\medido{\code{tests/espectro.c} --- e a corrida é inteira do princípio ao fim:
\emph{nenhuma raiz é avaliada, nenhum módulo é comparado, nenhum limiar decide}.
§E1: $\abs{\det C}=1$ nas $\mathbf{59\,046}$ cadeias. §E2: a dicotomia é
exaustiva, $\mathbf{184}$ fecham e $\mathbf{58\,862}$ escapam, e as duas classes
são não vazias em todo grau. §E3: os períodos da tabela, com o grau $2$ a dar
$\{2,3,4,6\}$ e \textbf{zero} ocorrências do $5$. §E4: $\ell=0$ em $184$ de
$184$ --- a companheira é invertível, logo não há o que gastar antes de ciclar.
§E5, o \textbf{controlo do teto}: varrido por cinco décadas, de $10^{3}$ a
$10^{7}$, a contagem das que fecham dá $\mathbf{184}$ em todas --- \emph{o teto
não decide nada}. §E6, o \textbf{controlo por outro caminho}: a reciprocidade
$a_i=\pm a_{N-i}$ é um teste sobre coeficientes inteiros que não vê órbita
nenhuma, e concorda --- das que fecham, $\mathbf{184}$ são recíprocas e
$\mathbf{0}$ não são. E a concordância não é vazia: $\mathbf{540}$ cadeias
recíprocas \emph{escapam}, pelo que ser recíproca não basta.}

\medskip\noindent\textbf{E é aqui que a taxonomia aparece sem ter sido
escolhida.} Não se perguntou qual cadeia é o elétron. Perguntou-se que objetos a
álgebra permite, e a resposta veio repartida: um lado que fecha e não tem massa
nenhuma, outro que escapa e tem; e dentro do que fecha, uma lista discreta de
períodos que cresce com o grau segundo uma regra que o próprio varrimento exibe.
\emph{Nada disto foi quantizado depois: a órbita é inteira desde o primeiro
passo, e o discreto é o que ela é.}

\section{O defeito de massa é a folga}\label{fis:defeito}

\noindent A massa de um núcleo é \textbf{menor} que a soma das massas dos seus
constituintes, e a diferença é a \emph{energia de ligação}. Na linguagem deste
documento isso não é uma segunda lei: é o Cor.~\ref{fis:cor:folga}.

\begin{teorema}[o defeito é o que a colagem identificou]\label{fis:thm:defeito}
Seja $\pi\colon I\to X$ a realização que leva os \emph{constituintes} nas
\emph{células ocupadas} do composto. Então
\[
\boxed{\;\underbrace{\abs{I}}_{\text{a soma das partes}}
\;-\;\underbrace{\abs{X}}_{\text{o composto}}
\;=\;\sum_{x}\bigl(G(x)-1\bigr)\;}
\]
e a diferença é exatamente a \textbf{dobra}: o número de constituintes que
vieram parar à mesma célula.
\end{teorema}

\begin{proof}
É o Cor.~\ref{fis:cor:folga}, que por sua vez é o Lema~\ref{fis:lem:conserva}: as
fibras particionam $I$, logo $\sum_x G(x)=\abs{I}$, e subtrair o suporte dá
$\sum_x(G-1)=\abs{I}-\abs{X}$.
\end{proof}

\medskip\noindent\textbf{Donde a leitura, e ela é literal.} Somar as partes é
contar $I$; pesar o todo é contar $X$; e o \emph{defeito} é a folga entre os
dois. O que a energia de ligação mede é \textbf{quanto se sobrepôs}, e é por a
sobreposição existir que há energia a cobrar ou a devolver.
\[
\boxed{\;\text{ligar é dobrar; a energia de ligação é a dobra, contada.}\;}
\]
E a folga é uma só, vista pelos dois mapas (§\ref{fis:folga}): do lado de quem
\emph{enche} ela aparece como constituintes na mesma célula --- o núcleo ligado;
do lado de quem \emph{conta} ela aparece como lugar por preencher --- a energia
que falta para o desmontar. \emph{São o mesmo número.}

\medskip\noindent\textbf{E a curva da ligação tem um fundo.} A energia de ligação
por constituinte sobe do mais leve até ao ferro e desce depois: o pico está em
$\mathrm{Fe}\text{-}56/\mathrm{Ni}\text{-}62$, com $\approx8{,}79$ MeV por
núcleon. \emph{Subir a curva liberta energia}, e é isso que o teorema seguinte
organiza.

\section{As duas faces do núcleo, e o ponto fixo que partilham}\label{fis:fusaofissao}

\noindent A §\ref{fis:corte} mostrou que o operador tem \emph{duas} faces, que
elas batem alternadamente, e que o processo \textbf{termina} no ponto fixo que as
duas passam a partilhar. O núcleo é essa figura com nomes próprios.

\begin{center}\small
\begin{tabular}{@{}lll@{}}
\toprule
 & \textbf{fusão} & \textbf{fissão} \\
\midrule
o que faz & \emph{junta} dois leves num pesado & \emph{quebra} um pesado em dois \\
de onde vem & dos leves, \textbf{a subir} & dos pesados, \textbf{a descer} \\
o que paga & vencer a repulsão: pressão e calor & absorver um neutrão \\
para onde vai & o vale do ferro & o vale do ferro \\
\bottomrule
\end{tabular}
\end{center}

\begin{teorema}[o vale do ferro é o ponto fixo comum]\label{fis:thm:ferro}
As duas operações do núcleo são as duas faces do Teor.~\ref{fis:thm:corte}: uma
sobe a curva da ligação a partir dos leves, a outra desce-a a partir dos pesados,
e \textbf{ambas libertam energia}. Batendo alternadamente, os intervalos
\emph{encaixam} --- o mais leve sobe, o mais pesado desce --- e o processo
termina num ponto que as duas partilham. Esse ponto é o \textbf{atrator} da
órbita nuclear.
\end{teorema}

\begin{proof}
É o Teor.~\ref{fis:thm:corte} com a curva da ligação por régua: a face que sobe
não passa o fundo (o produto da fusão nunca fica mais ligado que o fundo), a que
desce também não, e as duas larguras encaixam por
Cor.~\ref{fis:cor:medias}. A descida é bem-fundada porque o número de núcleons é
um inteiro e a grade tem grão --- logo o processo \textbf{chega}, e onde chega as
duas faces têm o mesmo ponto fixo, que é o Teor.~\ref{fis:thm:corte}(5).
\end{proof}

\medskip\noindent\textbf{E é por isso que ali a energia acaba.} O ponto fixo
comum é onde $\delta=0$: as duas faces deixam de se mover, e não há mais nada a
libertar. \emph{Uma face endereça, duas produzem} --- e aqui as duas produzem o
fundo do poço. Acima dele a fissão desce; abaixo dele a fusão sobe; nele, nem
uma nem outra.
\[
\boxed{\;\text{o vale do ferro é o \textbf{corte}: o ponto fixo que a fusão e a
fissão passam a partilhar}\;}
\]
E as estrelas percorrem exatamente esse encaixe: fundem o mais leve em hélio e,
nas massivas, sobem a curva \emph{até ao ferro} --- e param lá, porque é o ponto
fixo. \emph{A gravidade é o confinamento; a paragem é o corte.}

\section{O multiplicador é o espectro: subcrítico, crítico, supercrítico}\label{fis:kcritico}

\noindent Numa reação em cadeia, cada quebra devolve dois a três neutrões, e cada
um pode quebrar outro núcleo. O que decide o destino é um número só, e ele tem de
ser \emph{definido} antes de ser classificado.

\begin{definicao}[o operador de reprodução, e o seu raio]\label{fis:def:reproducao}
Seja $R$ o \textbf{operador de reprodução}: a aplicação linear que leva a
composição de uma geração na da seguinte, com $R_{ij}$ o número esperado de
quebras do tipo $i$ produzidas por uma quebra do tipo $j$. As entradas são
não-negativas. O \textbf{multiplicador} é o seu \textbf{raio espectral}
\[
k \;=\; \varrho(R)\;=\;\max_i\abs{\lambda_i(R)} ,
\]
e não o determinante nem o traço. \emph{É esta a quantidade que decide}, e é ela
que a tricotomia abaixo classifica.
\end{definicao}

\begin{teorema}[os três regimes da cadeia]\label{fis:thm:kappa}
Com $R$ da Def.~\ref{fis:def:reproducao} e $v_n$ a composição da geração $n$:
\[
\begin{array}{llll}
k<1 & \textbf{subcrítico} & v_n\to0\ \text{geometricamente} & \text{o decaimento}\\[2pt]
k=1 & \textbf{crítico} & v_n\ \text{nem cresce nem morre} & \text{o reator}\\[2pt]
k>1 & \textbf{supercrítico} & v_n\ \text{cresce geometricamente} & \text{a explosão}
\end{array}
\]
e a tricotomia é exaustiva porque $k$ é um real não-negativo.
\end{teorema}

\begin{proof}
$v_n=R^{n}v_0$, e o crescimento de $\lVert R^{n}\rVert$ é governado por
$\varrho(R)$:
$\varrho<1$ dá contração geométrica, $\varrho>1$ dá expansão geométrica, e
$\varrho=1$ é a fronteira, onde o crescimento é sub-geométrico --- polinomial no
tamanho do maior bloco de Jordan associado, e limitado quando esse bloco é
trivial. As três alternativas esgotam $k\ge0$.
\end{proof}

\medskip\noindent\textbf{E onde é que isto encontra a caixa --- e onde NÃO
encontra.} Em \emph{posto $2$} a correspondência é limpa e é a da
§\ref{fis:leidisc}: com $\abs{\det}=1$ o par de raízes é ou conjugado no círculo
($\Delta<0$, $k=1$), ou duplo ($\Delta=0$, $k=1$ com crescimento linear), ou real
com uma fora ($\Delta>0$, $k>1$) --- e é exactamente aí que a órbita de $(1,0)$
dá $1,1,1,1,1,1$ / $1,2,3,4,5,6$ / $1,1,2,3,5,8$.

\medskip\noindent\emph{Em posto $N$ arbitrário a identificação não se estende
por decreto.} O discriminante é um invariante de matrizes $2\times2$, e $R$ não é
uma companheira com $\abs{\det}=1$ --- é um operador de entradas não-negativas,
cujo raio espectral não se lê de $\Delta$. O que atravessa os dois casos é a
\textbf{dicotomia da órbita} do Teor.~\ref{fis:thm:orbita} --- fica ou escapa ---
e é ela, e não o discriminante, o enunciado geral.
\[
\boxed{\;\text{em posto }2\text{ o regime lê-se em }\Delta;\ \text{em posto }N
\text{ lê-se no \textbf{raio espectral}}\;}
\]

\medskip\noindent\textbf{E cada regime tem a sua máquina.}

\medskip\noindent\emph{Subcrítico.} O decaimento radioativo: um núcleo instável
emite e vira outro, e a população cai por metades,
$N(t)=N_0\cdot2^{-t/T}$. É a órbita que \textbf{contrai} --- e a meia-vida $T$ é
o grão da régua, o passo em que a largura se halva. É literalmente a cláusula~(2)
do Teor.~\ref{fis:thm:corte}: $\delta'\le\delta/2$.

\medskip\noindent\emph{Crítico.} O reator: barras de controlo absorvem neutrões e
o moderador desacelera-os, para segurar $k$ \emph{exatamente} em $1$. E é a
posição mais difícil de manter, porque é \textbf{a borda}: pelo
Cor.~\ref{fis:cor:borda} a borda é o conjunto dos espelhos cujo $\Delta$ é
quadrado, e ele é \emph{raro}; a classe parabólica está inteiramente dentro dela e
é a única que está. \emph{Um reator é uma máquina que vive na borda por
construção} --- na fronteira onde a matéria com massa toca o modo sem massa, que
é o Teor.~\ref{fis:thm:massa}.

\medskip\noindent\emph{Supercrítico.} Montada de repente uma massa com
$k\gg1$, a cadeia cresce geometricamente antes de se dispersar. É a órbita
\textbf{hiperbólica}, e os números do crescimento são os da família metálica ---
o mesmo $1,1,2,3,5,8$ que a §\ref{fis:leidisc} mede para o regime que expande.

\medskip\noindent\textbf{E a tricotomia não se transporta.} Pelo
Cor.~\ref{fis:cor:serie}, $\operatorname{sgn}\Delta$ \emph{não é função} das
classes dos factores: compor dois sistemas estáveis pode dar um instável. É a
razão estrutural de a criticidade ter de ser mantida \emph{elemento a elemento} e
não uma vez só --- \emph{a estabilidade não é uma propriedade que se transporte ao
longo de uma série}.

\section{A cadeia de decaimento é uma órbita: a cauda e o período}\label{fis:cadeia}

\noindent Um núcleo instável decai numa sequência de nuclídeos até parar num
estável. Isso é, palavra por palavra, uma trajetória determinista no sentido da
Def.~\ref{fis:def:determinista} --- e o campo lê-a.

\begin{teorema}[o campo lê a cadeia sem a seguir]\label{fis:thm:cadeia}
Seja $\pi$ a sucessão dos nuclídeos de uma cadeia de decaimento, com $\pi(t+1)$
função de $\pi(t)$. Se a cadeia termina num nuclídeo \textbf{estável}, então a
trajetória é periódica de período $p=1$ --- o estável é um ponto fixo --- e
\[
\abs{\{x:G(x)>1\}}=1,\qquad \abs{\{x:G(x)=1\}}=\ell ,
\]
com $\ell$ o número de passos da cadeia. \emph{A cauda conta a cadeia; o período
diz que ela parou.}
\end{teorema}

\begin{proof}
Um nuclídeo estável cumpre $f(x)=x$, logo a trajetória é constante a partir dele
e $p=1$. Os nuclídeos anteriores são distintos entre si --- uma cadeia de
decaimento não volta a um nuclídeo já visitado, porque cada passo baixa a massa
ou a carga --- pelo que o primeiro índice repetido é o do estável, isto é
$\ell$ é o comprimento da cadeia. Aplica-se então o Teor.~\ref{fis:thm:periodo},
cuja hipótese $N\ge\ell+2p-1=\ell+1$ vale assim que a trajetória atinge o estável
uma segunda vez.
\end{proof}

\medskip\noindent É o mesmo enunciado que a §\ref{fis:pi} usa para as somas
parciais: \textbf{período $1$, e a cauda é o custo}. Numa série o custo são as
somas distintas; numa cadeia de decaimento são os nuclídeos atravessados --- a
cadeia $\mathrm{U}\text{-}238\to\dots\to\mathrm{Pb}\text{-}206$ tem
$\ell=\mathbf{14}$ passos, e $\{G>1\}$ tem um elemento só, que é o chumbo.

\medskip\noindent\textbf{E a estabilidade é um ponto fixo, não uma ausência de
lei.} Um nuclídeo estável não é um que «não decai por acaso»: é um que a órbita
fixa. E o campo distingue-o sem seguir a trajetória --- pela
cláusula~(3) do Teor.~\ref{fis:thm:mult}, basta a marca no espaço:
\[
\boxed{\;\text{a estabilidade é }G>1\text{ numa célula só}\;}
\]

\medskip\noindent\textbf{E aqui é preciso separar duas coisas que a palavra
«decaimento» junta, porque só uma delas é o que este teorema trata.}
\[
\boxed{\;\text{a \textbf{cadeia} de nuclídeos é estrutural}
\qquad\neq\qquad
\text{o \textbf{instante} do evento}\;}
\]
O Teor.~\ref{fis:thm:cadeia} é sobre a \emph{sequência}: que nuclídeo se segue a
qual, quantos passos até ao estável, e que o estável é um ponto fixo. Essa
sequência é determinada --- $\mathrm{U}\text{-}238$ vai a
$\mathrm{Pb}\text{-}206$ pelos mesmos catorze passos, sempre. \emph{Quando} um
núcleo individual dá o seu passo é outra pergunta, e a resposta dela não está
neste teorema: um único núcleo não tem hora marcada.

\medskip\noindent\textbf{E a pergunta que isso abre é séria, e fica escrita
porque é interessante e não porque incomoda:} a construção dá a órbita, o
período e a cauda --- \emph{de onde sai a distribuição no tempo?} O que a
§\ref{fis:kcritico} já mostra é a forma da \emph{população}, $N(t)=N_0\cdot
2^{-t/T}$, que é a órbita a contrair com a meia-vida por grão; o que ela não
mostra é o mecanismo que reparte os instantes individuais dentro dessa forma.
\emph{É o primeiro sítio onde a construção pede mais do que já deu}, e vale mais
tê-lo dito do que ter deixado a palavra «determinista» a cobrir os dois sentidos.

\section{A carga é uma contagem}\label{fis:carga}

\noindent A última peça é a que mais se parece com o Lema~\ref{fis:lem:conserva}:
a carga não é um número posto à mão --- é um \textbf{número de ocupação}, e sai
quantizada porque uma contagem é inteira.

\begin{teorema}[a carga sai da contagem, e por isso é quantizada]\label{fis:thm:carga}
Seja $N$ o número de ocupação de um estado --- a contagem das escadas ocupadas,
que é exatamente uma fibra no sentido da Def.~\ref{fis:def:objeto}. Então $N$ toma
\textbf{quatro} valores, $N\in\{0,1,2,3\}$, sobre \textbf{oito} estados, com
multiplicidades $1,3,3,1$ --- e
\[
\boxed{\;Q=\tfrac{N}{3}\;}
\]
é quantizada, porque $N$ é uma \textbf{contagem} e não uma medida. Num ideal, os
valores são $Q\in\{0,\tfrac13,\tfrac23,1\}$; com o \emph{conjugado} juntam-se os
simétricos, e a lista fecha em
$Q\in\{0,\pm\tfrac13,\pm\tfrac23,\pm1\}$ --- dezasseis estados, que são uma
geração.
\end{teorema}

\begin{proof}
$N$ é o cardinal de uma fibra, logo inteiro e não-negativo por
Def.~\ref{fis:def:objeto}; e $Q=N/3$ herda a discretude. A contagem por nível dá
as multiplicidades $1,3,3,1$ para $N=0,1,2,3$ --- que é a linha do triângulo da
§\ref{fis:hipercubo} com $n=3$ --- e a soma $1+3+3+1=8$ é o
$\sum G=2^{n}=\abs{I}$ do Lema~\ref{fis:lem:conserva}. O conjugado duplica a
contagem sem a alterar, e $8+8=16$.
\end{proof}

\medskip\noindent\textbf{E há DUAS involuções, não uma --- e convém não as
colar.} A linha $1,3,3,1$ é simétrica porque $\binom{n}{k}=\binom{n}{n-k}$, e essa
simetria é uma involução do \emph{número de ocupação}; a que troca partícula por
antipartícula é outra, e é a que passa ao \emph{conjugado}. Contadas:
\[
\begin{array}{lll}
\textbf{ocupação} & N\mapsto3-N & Q\mapsto1-Q,\ \text{ponto fixo }Q=\tfrac12\\[2pt]
\textbf{conjugação} & \text{ideal}\mapsto\text{conjugado} & Q\mapsto-Q,\ \text{ponto fixo }Q=0
\end{array}
\]
\emph{A primeira troca ocupado por vazio; a segunda troca partícula por
antipartícula.} Confundi-las dá $Q\mapsto-Q$ onde só há $Q\mapsto1-Q$, e apaga
justamente o que cada uma tem de próprio: o seu ponto fixo.

\medskip\noindent\textbf{E os dois pontos fixos dizem coisas diferentes, ambas
verdadeiras.} O da ocupação é $N=\tfrac32$, isto é $Q=\tfrac12$: \emph{nenhuma
contagem o alcança}, e é por isso que não há estado de carga $\tfrac12$ --- é o
\emph{meio} do Teor.~\ref{fis:thm:duas}(2) a cair fora da grade, exatamente como
$g_{\ast}$ cai fora dela na Def.~\ref{fis:def:piso}. O da conjugação é $Q=0$, e
esse \textbf{é} alcançado: é o estado neutro, o \emph{vácuo} do ideal --- o único
que é a sua própria antipartícula.
\[
\boxed{\;\text{o neutro não sobra da conta: é o \textbf{ponto fixo} da conjugação}\;}
\]

\medskip\noindent\textbf{E compor as duas dá a terceira.} $Q\mapsto1-Q$ seguida
de $Q\mapsto-Q$ é $Q\mapsto Q-1$: uma \emph{translação}, não uma involução ---
e é o mesmo padrão do Teor.~\ref{fis:thm:pandora}(2), onde a travessia entre
espelhos de $\Delta$ igual também é uma translação do gerador e não uma
conjugação. \emph{Duas reflexões dão um deslocamento}, que é o
Cor.~\ref{fis:cor:serie} lido aqui: a orientação compõe-se, e o produto de duas
que refletem não reflete.

\medskip\noindent Em uma linha, e é a frase que esta parte inteira realiza:
\[
\boxed{\;\text{a carga é a \textbf{fibra}; a massa é o \textbf{logaritmo do
raio}; a ligação é a \textbf{folga}.}\;}
\]

\section{O problema não é falta de dados: é a leitura}\label{fis:leitura-dados}

\noindent Esta parte não termina a pedir experiências. Termina a dizer o que
fazer com as que já existem.

\medskip\noindent A física não sofre de escassez de observação. Há espectros,
massas, cargas, tempos de vida, secções de choque, cadeias de decaimento,
estados ligados, níveis atómicos, linhas espectrais, oscilações, abundâncias
nucleares --- um arquivo que já não cabe na intuição de uma pessoa. \emph{O que
falta não é mais um número.}

\medskip\noindent\textbf{A leitura convencional começa pelos objetos} ---
elétron, quark, núcleo, fotão --- mede as suas propriedades, e depois procura
relações entre as medidas. Aqui o movimento é o inverso:
\[
\boxed{\;\text{não se parte da partícula para explicar o dado;}\qquad
\text{parte-se da estrutura que organiza os dados}\;}
\]

\medskip\noindent E a troca é peça a peça, cada uma com o enunciado que já a
sustenta nesta parte: a massa deixa de ser etiqueta e passa a ser
\emph{coordenada espectral} (Teor.~\ref{fis:thm:massa}); a carga deixa de ser
número primitivo e passa a ser \emph{ocupação} (Teor.~\ref{fis:thm:carga}); a
estabilidade deixa de ser propriedade e passa a ser \emph{ponto fixo}
(Teor.~\ref{fis:thm:cadeia}); a ligação deixa de ser força acrescentada e passa a
ser \emph{multiplicidade} (Teor.~\ref{fis:thm:defeito}).

\medskip\noindent\emph{Os dados não mudaram. Mudou a pergunta feita aos dados.}
Ela deixa de ser «que partícula explica este resultado?» e passa a ser
\[
\boxed{\;\text{que estrutura produz simultaneamente estas famílias de resultados?}\;}
\]

\medskip\noindent\textbf{E é por já haver tanto arquivo que a pergunta se pode
atacar agora}, sem esperar máquina nova. As leituras cruzadas que a construção
pede são estas, e cada uma é uma reorganização do que já está medido:

\begin{center}\small
\begin{tabular}{@{}lll@{}}
\toprule
o arquivo & a leitura & contra o quê \\
\midrule
massas conhecidas & procurar $(\rho,\theta)$ & Teor.~\ref{fis:thm:massa} \\
estados de carga & reorganizar pela fibra $N$, não pelo nome & Teor.~\ref{fis:thm:carga} \\
tabelas de decaimento & como grafos: pontos fixos e órbitas & Teor.~\ref{fis:thm:cadeia} \\
energias de ligação & sobre a estrutura de folgas $\abs{I}-\abs{X}$ & Teor.~\ref{fis:thm:defeito} \\
linhas espectrais & os períodos e ângulos que a álgebra permite & §\ref{fis:mapa} \\
\bottomrule
\end{tabular}
\end{center}

\medskip\noindent E a disciplina do varrimento é a que a §\ref{fis:mapa} já
seguiu: \emph{não se ajusta a cadeia para reproduzir uma partícula} --- procura-se
a partícula entre as cadeias que a álgebra já permite. Se a correspondência
falhar, não há parâmetro livre onde a esconder: falha uma afirmação estrutural. E
se sobreviver, não terá sido por se escolher a resposta depois de olhar o
resultado, porque a lista de candidatos foi produzida \textbf{antes} da
comparação --- é o que aquelas cento e oitenta e quatro cadeias que fecham são.

\medskip\noindent Em uma linha:
\[
\boxed{\;\text{a diversidade dos fenómenos pode ser grande;
a arquitetura que os produz pode ser \textbf{pequena}}\;}
\]

\colunas{Prop.~\ref{fis:prop:neutro}; Teor.~\ref{fis:thm:massa},~\ref{fis:thm:orbita},~\ref{fis:thm:defeito},~\ref{fis:thm:ferro},~\ref{fis:thm:kappa},~\ref{fis:thm:cadeia},~\ref{fis:thm:carga}}
        {o átomo de grau $N$; a borda onde a matéria vira luz; o vale do ferro como corte; os três regimes de $k$}
        {\code{cristal/cristal\_fisica.tex} (energia nuclear, partículas minerais); \code{tests/neuronio.c}; \code{tests/pgwire.c} §W116, §W219--§W220}

% ═══════════════════════════════════════════════════════════════════════════
% ═══════════════════════════════════════════════════════════════════════════
\part{Termodinâmica}
% ═══════════════════════════════════════════════════════════════════════════

\section*{A tese desta parte}\label{fis:termo-tese}

\noindent A parte seguinte lê a expansão. Esta lê a \textbf{entropia}, e as duas
são o mesmo par --- o que se desfaz e o que se abre. Mas há uma coisa a dizer
primeiro, e ela vem antes de qualquer identificação: \emph{este é o terreno onde
é mais fácil chamar lei ao que é escolha de quem escreve}. A §\ref{fis:carnot}
fecha uma conjectura desta casa \textbf{como negativa}, e fecha-a por essa razão
exacta.

\begin{center}\small
\begin{tabular}{@{}lll@{}}
\toprule
& \textbf{a estrutura} & \textbf{o que a obriga} \\
\midrule
o rendimento máximo & o $\oint$ fechado, não a engenharia & Teor.~\ref{fis:thm:carnot} \\
o que custa & apagar --- e apagar é a \emph{dobra} & Teor.~\ref{fis:thm:landauer} \\
o piso do reversível & \textbf{zero exacto}, não pequeno & Teor.~\ref{fis:thm:landauer} \\
a entropia dual & o produto vale $1$; os expoentes somam $0$ & Teor.~\ref{fis:thm:entropiadual} \\
a segunda lei & vale de \textbf{um} lado --- o par não se move & Cor.~\ref{fis:cor:segundalei} \\
\bottomrule
\end{tabular}
\end{center}

\section{O rendimento é o $\oint$ fechado}\label{fis:carnot}

\begin{teorema}[o rendimento máximo, derivado de duas integrais]\label{fis:thm:carnot}
Num ciclo:
\[
\oint dU=0 \;\Longrightarrow\; W=Q_q-Q_f ,
\qquad
\oint \frac{dQ}{T}=0 \;\Longrightarrow\; \frac{Q_f}{Q_q}=\frac{T_f}{T_q} ,
\]
donde
\[
\boxed{\;\eta=\frac{W}{Q_q}=1-\frac{Q_f}{Q_q}=1-\frac{T_f}{T_q}\;}
\]
\end{teorema}

\begin{proof}
A primeira integral é o ciclo voltar ao estado: a energia interna é função de
estado, logo a sua variação num ciclo é nula, e o que sobra é $W=Q_q-Q_f$. A
segunda é a reversibilidade: $dQ/T$ é diferencial exacta num ciclo reversível,
logo a sua circulação anula-se, e daí $Q_q/T_q=Q_f/T_f$. Substituindo a razão dos
calores pela razão das temperaturas em $\eta=1-Q_f/Q_q$ vem o enunciado.
\end{proof}

\medskip\noindent\textbf{É o $\oint$ que faz o trabalho.} A primeira integral dá
$W$; a segunda converte a razão dos \emph{calores} na razão das
\emph{temperaturas}. Sem a segunda, $\eta$ fica em $Q$ e não em $T$. E abrir o
$\oint$ --- rejeitar calor a mais --- derruba $\eta$ abaixo do máximo,
\emph{sempre}:
\[
\boxed{\;\text{o tecto não é limite de engenharia: é o }\oint\text{ fechado}\;}
\]

\subsection*{E uma conjectura desta casa, fechada como negativa}

\noindent Houve aqui uma identificação tentadora: a razão focal de uma elipse,
$(1-e)/(1+e)$, como candidata a $T_f/T_q$. Ela \emph{funciona} --- dá um valor
válido para todo $e\in(0,1)$. E é por isso que não serve.

\medskip\noindent\textbf{O teste que a derruba é uma pergunta só:} \emph{que
parâmetros do teorema fui eu que escolhi?} Qualquer bijeção de $(0,1)$ em $(0,1)$
daria um $T_f/T_q$ igualmente válido. Sem \textbf{exibir} o ciclo cuja trajetória
\emph{é} a elipse do corte, aquilo não é um teorema --- é uma
\textbf{parametrização}, e uma parametrização não prova nada porque não pode
falhar.
\[
\boxed{\;\text{se qualquer bijeção servisse, a que se escolheu não é resultado}\;}
\]
\emph{Fica fechado como negativo}, e fica escrito, porque um negativo medido vale
mais do que uma coincidência bonita. O que \textbf{sobrevive} é o que não dependia
disto: o ciclo fecha se e só se $\Delta<0$, os dois focos são dois, e a janela é
aberta nos extremos --- \emph{isso é geometria da secção}
(Teor.~\ref{fis:thm:leidisc}), e vale. A temperatura é que não entrou.

\medido{\code{tests/carnot.c} --- $9$ blocos, $\mathbf{0}$ falhas, \emph{resíduo
$0$, exacto em racionais}. §C1: o ciclo fecha, $\oint dU=0$. §C2: o reversível dá
$Q_f/Q_q=T_f/T_q$. §C3: $\eta=1-T_f/T_q$ igual a $W/Q_q$ em toda a varredura. E o
veredito \textbf{negativo} sobre a razão focal está no próprio medidor, não numa
nota de rodapé.}

\section{O que custa é apagar}\label{fis:landauer}

\noindent A régua não é a memória reservada --- é o \textbf{bit apagado}. E as
duas separam-se de forma que não deixa dúvida: um vector grande que nunca é
escrito não dissipa nada, e um único bit reescrito mil milhões de vezes paga mil
milhões de vezes.

\begin{teorema}[o piso do involutivo é zero exacto]\label{fis:thm:landauer}
Apagar um bit tem um custo mínimo, e ele é positivo. Logo:
\begin{enumerate}\itemsep3pt
\item[(1)] uma operação \textbf{destrutiva} --- que identifica estados --- paga
por cada bit que perde;
\item[(2)] uma operação \textbf{involutiva} --- que aplicada duas vezes devolve
--- não perde bit nenhum, e o seu piso é
\[
\boxed{\;\text{zero \textbf{exacto}, e não um número pequeno}\;}
\]
\end{enumerate}
\end{teorema}

\begin{proof}
(1) O custo conta-se em bits perdidos, e perder um bit é ter dois estados de
entrada com a mesma saída --- é \emph{exactamente} $G>1$ na
Def.~\ref{fis:def:objeto}. Uma operação destrutiva tem fibra não trivial, logo
tem bits a perder.

(2) Uma involução é bijetiva: $I^{2}=\mathrm{id}$ obriga $I$ a ser injetiva, logo
$G\equiv1$ pela cláusula~(2) do Teor.~\ref{fis:thm:mult}. Sem fibra não trivial
não há estado identificado, logo não há bit apagado, e a soma dos custos é uma
soma vazia --- \emph{zero, e não um infinitésimo}.
\end{proof}

\medskip\noindent\textbf{E é a dobra outra vez, com unidade física.} O que a
§\ref{fis:folga} conta como $\sum_x(G(x)-1)$ --- os índices que vieram parar à
mesma célula --- é aqui o que se paga. \emph{A folga tem preço}, e o preço é o
mínimo termodinâmico vezes a folga.
\[
\boxed{\;\text{dissipar é dobrar}\;}
\]
Donde a leitura que a construção já vinha a fazer sem lhe chamar isto: as
operações que este documento constrói --- o gato e o esquilo, a subida e a
descida da torre, $\zeta$ e $\mu$ --- \emph{invertem}, e por isso o seu piso é
zero. \emph{Não é economia de implementação: é o Teor.~\ref{fis:thm:gato}(1),
$\oslash\circ\otimes=\mathrm{id}$, lido como termodinâmica.}

\medido{\code{tests/dissipa.c} --- $4$ blocos, $\mathbf{0}$ falhas. Contam-se os
bits apagados dos \emph{dois} lados, na mesma tarefa e com a mesma entrada: a
destrutiva paga $\mathbf{376\,304\,828\,416}$ rJ ($10^{-27}$ J), a involutiva
paga $\mathbf{0}$ --- e os dois lados têm de ser \emph{visíveis} para a frase
valer, senão a distinção não aparece em lado nenhum. E a involutiva é mesmo
involutiva: repetida, devolve o original em $\mathbf{4096}$ de $4096$ posições.
\emph{Nota:} este bloco imprimia um custo por segundo que transbordava num alvo
de inteiro largo a $32$ bits --- saía \emph{negativo}, e por não ser asserido
nenhum bloco o apanhava. Corrigido para $64$ bits explícitos, dá
$3{,}08\cdot10^{15}$ rJ/s. \emph{O tecto diz-se no tipo}, e um número publicado
errado é pior do que um bloco vermelho.}

\section{A entropia dual, e a segunda lei lida como ela é}\label{fis:entropiadual}

\noindent A entropia de um horizonte é a sua \textbf{área}. E no universo dual o
raio é o inverso --- não há nada a inventar: aplica-se a mesma lei ao raio dual.

\begin{teorema}[o par fecha na unidade]\label{fis:thm:entropiadual}
Com $S=r^{\,d-1}$ a área em codimensão $1$, e o dual a inverter o raio:
\[
S_{\text{negro}}=r^{\,d-1}\ \ (\text{cresce}),
\qquad
S_{\text{branco}}=r^{-(d-1)}\ \ (\text{decresce}),
\]
e então
\[
\boxed{\;S_{\text{negro}}\cdot S_{\text{branco}}=1
\qquad\text{e}\qquad
\log S_{\text{negro}}+\log S_{\text{branco}}=0\;}
\]
--- o \textbf{produto} vale a unidade, e a \textbf{soma dos expoentes} é nula.
\end{teorema}

\begin{proof}
$r^{\,d-1}\cdot r^{-(d-1)}=r^{0}=1$, e o logaritmo de $1$ é $0$. As duas
igualdades são a mesma escrita nas duas faces: \emph{o que é produto de um lado é
soma do outro}, e a exponencial é o dicionário entre elas
(§\ref{fis:transf}). E não é acidente que a entropia caiba na segunda: ela
\emph{já é um logaritmo} --- conta caminhos, e contar caminhos é somar expoentes.
\end{proof}

\medskip\noindent\textbf{E isto é a neutralidade, outra vez.} A
Prop.~\ref{fis:prop:neutro} diz que num átomo $\sum_i\log\abs{\rho_i}=0$ porque o
determinante é unidade. Aqui a soma dos dois expoentes é nula pela mesma razão
estrutural: \emph{o par conserva, e nenhum dos lados conserva sozinho}.

\begin{corolario}[a segunda lei vale de um lado]\label{fis:cor:segundalei}
A entropia do lado negro cresce em todos os passos; a do branco decresce em
todos; e a \textbf{soma não se move em nenhum}. Logo
\[
\boxed{\;\text{«a entropia cresce» é verdade, e é lei de \textbf{metade} do par}\;}
\]
e a seta do tempo não é propriedade do tempo: \emph{é uma polaridade lida sem a
outra}.
\end{corolario}

\begin{proof}
$r\mapsto r^{\,d-1}$ é crescente e $r\mapsto r^{-(d-1)}$ é decrescente em $r>0$,
para $d>1$. E a soma dos logaritmos é identicamente nula pelo
Teor.~\ref{fis:thm:entropiadual}, logo não varia com $r$: o que cresce de um lado
é exactamente o que decresce do outro.
\end{proof}

\medskip\noindent\textbf{E o ponto fixo é a borda.} Há um só raio onde as duas
entropias coincidem --- $r=1$ --- e aí ambas valem a unidade, e no logaritmo,
zero. \emph{É aí que a seta não aponta para lado nenhum}, porque nem cresce nem
decresce. É o círculo unitário do Teor.~\ref{fis:thm:massa}, o $\delta=0$ do
Teor.~\ref{fis:thm:corte}, e o $\Delta=0$ da §\ref{fis:leidisc} --- \emph{a mesma
borda, pelo quarto caminho diferente}.

\medido{\code{tests/entropia\_dual.c} --- $6$ blocos, $\mathbf{0}$ falhas. §X5:
há \textbf{um} raio em $40$ onde as duas entropias coincidem, e é $r=1$. §X6, o
\textbf{controlo}: o par fecha com o expoente da \emph{área} com $\mathbf{0}$
falhas; trocando a codimensão pelo \emph{volume} falha em $\mathbf{9}$; e com um
dual \emph{falso} --- o inverso trocado por um múltiplo dele --- falha em
$\mathbf{10}$. \emph{Logo o zero não veio de a álgebra ser generosa}: veio de a
entropia ser a fronteira e de os dois universos serem duais.}

\section{E porque esta parte vem antes da seguinte}\label{fis:termo-cosmo}

\noindent As duas leituras são o mesmo par, e é por isso que estão lado a lado:
\[
\boxed{\;\textbf{entropia} \text{ é o que se desfaz};\qquad
\textbf{expansão} \text{ é o que se abre}\;}
\]
Uma conta estados; a outra cria lugar. \emph{O que se solta de um lado é o espaço
que nasce do outro.}

\medskip\noindent E a forma é a mesma nas duas partes, o que se vê pondo-as em
paralelo:
\begin{center}\small
\begin{tabular}{@{}lll@{}}
\toprule
& \textbf{termodinâmica} & \textbf{cosmologia} \\
\midrule
o invariante & $S_{\text{n}}\!\cdot\!S_{\text{b}}=1$ & $r\,a^{3(1+w)}=C$ \\
como se lê & soma de logaritmos $=0$ & entra aditiva, sai multiplicativa \\
o ponto fixo & $r=1$: a seta não aponta & $w=-\tfrac13$: a aceleração anula \\
a lei de um lado & «a entropia cresce» & «o universo expande» \\
a lei do par & a soma não muda & a constante não muda \\
\bottomrule
\end{tabular}
\end{center}

\medskip\noindent\textbf{E aqui é preciso ser exacto sobre o que está e o que não
está demonstrado}, porque é o terreno onde esta casa já se enganou uma vez. O que
está provado é cada coluna \emph{por si}: os teoremas desta parte e os da
seguinte, cada um com a sua prova e o seu controlo. O que a tabela mostra é que
as duas colunas têm a \textbf{mesma forma} --- um invariante multiplicativo que
se lê como soma nula, com o ponto fixo na borda.

\medskip\noindent\emph{Que sejam a mesma lei, e não duas com a mesma forma, é o
que falta.} E o teste que o diria é o mesmo que fechou a §\ref{fis:carnot} como
negativa: \textbf{exibir} a construção onde uma se torna a outra --- não uma
bijeção que as ligue, porque uma bijeção há sempre.
\[
\boxed{\;\text{a mesma forma não é a mesma lei --- e a diferença mede-se
exibindo, não emparelhando}\;}
\]

\colunas{Teor.~\ref{fis:thm:carnot},~\ref{fis:thm:landauer},~\ref{fis:thm:entropiadual}; Cor.~\ref{fis:cor:segundalei}}
        {o $\oint$ fechado; o piso zero do involutivo; o par que fecha na unidade}
        {\code{tests/carnot.c}; \code{tests/dissipa.c}; \code{tests/entropia\_dual.c}}

% ═══════════════════════════════════════════════════════════════════════════
\part{Cosmologia}
% ═══════════════════════════════════════════════════════════════════════════

\section*{A tese desta parte}\label{fis:cosmo-tese}

\noindent A parte anterior leu o pequeno. Esta lê o grande, e com os mesmos
objetos --- \emph{não com outros}. O que muda é a escala a que se aplica a mesma
construção:
\begin{center}\small
\begin{tabular}{@{}lll@{}}
\toprule
& \textbf{a estrutura} & \textbf{o que a obriga} \\
\midrule
a taxa de expansão & $H^{2}=1/D$: a régua, invertida & Teor.~\ref{fis:thm:taxa} \\
os conteúdos & os três valores do trial & Teor.~\ref{fis:thm:conteudos} \\
a diluição & a involução entre o aditivo e o multiplicativo & Teor.~\ref{fis:thm:diluicao} \\
o destino & o mesmo sinal, lido em $-(1+3w)$ & Cor.~\ref{fis:cor:destino} \\
a equação de estado & uma órbita de quatro, duas no ponto fixo & Teor.~\ref{fis:thm:bidual} \\
\bottomrule
\end{tabular}
\end{center}

\medskip\noindent\textbf{E nenhuma constante física fundamental entra.} A parte
corre em \emph{por-unidade} e em inteiros: em por-unidade as constantes valem $1$
por construção, e o que sobra é a estrutura. Onde o sistema \emph{sai} do
por-unidade está dito, e é num sítio só --- a \emph{constante de integração} da
lei de conservação, que não é uma constante fundamental introduzida pela teoria:
é o que carrega a normalização dimensional.

\section{A taxa é a régua}\label{fis:taxa}

\noindent A expansão não é um parâmetro que se ajuste. Ela é o inverso do
discriminante --- o mesmo $\Delta$ que a §\ref{fis:leidisc} usa para classificar
espelhos, aqui lido na família metálica.

\begin{definicao}[a régua cosmológica]\label{fis:def:regua-cosmo}
A escala de expansão \textbf{não é uma quantidade independente}: define-se como a
recíproca do discriminante do operador,
\[
\boxed{\;H^{2}\,D\;=\;1\;}
\]
com $D=\operatorname{tr}^{2}-4\det$ o mesmo discriminante que classifica os
espelhos na §\ref{fis:leidisc}. \emph{É esta a régua}, e é ela --- e não uma
coincidência --- que faz a expansão viver no mesmo invariante que o regime.
\end{definicao}

\begin{teorema}[e daí o valor, que sai do operador]\label{fis:thm:taxa}
Seja $M_m=\left(\begin{smallmatrix}m&1\\1&0\end{smallmatrix}\right)$ o gato de
assinatura $m$. Então
\[
D_m=m^{2}+4
\qquad\text{e portanto}\qquad
\boxed{\;H_m^{2}=\frac{1}{m^{2}+4}\;}
\]
--- lido \textbf{do operador}, e não de uma fórmula escrita ao lado dele. Em
particular $H^{2}$ \textbf{decresce estritamente} quando o grau sobe.
\end{teorema}

\begin{proof}
$\operatorname{tr}M_m=m$ e $\det M_m=-1$, donde
$D_m=m^{2}-4(-1)=m^{2}+4$; o valor de $H_m^{2}$ é a
Def.~\ref{fis:def:regua-cosmo} aplicada a esse $D_m$. O decrescimento é a
monotonia de $m\mapsto m^{2}+4$ em $m\ge0$, e lê-se por produto cruzado, sem
dividir.
\end{proof}

\medskip\noindent\emph{O que está demonstrado e o que é definição, para não se
confundirem:} $H^{2}D=1$ é a \textbf{régua} --- é onde a construção diz o que a
expansão \emph{é}. O que o teorema prova é que, nesta família, esse $D$ vale
$m^{2}+4$, e que a taxa daí resultante decresce com o grau. \emph{A taxa
cosmológica é uma função do mesmo invariante estrutural}, e não uma escala nova.

\medskip\noindent\textbf{E aqui a cosmologia encontra a caixa --- com uma
ressalva que é preciso dizer.} \emph{Não é o sinal de $D$ isoladamente que fixa o
regime: é o par $(\det,\Delta)$.} Na classe normalizada da caixa, com
$\det=+1$, o discriminante separa os três regimes. Na família metálica, porém,
\[
\det M_m=-1
\qquad\Longrightarrow\qquad
\Delta=m^{2}+4>0\ \text{ para todo } m ,
\]
de modo que \emph{esta família pertence inteiramente ao regime hiperbólico} ---
e é por isso que ela expande, para toda a assinatura. É o
Cor.~\ref{fis:cor:optica} lido na cosmologia: \emph{nesta classe de operadores},
o determinante $-1$ força $\Delta=\operatorname{tr}^{2}+4>0$ e, portanto, o
regime hiperbólico.

\medido{\code{tests/cosmologia.c} §C3: $D$ lido da matriz como
$\operatorname{tr}^{2}-4\det$ bate com $m^{2}+4$ em $\mathbf{5}$ de $5$ graus, e
a taxa decresce estritamente em $\mathbf{4}$ dos $4$ passos --- medido por
produto cruzado, \emph{sem dividir}.}

\section{Os três conteúdos são os três valores do trial}\label{fis:conteudos}

\begin{teorema}[e não há quarto]\label{fis:thm:conteudos}
A equação de estado $w=p/r$ toma, nos conteúdos que a construção admite,
exactamente três valores, e eles são os três \textbf{sinais}:
\[
\begin{array}{llcl}
\textbf{radiação} & w=+\tfrac13 & \text{sinal }+ & \text{três eixos, a parte de um}\\[2pt]
\textbf{matéria}  & w=0         & \text{sinal }0 & \text{sem pressão}\\[2pt]
\textbf{vácuo}    & w=-1        & \text{sinal }- & \text{o ponto fixo}
\end{array}
\]
\emph{Não há quarto conteúdo porque não há quarto sinal.}
\end{teorema}

\begin{proof}
A lei de diluição do Teor.~\ref{fis:thm:diluicao} dá $r\propto a^{-3(1+w)}$, e o
expoente é \textbf{inteiro} exactamente nos $w$ para os quais $3(1+w)$ o é.
Exigindo além disso que o conteúdo seja um dos três estados que o trial admite
--- e o trial é $\{-1,0,+1\}$, com $+$ a pressão positiva, $0$ a ausência de
pressão e $-$ o ponto fixo --- os expoentes saem $4$, $3$ e $0$, isto é
$w=+\tfrac13$, $w=0$ e $w=-1$.

Que \emph{não há quarto} é a contagem: $w$ entra na dinâmica apenas pelo sinal de
$1+3w$ (Cor.~\ref{fis:cor:destino}) e pelo expoente $3(1+w)$, e a aplicação
$w\mapsto\operatorname{sgn}$ tem \textbf{três} valores --- é a tricotomia da
ordem, a mesma do Teor.~\ref{fis:thm:leidisc}. Um quarto conteúdo exigiria um
quarto valor do sinal, e a ordem não o tem.
\end{proof}

\medskip\noindent Este é o \textbf{trial} $\{-1,0,+1\}$ --- o mesmo conjunto de
coeficientes que varre as cadeias na §\ref{fis:mapa}. A radiação está do lado
positivo, o vácuo do negativo, e a matéria no zero que os separa.

\medido{\code{tests/cosmologia.c} §C2: os três conteúdos, com os expoentes de
diluição a saírem \textbf{inteiros} --- $a^{-4}$, $a^{-3}$ e $a^{0}$ --- e o
sinal a repartir-se por $+$, $0$, $-$ sem sobrar nem faltar.}

\section{A diluição, e a constante que sai do por-unidade}\label{fis:diluicao}

\begin{teorema}[a lei de diluição, com a sua constante]\label{fis:thm:diluicao}
Da continuidade $r'+3H(r+p)=0$ com $p=w\,r$ integra-se
\[
\boxed{\;r\cdot a^{3(1+w)}=C\;}
\]
e $C$ é o \textbf{invariante} que a expansão conserva. Ela entra
\emph{aditiva} --- integrar soma --- e sai \emph{multiplicativa} em $r$: é a
involução a passar de um lado ao outro, e é a mesma travessia
soma\,$\leftrightarrow$\,produto que a forma polar faz na §\ref{fis:transf}.
\end{teorema}

\begin{proof}
Com $p=w\,r$ e $w$ constante, a continuidade é
$r'=-3H\,r\,(1+w)$. Como $H=a'/a$, isto é
\[
\frac{r'}{r}=-3(1+w)\,\frac{a'}{a} .
\]
Ambos os lados são derivadas logarítmicas, e integrar uma soma de derivadas
logarítmicas dá o logaritmo de um produto: somando, $\log r+3(1+w)\log a$ é
constante, isto é $r\,a^{3(1+w)}=C$. \emph{A constante entra aditiva} --- é a
constante de integração, que se soma --- \emph{e sai multiplicativa} em $r$,
porque a exponencial leva soma em produto. É a mesma travessia que a forma polar
faz na §\ref{fis:transf}: o dicionário entre a face que soma e a que multiplica.
E que $C$ é invariante é o que a igualdade diz: o membro esquerdo não depende da
escala.
\end{proof}

\medskip\noindent\textbf{E o que $C$ é.} Não é uma constante fundamental
introduzida pela teoria: é a \emph{constante de integração} da lei de
conservação, e é ela que carrega a normalização dimensional quando se sai do
por-unidade. O que ela vale, medido, é a \textbf{força em por-unidade}. \emph{É
por aqui, e só por aqui, que o sistema sai do por-unidade} --- todo o resto desta
parte corre com as constantes a valerem $1$.

\medido{\code{tests/cosmologia.c} §C1: $r\cdot a^{3(1+w)}=C$ conservado em
$\mathbf{36}$ escalas para seis valores de $C$, com $\mathbf{0}$ desvios; a
involução a levar o aditivo ao multiplicativo em $\mathbf{91}$ pares, com
$\mathbf{0}$ desvios, e \emph{medida sem um logaritmo} --- com $x=2^{i}$ o
produto \textbf{é} a soma dos expoentes. §C6, o \textbf{controlo}: trocado o
expoente da radiação pelo da matéria, a invariância cai --- $\mathbf{4}$ falhas
em $5$ com o expoente errado contra $\mathbf{0}$ em $5$ com o certo. \emph{A lei
escolhe o expoente, e não se escolhe a lei.}}

\medskip\noindent\textbf{E uma nota que vale a pena ficar escrita}, porque foi
um defeito real e foi o controlo que o apanhou: este bloco usa $2^{40}$ como
base, e num alvo onde o inteiro largo tem $32$ bits ele transborda --- o resíduo
dá zero e o controlo deixa de discriminar, com o expoente errado a «conservar»
em todas as escalas. \emph{O tecto tem de ser dito no tipo}, e é. Sem o
controlo, a falha teria passado por sucesso.

\section{O destino lê-se no mesmo sinal}\label{fis:destino}

\begin{corolario}[o destino, e a fronteira exacta]\label{fis:cor:destino}
A aceleração segue $-(1+3w)$. Logo
\[
\begin{array}{lll}
w=+\tfrac13 & \text{radiação} & \textbf{trava}\\
w=0 & \text{matéria} & \textbf{trava}\\
w=-1 & \text{vácuo} & \textbf{acelera}\\[2pt]
w=-\tfrac13 & \multicolumn{2}{l}{\textbf{neutro} --- a \emph{fronteira}}
\end{array}
\]
\end{corolario}

\begin{proof}
O sinal da aceleração é o de $-(1+3w)$, logo a tricotomia dela é a de $1+3w$, que
é uma função \emph{estritamente crescente} de $w$. Ela anula-se num ponto só,
$w=-\tfrac13$, é negativa abaixo e positiva acima. Avaliando nos três conteúdos
do Teor.~\ref{fis:thm:conteudos}: $w=+\tfrac13$ dá $1+3w=2>0$ e $w=0$ dá $1>0$
--- ambos travam; $w=-1$ dá $-2<0$ --- acelera. E em $w=-\tfrac13$ vale $0$: a
aceleração anula-se \textbf{exactamente}, sem aproximação, porque $1+3w=0$ é uma
igualdade entre inteiros depois de multiplicar por $3$.
\end{proof}

\medskip\noindent\textbf{E daí a leitura em três linhas}, que não precisa de
integrar a história para saber de que lado se está:
\[
\begin{array}{ccl}
w>-\tfrac13 &\Longrightarrow& \textbf{travagem}\\[2pt]
w=-\tfrac13 &\Longrightarrow& \textbf{fronteira}\\[2pt]
w<-\tfrac13 &\Longrightarrow& \textbf{aceleração}
\end{array}
\]
\emph{O sinal decide antes de a trajetória ser resolvida.} «O universo acelera»
não é então uma propriedade misteriosa de uma componente: é a afirmação de que o
estado caiu do lado negativo da mesma expressão que já decide a fronteira.

\medskip\noindent A fronteira $w=-\tfrac13$ é onde $1+3w=0$: as duas coordenadas
coincidem e \textbf{deixam de separar}. É a mesma figura do
Teor.~\ref{fis:thm:corte}(5) --- o ponto onde as duas faces partilham o ponto
fixo --- e do Cor.~\ref{fis:cor:medias}, onde a igualdade só se dá em
$\delta=0$. \emph{O destino não é uma escolha de modelo: é o sinal de uma soma.}

\medido{\code{tests/cosmologia.c} §C4: os quatro casos com o sinal a decidir, e a
fronteira $w=-\tfrac13$ a sair \textbf{exactamente} neutra.}

\section{A equação de estado fecha em quatro}\label{fis:bidual}

\noindent\textbf{Antes do enunciado, as duas involuções têm de ser nomeadas} ---
e a razão é aritmética: \emph{uma involução sozinha tem órbitas de tamanho $1$ ou
$2$; para obter quatro estados são precisas duas involuções \textbf{distintas},
\textbf{comutantes} e \textbf{independentes}}. Duas que comutem podem gerar um
subgrupo de ordem $1$, $2$ ou $4$ --- é a independência que fixa o quatro.

\begin{definicao}[as duas involuções cosmológicas]\label{fis:def:duas-inv}
Escreva-se a equação de estado como
\[
w(m)=-1+f(m),\qquad f\ \text{\textbf{ímpar} em } m ,
\]
--- é $w+1$, e não $w$, que é ímpar no grau. As duas involuções são
\[
\begin{array}{lll}
I_1: & w\longmapsto -w & \text{o \textbf{sinal da pressão}}\\[2pt]
I_2: & m\longmapsto -m & \text{o \textbf{sentido do grau}}
\end{array}
\]
\emph{Elas não agem na mesma coordenada}: $I_1$ troca o sinal do valor, $I_2$
troca o sinal do parâmetro. É daí que vem tudo o que segue.
\end{definicao}

\begin{teorema}[a bidualidade em $w$: quatro, e dois no ponto fixo]\label{fis:thm:bidual}
Com a Def.~\ref{fis:def:duas-inv}:
\begin{enumerate}\itemsep3pt
\item[(1)] \textbf{são involuções}: $I_1^{2}=I_2^{2}=\mathrm{id}$;
\item[(2)] \textbf{comutam}: $I_1I_2=I_2I_1$;
\item[(3)] \textbf{o subgrupo gerado é o de Klein}, $\langle I_1,I_2\rangle
=\{\mathrm{id},I_1,I_2,I_1I_2\}\cong C_2\times C_2$, e a órbita é
\[
\mathcal{O}(w)=\{\,w,\ I_1w,\ I_2w,\ I_1I_2w\,\}
=\{\,-1+f,\ \ 1-f,\ \ -1-f,\ \ 1+f\,\},
\]
com os quatro distintos exactamente quando
\[
\boxed{\;f\notin\{-1,\,0,\,1\}\;}
\]
\item[(4)] \textbf{e degenera em dois no ponto fixo} $m=0$: aí $f=0$, $I_2$ não
move nada, e sobram $\{-1,\,+1\}$.
\end{enumerate}
\end{teorema}

\begin{proof}
(1) $-(-w)=w$, e $-(-m)=m$.

(2) Escreva-se $w=-1+f(m)$. Então $I_2w=-1+f(-m)=-1-f$, porque $f$ é ímpar; e
$I_1w=-(-1+f)=1-f$. Compondo nas duas ordens:
\[
I_1I_2w=-(-1-f)=1+f,
\qquad
I_2I_1w=\bigl.(1-f(m))\bigr|_{m\to-m}=1-f(-m)=1+f .
\]
São iguais. \emph{$I_1$ age no valor $w$ e $I_2$ age no parâmetro $m$} --- são
coordenadas distintas, e é por isso que comutam. E são \textbf{independentes}:
nenhuma é a identidade e nenhuma é a outra, porque $I_1$ não move $m$ e $I_2$
não troca o sinal do valor. Logo $\langle I_1,I_2\rangle$ tem os quatro
elementos $\mathrm{id},I_1,I_2,I_1I_2$, e é o grupo de Klein
$C_2\times C_2$ --- \emph{a composição fecha em quatro, e não produz translação
nenhuma}.

(3) A lista sai de aplicar $I_1$, $I_2$ e a composta a $-1+f$. As colisões
possíveis são seis, e verificam-se todas:
\[
\begin{array}{lll}
-1+f=1-f &\iff& f=1\\
-1+f=-1-f &\iff& f=0\\
-1-f=1+f &\iff& f=-1\\
1-f=1+f &\iff& f=0\\
-1+f=1+f & & \text{impossível}\\
1-f=-1-f & & \text{impossível}
\end{array}
\]
Donde os quatro são distintos exactamente fora de $f\in\{-1,0,1\}$.

(4) Em $m=0$ a imparidade dá $f(0)=-f(0)$, logo $f=0$: $I_2$ é a identidade nesse
ponto, e a órbita é $\{-1,\,1\}$.
\end{proof}

\medskip\noindent É a mesma estrutura do rotor $R^{4}=\mathrm{id}$
(Teor.~\ref{fis:thm:fecho}) e do ciclo de quatro de $J$ na §\ref{fis:pi}: duas
involuções que comutam geram o grupo de quatro elementos, e quando uma delas fixa
o estado o ciclo cai a dois. \emph{O quatro não foi escolhido: é o que duas
involuções independentes dão.}

\medskip\noindent E $w+1$ sobe para $\tfrac23$ \textbf{sempre por baixo}, com a
distância ao limite a valer $\dfrac{16}{9D}$ --- o discriminante a aparecer outra
vez, e a conta exacta em inteiros.

\medido{\code{tests/cosmologia.c} §C5: $w\mapsto-w$ é involução com
$\mathbf{0}$ desvios em $17$ graus; a órbita dá $\mathbf{4}$ estados em $m=5$ e
$\mathbf{2}$ em $m=0$, o ponto fixo; e a distância ao limite $\tfrac23$ vale
$16/(9D)$ em $12$ graus, com $\mathbf{0}$ desvios --- \emph{sem tolerância
nenhuma, porque a conta é inteira}. E a troca de sinal é verificada pela
\textbf{definição}, $w+(-w)=0$ nas duas coordenadas, porque a contagem sozinha
aceitaria qualquer involução que separasse quatro estados.}

\section{O universo é uma órbita}\label{fis:universo-orbita}

\noindent A esta altura a leitura pode ser invertida. Já não é preciso dizer que
a construção \emph{se aplica} à cosmologia: a cosmologia inteira escreve-se como
uma órbita da mesma máquina.
\[
\boxed{\;
\text{estado}\;\xrightarrow{\;C\;}\;\text{estado}\;\xrightarrow{\;C\;}\;
\text{estado}\;\xrightarrow{\;C\;}\;\cdots\;}
\]
O que se chama evolução cósmica é uma trajetória no espaço de estados: a
expansão é o crescimento da coordenada de escala; a diluição é a conservação
atravessada pela mudança de escala (Teor.~\ref{fis:thm:diluicao}); a
\emph{classificação} do regime é a posição na tricotomia do operador
(Teor.~\ref{fis:thm:leidisc}); e o destino é o sinal que decide se a órbita cai
num atrator ou escapa (Cor.~\ref{fis:cor:destino}).

\medskip\noindent\emph{E uma ressalva, para não contradizer o que ficou provado
duas secções acima:} dentro da \textbf{família metálica} não há passagem entre os
três regimes --- ela é inteiramente hiperbólica, porque $\det=-1$ força
$\Delta>0$ para toda a assinatura. A mudança de regime, quando ocorre, é entre
\emph{famílias} --- entre operadores de determinante diferente --- e não ao longo
da órbita de um só.

\medskip\noindent\textbf{E é a mesma dicotomia do Teor.~\ref{fis:thm:orbita}}, o
que dá à parte anterior e a esta um enunciado só: a órbita \emph{fecha} ou
\emph{escapa}, e o que muda com a escala é o nome que se lhe dá --- período, num
átomo; destino, num universo.

\medskip\noindent O universo não precisa de carregar uma lei diferente em cada
época. \emph{A mesma transformação produz regimes diferentes porque a órbita
atravessa classes diferentes do mesmo operador.} Daí a cadeia:
\[
\boxed{\;\text{discreto}\longrightarrow\text{espectral}
\longrightarrow\text{dinâmico}\longrightarrow\text{cosmológico}\;}
\]
\emph{A escala muda. O operador não.}

\section{O que a cosmologia herda das outras partes}\label{fis:cosmo-herda}

\noindent Esta parte não trouxe objeto novo. Trouxe a mesma construção lida numa
escala diferente, e o quadro fecha isso:

\begin{center}\small
\begin{tabular}{@{}lll@{}}
\toprule
na cosmologia & é, noutra parte & onde \\
\midrule
$H^{2}=1/D$ & o discriminante da caixa & Teor.~\ref{fis:thm:leidisc} \\
os três conteúdos & o trial $\{-1,0,+1\}$ & §\ref{fis:mapa} \\
$C$ aditiva $\to$ multiplicativa & a forma polar & §\ref{fis:transf} \\
$w=-\tfrac13$ neutro & o ponto fixo comum às duas faces & Teor.~\ref{fis:thm:corte} \\
a órbita de quatro & o rotor $R^{4}=\mathrm{id}$ & Teor.~\ref{fis:thm:fecho} \\
a expansão da família metálica & o regime hiperbólico & Teor.~\ref{fis:thm:leidisc} \\
\bottomrule
\end{tabular}
\end{center}

\medskip\noindent\textbf{E é a mesma economia que a §\ref{fis:leitura-dados}
pede.} O grande e o pequeno não recebem duas teorias: recebem a mesma, com o
mesmo discriminante a decidir o regime, o mesmo trial a dar os conteúdos, e a
mesma involução a atravessar o par soma\,/\,produto. Em três escalas, e é a mesma
máquina:
\[
\boxed{\;
\begin{array}{c}
\textbf{PARTÍCULA}\\ \downarrow\\ \text{modo espectral}
\end{array}
\qquad
\begin{array}{c}
\textbf{NÚCLEO}\\ \downarrow\\ \text{órbita de ligação}
\end{array}
\qquad
\begin{array}{c}
\textbf{UNIVERSO}\\ \downarrow\\ \text{órbita cosmológica}
\end{array}
\;}
\]

\medskip\noindent\textbf{Na leitura desta construção}, a partícula corresponde à
órbita mínima; o núcleo, à órbita composta; e o universo, à órbita cosmológica.
\emph{Não há uma física para o pequeno e outra para o grande --- há uma
arquitetura, e a escala é a coordenada em que ela se lê.}
\[
\boxed{\;\text{a matéria é \textbf{espectro}}\qquad
\text{a dinâmica é \textbf{órbita}}\qquad
\text{a cosmologia é a \textbf{órbita do espectro}}\;}
\]

\medskip\noindent E daí a pergunta que governa o que vem a seguir deixa de ser
«que modelo cosmológico escolher»:
\[
\boxed{\;\text{que universo sobrevive às leis da construção?}\;}
\]

\colunas{Teor.~\ref{fis:thm:taxa},~\ref{fis:thm:conteudos},~\ref{fis:thm:diluicao},~\ref{fis:thm:bidual}; Cor.~\ref{fis:cor:destino}}
        {$H^{2}=1/D$; os três conteúdos no trial; a diluição com a sua constante; a órbita de quatro}
        {\code{tests/cosmologia.c} §C1--§C6}

% ═══════════════════════════════════════════════════════════════════════════
\part{Redes}
% ═══════════════════════════════════════════════════════════════════════════

\section*{A tese desta parte}\label{fis:redes-tese}

\noindent Um neurónio de McCulloch e Pitts é uma soma pesada com limiar. Uma rede
de Hopfield é uma matriz de pesos com uma regra local de gravação e uma descida
até um mínimo. \emph{Nenhuma das duas precisa de objeto novo aqui:} a primeira é
uma \textbf{cisão seguida de soma}, e a segunda é uma \textbf{realização} cuja
multiplicidade tem nome próprio --- \emph{interferência}.

\begin{center}\small
\begin{tabular}{@{}lll@{}}
\toprule
na rede & é, nesta construção & onde \\
\midrule
o neurónio & a cisão $\oplus$ seguida da soma $\sum$ & §\ref{fis:neuronio} \\
a sobreposição & o \textbf{prefixo comum} & Teor.~\ref{fis:thm:sobrepoe} \\
recuperar & \emph{descer} --- energia, ou profundidade & Teor.~\ref{fis:thm:desce} \\
a interferência & a \textbf{dobra}: $G>1$ & Teor.~\ref{fis:thm:interfere} \\
a árvore & o \textbf{levantamento}: $\tilde G\equiv1$ & Teor.~\ref{fis:thm:levant} \\
a parte simétrica & \emph{mede} --- converge a ponto fixo & Teor.~\ref{fis:thm:duastorres} \\
a parte antissimétrica & \emph{ordena} --- cicla com período $2$ & Teor.~\ref{fis:thm:duastorres} \\
\bottomrule
\end{tabular}
\end{center}

\section{O neurónio é a cisão seguida da soma}\label{fis:redes-neuronio}

\noindent O neurónio booleano soma entradas pesadas e compara com um limiar. Na
construção isso não é uma peça nova: é o algoritmo da §\ref{fis:neuronio}
aplicado uma vez.
\[
e=\operatorname{pop}(b\wedge\mathtt{0xAA}),
\qquad
o=\operatorname{pop}(b\wedge\mathtt{0x55}),
\qquad
\text{sai } [\,e+o,\ e\,] .
\]
As duas máscaras \textbf{particionam} --- é a cisão $\oplus$ --- e o
\emph{popcount} é o $\sum$. \emph{Os dados ficam no disco; só a conta vai à
memória.} E o par que sai é o gato aplicado uma vez, logo é um
\textbf{levantamento} (Teor.~\ref{fis:thm:levant}): guardado $(e+o,\,e)$, a fase
que não aparece recupera-se por $o=(e+o)-e$.

\medskip\noindent\textbf{Donde o neurónio da casa não perde informação}, ao
contrário do neurónio com limiar: o limiar é uma projeção com dobra --- muitas
entradas dão a mesma saída --- e o par $[\,e+o,\ e\,]$ é a mesma conta com a
folha guardada. \emph{É a diferença entre $\pi$ e $\tilde\pi$, dita numa rede.}

\section{A sobreposição é o prefixo comum}\label{fis:sobrepoe}

\begin{teorema}[a sobreposição é a contagem dos níveis que concordam]\label{fis:thm:sobrepoe}
Sejam $\xi,\eta$ dois padrões de $N$ níveis, cada um a valer $\pm1$. Então a
sobreposição normalizada é
\[
\boxed{\;\langle\xi,\eta\rangle=\frac{\#\{\text{concordam}\}-\#\{\text{discordam}\}}{N}\;}
\]
--- \emph{uma contagem, e nada mais}. E \textbf{sob a hipótese da árvore} --- as
duas folhas coincidem nos primeiros $q$ níveis e \emph{discordam em todos os
seguintes}, que é o que duas folhas que se separam e seguem ramos complementares
fazem --- vale
\[
\langle\xi,\eta\rangle=\frac{2q-N}{N} ,
\]
que é \textbf{estritamente crescente} em $q$: partilhar mais prefixo é sobrepor
mais.
\end{teorema}

\begin{proof}
$\langle\xi,\eta\rangle=\tfrac1N\sum_i\xi_i\eta_i$, e cada termo vale $+1$ se os
níveis concordam e $-1$ se discordam, porque $\xi_i,\eta_i\in\{\pm1\}$; somar dá
a diferença das contagens sobre $N$. Sob a hipótese da árvore há $q$ concordâncias
e $N-q$ discordâncias, donde $\bigl(q-(N-q)\bigr)/N=(2q-N)/N$, cuja derivada em
$q$ é $2/N>0$.

\emph{A hipótese é necessária e não é decorativa}: sem ela, dois padrões podem
partilhar um prefixo longo e voltar a concordar mais abaixo, e a sobreposição
deixa de ser função só de $q$.
\end{proof}

\medskip\noindent\textbf{E isto liga a rede à ultramétrica, sem analogia.} A
distância da árvore é $d=2^{-\operatorname{prof}}$ (Def.~\ref{fis:def:arvore}), e
a sobreposição é uma função crescente do mesmo $\operatorname{prof}$. \emph{Duas
memórias próximas na rede são duas memórias na mesma bola} --- e as bolas
particionam (Teor.~\ref{fis:thm:bolas}).

\section{Recuperar é descer}\label{fis:desce}

\begin{teorema}[a descida, e porque ela não sobe]\label{fis:thm:desce}
Seja $W$ \textbf{simétrica} com diagonal nula e a actualização assíncrona
$s_i\leftarrow\operatorname{sgn}\bigl(\sum_j W_{ij}s_j\bigr)$. Então a energia
$E=-\tfrac12\sum_{ij}W_{ij}s_is_j$ \textbf{nunca sobe}, e a trajetória atinge um
ponto fixo em número finito de passos.
\end{teorema}

\begin{proof}
Trocando um só $s_k$, a variação é
$\Delta E=-\Delta s_k\sum_{j}W_{kj}s_j$, porque a simetria faz as duas somas
sobre $k$ coincidirem e a diagonal nula tira o termo próprio. Ora $s_k$ só muda
quando o seu sinal difere do de $h_k=\sum_jW_{kj}s_j$, e nesse caso
$\Delta s_k$ tem o sinal de $h_k$; logo $\Delta s_k\,h_k\ge0$ e $\Delta E\le0$.
Como o espaço de estados é \emph{finito} --- são $2^{N}$ configurações --- e $E$
não sobe, a descida é bem-fundada e pára. \emph{É a mesma forma do
Teor.~\ref{fis:thm:corte}: uma quantidade que desce estritamente num conjunto
finito termina, e termina num ponto fixo.}
\end{proof}

\medskip\noindent E na árvore recuperar é também \emph{descer} --- a
profundidade, até uma folha. \emph{São a mesma operação}: a Hopfield desce a
energia até um mínimo, a árvore desce o prefixo até uma folha, e o
Teor.~\ref{fis:thm:sobrepoe} diz que a régua é a mesma.

\section{A interferência é a dobra}\label{fis:interfere}

\noindent Aqui está a diferença, e ela é o resultado.

\begin{teorema}[somar identifica; separar não]\label{fis:thm:interfere}
A regra de Hebb grava $P$ padrões numa \textbf{única} matriz,
\[
W_{ij}=\sum_{p=1}^{P}\xi_i^{p}\xi_j^{p} ,
\]
que é uma \textbf{soma}. Logo $W$ é a imagem de uma realização
$\pi\colon\{\text{padrões}\}\to\{\text{matrizes}\}$ que \emph{não é injetiva}:
conjuntos distintos de padrões podem dar a mesma $W$, e a sua multiplicidade
$G>1$ é exactamente o que se chama \textbf{interferência}. Já a árvore põe cada
memória no seu \emph{ramo}, uma por prefixo: aí a realização é injetiva e
$G\equiv1$.
\end{teorema}

\begin{proof}
A soma é comutativa e não retém as parcelas: pelo Teor.~\ref{fis:thm:palavra}, um
morfismo que soma tem \textbf{núcleo}, e dois conjuntos de padrões cuja diferença
esteja no núcleo têm a mesma imagem --- que é $G>1$ pela cláusula~(2) do
Teor.~\ref{fis:thm:mult}. Na árvore, cada memória ocupa um prefixo distinto, e
prefixos distintos são bolas disjuntas (Teor.~\ref{fis:thm:bolas}), donde a
aplicação é injetiva e $G\le1$.
\end{proof}

\medskip\noindent\textbf{E daí a saturação, com o seu preço.} Porque soma, a
Hopfield deixa de recuperar acima de uma capacidade --- e isso mede-se. A árvore
\emph{não sofre a mesma saturação por interferência}: ao acrescentar uma
coordenada estrutural --- o ramo --- ela separa memórias em vez de as sobrepor.
\emph{Mas isso não é capacidade infinita de graça}: ela preserva distinções
\textbf{enquanto houver espaço para os ramos}, e cada distinção preservada ocupa
lugar na hierarquia.
\[
\boxed{\;\text{a árvore troca \textbf{interferência} por \textbf{espaço}}\;}
\]
A Hopfield é $N^{2}$ \textbf{fixo}; a árvore \textbf{cresce} com o que guarda. É
o Cor.~\ref{fis:cor:folga} outra vez: quem enche paga com dobra, quem conta paga
com buraco --- \emph{a folga é uma só, vista pelos dois mapas}.

\medskip\noindent E a frase de síntese diz-se com cuidado, porque «é a Hopfield»
sugeriria equivalência funcional entre dois mecanismos que não a têm:
\emph{a árvore realiza explicitamente, pela hierarquia, a separação que a soma de
Hebb comprime e que aparece como interferência}.

\medskip\noindent\textbf{E é o levantamento, dito numa rede.} Desfazer a
interferência é exactamente o Teor.~\ref{fis:thm:levant}: acrescentar
\emph{uma} coordenada --- aqui, o ramo --- e a dobra desaparece,
$\tilde G\equiv1$. \emph{A hierarquia é a folha.}

\section{As duas torres: o que mede e o que ordena}\label{fis:duastorres}

\noindent A descida do Teor.~\ref{fis:thm:desce} usou a \textbf{simetria} de $W$
--- e usou-a de forma essencial. A outra metade da matriz faz outra coisa.

\begin{teorema}[a partição é única, e cada metade faz uma coisa só]\label{fis:thm:duastorres}
Toda matriz decompõe-se de \textbf{maneira única} em
\[
W=W_{s}+W_{a},\qquad
W_{s}=\tfrac12(W+W^{\!\top}),\qquad
W_{a}=\tfrac12(W-W^{\!\top}),
\]
com $W_s$ simétrica e $W_a$ antissimétrica. E as duas metades fazem coisas
disjuntas:
\begin{enumerate}\itemsep3pt
\item[(1)] \textbf{$W_s$ mede}: a energia não sobe e a órbita converge a um
\emph{ponto fixo} --- é o Teor.~\ref{fis:thm:desce};
\item[(2)] \textbf{$W_a$ não mede}: a forma quadrática é \emph{identicamente
nula},
\[
\boxed{\;W_a^{\!\top}=-W_a\ \Longrightarrow\ s^{\!\top}W_a\,s=0\ \text{ para todo } s\;}
\]
e a dinâmica deixa, portanto, de ter a função de Lyapunov que a energia fornecia.
\end{enumerate}
\emph{E é só isto que a antissimetria dá.} O que a órbita faz depois --- ciclar,
e com que período --- \textbf{não se segue daqui}, e precisa de hipótese sobre a
actualização.
\end{teorema}

\begin{proof}
A existência é a identidade $W=\tfrac12(W+W^{\!\top})+\tfrac12(W-W^{\!\top})$; a
unicidade é que se $W=S+A$ com $S$ simétrica e $A$ antissimétrica, transpor dá
$W^{\!\top}=S-A$, e somar e subtrair devolvem $S$ e $A$.

(1) é o Teor.~\ref{fis:thm:desce}.

(2) Para $W_a$ antissimétrica e qualquer estado $s$,
\[
\sum_{ij}(W_a)_{ij}s_is_j=0 ,
\]
porque cada par $(i,j)$ entra com $(W_a)_{ij}s_is_j$ e $(W_a)_{ji}s_js_i$, e
$(W_a)_{ji}=-(W_a)_{ij}$: os termos cancelam aos pares, e os diagonais são nulos.
Logo a forma quadrática é identicamente nula --- \emph{não existe energia a
descer}, e o argumento de (1) não tem sobre o que operar.
\end{proof}

\medskip\noindent\textbf{E aqui é preciso separar duas coisas que se parecem e
não são.} Uma é uma propriedade do \emph{operador}, $I^{2}=\mathrm{id}$; a outra
é uma propriedade da \emph{órbita}, $s_{t+2}=s_t$. A primeira é álgebra e está
provada acima; a segunda é dinâmica, e \textbf{não se segue} da antissimetria
sozinha --- pede hipótese sobre a regra de actualização.
\[
\boxed{\;W_a^{\!\top}=-W_a\ \Longrightarrow\ s^{\!\top}W_a s=0
\qquad\text{mas}\qquad
\text{o período é questão da \textbf{dinâmica}}\;}
\]

\medskip\noindent\textbf{E o que se mede diz exactamente onde a fronteira está.}
Com a actualização considerada nesta construção e $W$ \emph{puramente}
antissimétrica, a órbita observada fecha com período $2$ (§F7). Mas somando as
duas metades --- $W_s+W_a$, que é a rede completa --- \emph{a órbita medida não
passa a ciclar}: o bloco que o afirmava \textbf{falha} (§F10). Fica escrito
porque é o resultado: \emph{a antissimetria tira a função de Lyapunov, e não põe
um ciclo no lugar dela}.

\medskip\noindent\textbf{Feita a separação, o par que atravessa o documento
mantém-se.} A parte simétrica é a que \emph{mede} e não vê a ordem; a
antissimétrica é a que \emph{ordena} e não dá tamanho --- exactamente as duas
faces do Teor.~\ref{fis:thm:duas}(4). E o que se mede das duas \emph{interfaces}
é a diferença de ordem: a branca tem ordem $2$ e espelha, a negra tem ordem $4$ e
roda (§F12) --- que são o $\nu$ e o $i$ da §\ref{fis:borda}, com $\det-1$ e
$\det+1$.
\[
\boxed{\;\text{a \textbf{memória} é a parte que mede; a \textbf{ordem} é a parte
que não mede}\;}
\]

\medido{\code{tests/hopfield.c} --- \textbf{32 unidades, 2 falhas}, e as duas
falhas são resultado e não acidente. §F1: a energia nunca sobe num passo
assíncrono, em milhares de passos, e $8$ padrões corrompidos a $10\%$ voltam
todos ao original. §F3: a sobreposição é \emph{exactamente} a conta dos níveis
que concordam, em $4096$ pares sem um erro, e cresce com o prefixo. §F4:
recuperar é descer nos dois, e os dois caem no mesmo padrão. §F5: a árvore nunca
fica abaixo da Hopfield e fica acima em pelo menos um $P$, \emph{sem limiar}; e o
preço em espaço está medido, não escondido. §F7: Hebb dá matriz \textbf{simétrica}
nas $16\,384$ entradas --- com o controlo de a matriz não ser nula, sem o que a
simetria valeria por $0=0$. §F8: a partição $W=W_s+W_a$ é exacta, medida
\emph{em inteiros e por igualdade}, sem limiar. §F12: as duas interfaces são
duais e a diferença é o \textbf{sinal} --- a branca com $u\cdot u=+1$, ordem $2$,
espelha; a negra com $i\cdot i=-1$, ordem $4$, roda; $1024$ pares na tabela,
$1024$ fecham.
\smallskip\emph{E as duas que falham.} §F2: a Hopfield \textbf{satura} --- isso
passa --- mas a capacidade medida \emph{não} cai na vizinhança do $0{,}138\,N$ da
literatura: com $N=128$ ela recupera $100\%$ até $P=38$, que é cerca de
$0{,}30\,N$. O próprio bloco anota que $N=128$ é pequeno; \emph{a saturação é o
resultado, o valor da constante não}. §F10: o bloco que afirmava que somar a
torre negra à branca faz a rede \textbf{passar a ciclar} falha --- e é
exactamente o que o Teor.~\ref{fis:thm:duastorres} agora não afirma: a
antissimetria tira a função de Lyapunov e não põe um ciclo no lugar dela.
\smallskip\emph{Nota de execução:} estes medidores pedem o vector em disco ---
«o ficheiro é o vector» --- que em POSIX é \code{mmap}. O ramo equivalente para
hospedeiros sem \code{mmap} está agora no \code{lib/disco.h}, com
\code{CreateFileMapping}\,/\,\code{MapViewOfFile}: é a mesma semântica de
partilha, e não uma segunda implementação do disco. Foi por não existir esse ramo
que estes dois blocos ficaram por correr até aqui --- \emph{e as duas falhas
acima são o que aparece quando eles correm}.}

\colunas{Teor.~\ref{fis:thm:sobrepoe},~\ref{fis:thm:desce},~\ref{fis:thm:interfere},~\ref{fis:thm:duastorres}}
        {o neurónio como cisão$+$soma; a árvore como levantamento; as duas torres}
        {\code{tests/hopfield.c} §F1--§F10; \code{tests/neuronio.c}}

% ═══════════════════════════════════════════════════════════════════════════
\section{As dualidades, reunidas}\label{fis:dualidades}

\noindent Antes da tabela, a dualidade que este documento não escolheu ---
\emph{deduziu}. O Teor.~\ref{fis:thm:duas}(4) mostra que uma composição não chega
e obriga a \textbf{duas}. Tudo o que se seguiu vive nessa divisão, e ela tem um
nome clássico em cada andar:

\begin{center}\footnotesize
\begin{tabular}{@{}lll@{}}
\toprule
face \textbf{aditiva} & face \textbf{multiplicativa} & onde \\
\midrule
$\partial x+x=0$: o oposto & $\partial x\cdot x=1$: o inverso & Teor.~\ref{fis:thm:duas}(4) \\
meio $m=\frac{a+b}{2}$ & meio $g_{\ast}=\sqrt{ab}$ & Teor.~\ref{fis:thm:duas}(2), Lema~\ref{fis:lem:geo} \\
o \textbf{comprimento}: o raio & a \textbf{área}: a altura & Cor.~\ref{fis:cor:medias} \\
a \textbf{reta}: os diádicos $k/2^{n}$ & o \textbf{círculo}: as raízes & Teor.~\ref{fis:thm:duas}(6) \\
as fases \textbf{somam-se} & os módulos \textbf{multiplicam-se} & §\ref{fis:transf} \\
cartesiano & polar & §\ref{fis:transf} \\
$\oplus$: a cisão & $\otimes$: o gato & Def.~\ref{fis:def:gato} \\
\emph{endereça} (um bit por nível) & \emph{produz} (o par $(a,b)$ basta) & Teor.~\ref{fis:thm:corte} \\
\bottomrule
\end{tabular}
\end{center}

\medskip\noindent Duas coisas a reter. \textbf{Primeira:} as duas faces não são
simétricas no que \emph{fazem} --- uma endereça e a outra produz --- mas são-no no
que \emph{são}: cada uma é uma dobra, com o seu invariante e o seu ponto fixo, e o
invariante de cada uma é o neutro da outra. \textbf{Segunda:} juntas elas
\emph{fecham}, e o ponto onde fecham é o ponto fixo que as duas passam a partilhar:
\[
\boxed{\;\text{o corte é o \textbf{ponto fixo comum} às duas faces}\;}
\]
--- não um limite entre elas, nem um habitante trazido de fora.

\medskip\noindent E a tabela que segue é a mesma figura noutro registo: um lado
\emph{comprime} e perde uma distinção; o outro \emph{repõe-na} noutro sítio, e
paga por isso. Nenhuma linha é analogia --- cada uma é um enunciado deste
documento, com a sua prova.

\begin{center}\footnotesize
\begin{tabular}{@{}llll@{}}
\toprule
o que comprime & o que repõe & o preço & onde \\
\midrule
$\pi$: a realização & $\tilde\pi$: o levantamento & uma coordenada de folha & Teor.~\ref{fis:thm:levant} \\
$\zeta$: acumular & $\mu=\zeta^{-1}$: diferenciar & --- (inverte sempre) & Teor.~\ref{fis:thm:mu},~\ref{fis:thm:zetamu} \\
encher: $\pi$ sobrejetiva & contar: $\rho$ injetiva & dobra \emph{ou} buraco & Teor.~\ref{fis:thm:contraria} \\
posição & palavra de arestas & a ordem & Teor.~\ref{fis:thm:palavra} \\
vértice: $G$ & aresta: $G_{\mathrm{ar}}$ & --- (uma determina a outra) & Teor.~\ref{fis:thm:handshake} \\
base $\{\delta_x\}$ & base $\{\chi_k\}$ & o fator $2^{m}$ & Def.~\ref{fis:def:dual}, Teor.~\ref{fis:thm:H} \\
$S$ nilpotente & $C_f$ diagonalizável & inverte sempre \emph{vs} se $\widehat f\neq0$ & Lema~\ref{fis:lem:nilp}, Teor.~\ref{fis:thm:proprios} \\
ocupação: muitos agentes & expansão: um agente & o identificador & Teor.~\ref{fis:thm:agentes} \\
construção & construtor & --- (não se repõe) & Teor.~\ref{fis:thm:agentes} \\
\bottomrule
\end{tabular}
\end{center}

\medskip\noindent Três observações, e cada uma foi demonstrada acima.
\textbf{Primeira:} os pares não são involuções --- são \emph{retrações}. Só um
lado fecha, e a exceção é $\zeta/\mu$, que inverte dos dois lados por
nilpotência. \textbf{Segunda:} o preço é sempre \emph{uma} coordenada, nunca $n$
delas --- a folha sobe um andar seja qual for a dimensão. \textbf{Terceira:} a
folga é uma só, e conta-se: o que de um lado aparece como colagem em excesso é, do
outro, lugar por preencher, e o número é $\abs{I}-\abs{X}$ nos dois.

\medskip\noindent Em uma frase:
\[
\boxed{\ \textbf{multiplicidade} \text{ é o que o campo vê};\qquad
\textbf{memória} \text{ é o que a construção preserva.}\ }
\]

\medskip\noindent E a leitura que liga o primeiro teorema ao último:
\[
\text{caminhar é \textbf{realizar}};\qquad
\text{codificar é \textbf{conservar}};\qquad
\text{a dobra \textbf{mede} a diferença.}
\]
Cada realização deste documento --- a bissecção, a trie, o índice, a órbita, a
polar, a métrica, o espectro --- é uma caminhada com duas faces, e o que dela fica
é sempre o mesmo par: o \textbf{centro} que ela conserva e o \textbf{desvio} que
ela mede. \emph{Não se postula o ponto: constrói-se o caminho que o determina.}

% ═══════════════════════════════════════════════════════════════════════════
\section{Medido}\label{fis:medido}

\medido{Teor.~\ref{fis:thm:mult} e Teor.~\ref{fis:thm:dim}:
\code{tests/aranha\_n.c} §AN1--§AN2 (campo incremental $=$ recontagem e
$\sum G=\abs{I}$ nos postos $1$ a $4$), §AN7 ($\abs{\mathcal{V}(x)}=2n$, vizinhos
distintos);
Teor.~\ref{fis:thm:dragao} §AN3 (bits $=$ dobra, $k=1..12$, com os dois sinais de
$R$ corridos e exatamente um a bater); dobra e controlos §AN4;
Teor.~\ref{fis:thm:bolas} §AN6 (bolas particionam; todo ponto é centro; a
euclidiana falha ambas com o mesmo raio);
Teor.~\ref{fis:thm:periodo} §AN8 ($\abs{\{G>1\}}=p$ e $\abs{\{G=1\}}=\ell$, com o
período obtido por busca independente e a órbita que escapa como controlo);
ciclo em qualquer posto e Teor.~\ref{fis:thm:fecho} §AN13;
Teor.~\ref{fis:thm:vizarvore} e o custo \textbf{contado} §AN14 ($2n$ leituras por
passo, nem uma a mais --- o autómato não varre);
Teor.~\ref{fis:thm:espaco} e \ref{fis:thm:base} §AN19, §AN24 (Gram
$=\mathrm{Id}$; as oito réguas cegas e compostas; os $256$ elementos a
reconstruírem-se);
Teor.~\ref{fis:thm:euler} §AN25 ($b_1=D=335$ no dragão, a espiral a reusar
arestas, a reta a ser árvore);
a tabela da §\ref{fis:realizacoes} §AN23 (medida, não escrita: os nove $\max G$,
as células e a coluna \emph{percurso});
Teor.~\ref{fis:thm:gcanon} e \ref{fis:thm:geodesica} §AN21--§AN22
($(\partial\tau)^2$ inteiro por aresta com três $g$; $g_{uv}(x)=G(x)\delta_{uv}$
curva onde há dobra e é plana com $G\equiv1$; e a geodésica a divergir do
gradiente quando $h$ é anisotrópica);
Def.~\ref{fis:def:tempo} §AN20 ($\mu D=\theta$ com resíduo $0$; e a mesma
trajetória com salto uniforme $1$ na grade e $11$ comprimentos distintos na
árvore, de $2$ a $2048$);
Teor.~\ref{fis:thm:shift} e \ref{fis:thm:handshake} §AN18 ($(1-S)\zeta=\mathrm{id}$;
quais realizações são percursos; e a lei local com $0$ violações em $513$
vértices);
Teor.~\ref{fis:thm:hiper} e Cor.~\ref{fis:cor:modos} §AN17 (grau nos três grafos;
$A\chi_k=(m-2\abs{k})\chi_k$ exato em inteiros);
Teor.~\ref{fis:thm:medida} §AN16 (a mesma medida com dobras diferentes: Heighway e
a espiral, e a reta como controlo);
Lema~\ref{fis:lem:nilp} §AN27 ($S^{12}=0$ com $12$ mínimo; $\sum S^j$ a ser a
triangular de uns; $\dim\ker S=1$);
Teor.~\ref{fis:thm:mu}, \ref{fis:thm:zetamu} e Cor.~\ref{fis:cor:mobius} §AN15
($\mu=\zeta^{-1}$ em $4096$ pares; o directo e o inverso com resíduo $0$; e a
mesma inversão na árvore dos divisores);
Def.~\ref{fis:def:doisandares}, Teor.~\ref{fis:thm:dobranorma} e \ref{fis:thm:conv}
§AN12 (as coordenadas do campo na base canónica; a dobra na norma por três
caminhos; convolução diagonalizada e deconvolução com resíduo $0$; e a
degenerescência exibida);
Def.~\ref{fis:def:op} e Teor.~\ref{fis:thm:duas} §AN36 (as $81$ soluções com
$\partial e=e$; zero soluções sobre dois com uma composição só; e exatamente $2$
com duas composições de invariantes distintos);
Teor.~\ref{fis:thm:coord} §AN35 (as duas nomeações isomorfas pela troca, com a
identidade a NÃO servir de controlo; e permutar a tupla a mudar o objeto em $32$
das $64$);
Teor.~\ref{fis:thm:B} e \ref{fis:thm:BI} §AN33 (as quatro aplicações, as duas
invertíveis e involutivas, a única tabela de produto em dezasseis, e o encaixe por
prefixos);
a codificação e a subida §AN34 ($\lceil\log_2 c\rceil$ bits para $c$ de $1$ a $64$
e nenhum a menos);
Teor.~\ref{fis:thm:escada} §AN32 (os quatro casos, com $\sum G=\abs{I}$);
Teor.~\ref{fis:thm:agentes} §AN31 (um agente com $4$ voltas e $4$ agentes com uma:
campos iguais nas $64$ células; e a ordem a separar $96$ de $96$ pares contra o
identificador a separar $64$);
Teor.~\ref{fis:thm:palavra} §AN30 (as $1024$ letras do dragão; reconstrução com
divergência $0$; e a mesma palavra com três $T$ a dar $\max G$ de $2$, $1$ e $28$);
Teor.~\ref{fis:thm:contraria} e Cor.~\ref{fis:cor:folga} §AN29
($\sum(G-1)=256=\abs{I}-\abs{X}$; $\pi\circ\rho=\mathrm{id}$ e $\rho\circ\pi$ a
falhar em $256$ de $512$; e o grau $4$ contra $2$);
Teor.~\ref{fis:thm:bijeccao} e \ref{fis:thm:enumera} §AN37 (dobra $0$ e buraco $0$
em $M=2,4,8,16,32$);
Lema~\ref{fis:lem:comuta}, Teor.~\ref{fis:thm:proprios} e \ref{fis:thm:H} §AN26;
Teor.~\ref{fis:thm:leitura} e Cor.~\ref{fis:cor:duais} §AN10;
Teor.~\ref{fis:thm:serie} §AN11 (período $1$, cauda $=$ custo, e as três
combinações de $\arctan$ a coincidirem sem um dígito escrito à mão);
Teor.~\ref{fis:thm:levant} §AN5 e §AN9;
Teor.~\ref{fis:thm:corte} \code{tests/meio.c} §M8 e \code{tests/agm.c}
§A1--§A3;
Teor.~\ref{fis:thm:gato} e \ref{fis:thm:bitunico} \code{papers/aranha.c} e
\code{tests/neuronio.c};
a caixa e a óptica \code{tests/pgwire.c} §W116, §W211--§W215, §W217--§W220,
§W222--§W225 e \code{banco/sql.c};
o plano \code{tests/aranha\_g.c} §AG1--§AG11 e \code{tests/aranha\_inversa.c}
§AI1--§AI7.}

\colunas{Teor.~\ref{fis:thm:mult} (só igualdade); Teor.~\ref{fis:thm:levant} (a fibra volta como uma coordenada); Teor.~\ref{fis:thm:bolas} (partição); Teor.~\ref{fis:thm:periodo} (a cauda e o período no campo); Teor.~\ref{fis:thm:leitura}}
        {campo esparso, $\Theta(n)$/passo, memória $\Theta(\abs{I})$, sem alocação dinâmica; postos $1$ a $4$, árvore, órbitas}
        {\code{tests/aranha\_n.c} §AN1--§AN37; \code{tests/aranha\_g.c} §AG1--§AG11; \code{tests/aranha\_inversa.c} §AI1--§AI7}

\vfill
{\footnotesize\color{cinza}\noindent
Documento autocontido: todo enunciado usado está demonstrado acima.
Números: \code{bash tools/bateria.sh}.
Proprietário --- uso mediante contrato e pagamento (ver \texttt{LICENSE}).\par}

\end{document}
